\documentclass[a4paper,11pt,titlepage]{scrbook}

\let\cleardoublepage\clearpage

\usepackage[utf8]{inputenc}
\usepackage[spanish]{babel}
\usepackage[strings]{underscore}

% \usepackage[style=list, number=none]{glossary} %si se va a usar glosario, quitar marca de comentario
%\usepackage{titlesec}
%\usepackage{palatino} %usar fot palatino en vez de times roman

%\decimalpoint %revisar
%\usepackage{dcolumn} %revisat
%\newcolumntype{.}{D{.}{\esperiod}{-1}}
%\makeatletter
%\addto\shorthandsspanish{\let\esperiod\es@period@code}
%\makeatother


%\usepackage[chapter]{algorithm}
%\RequirePackage{verbatim}
%\RequirePackage[Glenn]{fncychap}
\usepackage{longtable}
\usepackage{array}
\usepackage{float}
\usepackage{fancyhdr}
\usepackage{placeins}
\usepackage{afterpage}

% Estilo común para todas las tablas
\newcolumntype{L}[1]{>{\raggedright\arraybackslash\hspace{0pt}}p{#1}}
\newcolumntype{C}[1]{>{\centering\arraybackslash\hspace{0pt}}p{#1}}

\newcommand{\estilotabla}{%
    \small
    \setlength{\tabcolsep}{4pt}
    \renewcommand{\arraystretch}{1.15}
}

\raggedbottom
\usepackage{float}
\usepackage{longtable}
\usepackage{subcaption}
\usepackage{xcolor}
\definecolor{portada}{RGB}{239,206,53}
\definecolor{base}{RGB}{35,31,32}
\usepackage{pdfpages}

\usepackage[acronym]{glossaries}
\usepackage{acronym}

\makenoidxglossaries
\input{glosario/entradas_glosario.tex}


%Instrucciones para poder escribir código y mostrarlo de manera elegante:
\definecolor{gray97}{gray}{.97}
\definecolor{gray75}{gray}{.75}
\definecolor{gray45}{gray}{.45}
\definecolor{gray30}{gray}{.94}

\usepackage{listings}
\lstset{ frame=Ltb,
framerule=0pt,
aboveskip=0.5cm,
framextopmargin=3pt,
framexbottommargin=3pt,
framexleftmargin=0.4cm,
framesep=0pt,
rulesep=.4pt,
backgroundcolor=\color{gray97},
rulesepcolor=\color{black},
%
stringstyle=\ttfamily,
showstringspaces = false,
basicstyle=\small\ttfamily,
commentstyle=\color{gray45},
keywordstyle=\bfseries,
%
numbers=left,
numbersep=15pt,
numberstyle=\tiny,
numberfirstline = false,
breaklines=true,
literate={á}{{\'a}}1 {Á}{{\'A}}1 {é}{{\'e}}1 {É}{{\'e}}1 {í}{{\'i}}1  {Í}{{\'I}}1  {ó}{{\'o}}1  {Ó}{{\'O}}1  {ú}{{\'u}}1  {Ú}{{\'U}}1  {Ñ}{{\~N}}1 {ñ}{{\~n}}1 ,
}



% minimizar fragmentado de listados
\lstnewenvironment{listing}[1][]
   {\lstset{#1}\pagebreak[0]}{\pagebreak[0]}

\lstdefinestyle{Consola}
   {basicstyle=\scriptsize\bf\ttfamily,
    backgroundcolor=\color{gray30},
    frame=single,
    numbers=none
   }
\lstdefinestyle{C}
	{basicstyle=\scriptsize,
	frame=single,
	language=C,
	numbers=left
	}
\lstdefinestyle{CodigoC++}
        {basicstyle=\small,
	frame=single,
	backgroundcolor=\color{gray30},
	language=C++,
	numbers=left
 	}
\lstdefinestyle{PHP}
	{basicstyle=\scriptsize,
%        {basicstyle=\small,
	frame=single,
	language=PHP,
	numbers=left
	}
	



% ********************************************************************
% Información sobre el TFG. Comentar lo que NO se desee añadir y sustituir con la información correcta.
% ********************************************************************
\newcommand{\myTitle}{Diseño de un pipeline de entrenamiento de modelos generativos orientado a la producción de imágenes en estilo pixel-art}
\newcommand{\myDegree}{Grado en Ingeniería Multimedia}
\newcommand{\myName}{Izan Quiles Jiménez}
\newcommand{\myProf}{Pedro Agustín Pernías Peco}
\newcommand{\myOtherProf}{Maria Pilar Escobar Esteban}
\newcommand{\myFaculty}{Escuela Politécnica Superior de la Universidad de Alicante}
\newcommand{\myFacultyShort}{EPS UA}
\newcommand{\depTutorOne}{Departamento del tutor}
\newcommand{\depTutorTwo}{Departamento del cotutor}


\newcommand{\myUni}{\protect{Universidad de Alicante}}
\newcommand{\myLocation}{Alicante}
\newcommand{\myTime}{\today}
%\newcommand{\myVersion}{Version 0.1}

\newcommand{\logoGrado}{imagenes/logoim.jpg}
\newcommand{\logoFacultad}{imagenes/logoeps.jpg}
\newcommand{\logoUniversidad}{imagenes/logoua.jpg}

\usepackage{url}

% Definición de comandos que me son útiles:
%\renewcommand{\indexname}{Índice alfabético}
%\renewcommand{\glossaryname}{Glosario}

\pagestyle{fancy}
\fancyhf{}
\fancyhead[LO]{\leftmark}
\fancyhead[RE]{\rightmark}
\fancyhead[RO,LE]{\textbf{\thepage}}
\renewcommand{\chaptermark}[1]{\markboth{\textbf{#1}}{}}
\renewcommand{\sectionmark}[1]{\markright{\textbf{\thesection. #1}}}


\setlength{\headheight}{1.5\headheight}

\newcommand{\HRule}{\rule{\linewidth}{0.5mm}}
%Definimos los tipos teorema, ejemplo y definición podremos usar estos tipos
%simplemente poniendo \begin{teorema} \end{teorema} ...
\newtheorem{teorema}{Teorema}[chapter]
\newtheorem{ejemplo}{Ejemplo}[chapter]
\newtheorem{definicion}{Definición}[chapter]
 
\newcommand{\bigrule}{\titlerule[0.5mm]}


%Para conseguir que en las páginas en blanco no ponga cabeceras
\makeatletter
\def\clearpage{%
  \ifvmode
    \ifnum \@dbltopnum =\m@ne
      \ifdim \pagetotal <\topskip
        \hbox{}
      \fi
    \fi
  \fi
  \newpage
  \thispagestyle{empty}
  \write\m@ne{}
  \vbox{}
  \penalty -\@Mi
}
\makeatother

\usepackage[pdfborder={000}]{hyperref} %referencia
\hypersetup{
pdfauthor = {\myName (email (en) ua (punto) es)},
pdftitle = {\myTitle},
pdfsubject = {},
pdfkeywords = {palabra_clave1, palabra_clave2, palabra_clave3, ...},
pdfcreator = {LaTeX con el paquete ....},
pdfproducer = {pdflatex}
}
%AQUI COMIENZA LA LISTA DE FICHEROS A INCLUIR



\begin{document}
\renewcommand{\listtablename}{Índice de tablas} %para sustituir la palabra cuadro por tabla
\renewcommand{\tablename}{Tabla}
\renewcommand{\lstlistingname}{Listado}
\renewcommand{\lstlistlistingname}{Índice de \lstlistingname s}

\frontmatter
\input{portada/portada} %la portada en color
\input{portada/portada_2} %la portada en b/n

\chapter{Resumen}
El presente Trabajo de Fin de Grado consiste en el desarrollo de un pipeline completo para la generación de imágenes en estilo pixel-art mediante modelos generativos de difusión. El proyecto parte del uso de un modelo FLUX y de la técnica LoRA, con el objetivo de adaptar un modelo previamente entrenado a un estilo visual concreto sin necesidad de realizar un reentrenamiento completo. El trabajo aborda la construcción de un sistema reproducible que integra todas las fases necesarias para obtener, preparar, entrenar, generar y evaluar imágenes pixel-art.

\medskip

Para ello, se ha partido del repositorio base \textit{ai-toolkit}, ampliándolo con una capa propia de herramientas, scripts e interfaz web. El sistema desarrollado permite organizar datasets por iteraciones, preparar imágenes para el entrenamiento, sincronizar captions, configurar entrenamientos y conservar la trazabilidad de los modelos generados. Las primeras pruebas se realizaron con datasets reducidos y curados manualmente, lo que permitió validar el funcionamiento básico del flujo de trabajo y detectar las principales limitaciones iniciales.

\medskip

Posteriormente, se incorporó un sistema de obtención de datos mediante scraping de Pixilart, con filtros automáticos para descartar imágenes no válidas o poco adecuadas para el entrenamiento. Este proceso permitió ampliar el volumen y la diversidad del dataset, manteniendo una fase de revisión manual para asegurar la calidad de las imágenes seleccionadas. Además, se implementó un sistema de generación automática de captions mediante un modelo vision-language, permitiendo crear descripciones adaptadas al entrenamiento LoRA y asociadas al estilo pixel-art.

\medskip

El proyecto incluye también una interfaz web desarrollada con Next.js, TypeScript, Prisma y SQLite, desde la que se pueden gestionar datasets, lanzar procesos de captioning, revisar resultados, configurar entrenamientos, consultar iteraciones previas y generar imágenes con los LoRAs entrenados. Esta interfaz centraliza el flujo de trabajo y facilita el seguimiento de cada fase del proceso.

\medskip

A nivel experimental, el pipeline se ha ejecutado tanto en local como en la nube mediante Modal, permitiendo realizar entrenamientos con GPUs de mayor capacidad. A lo largo de las iteraciones se han detectado y corregido problemas de configuración, se han comparado resultados visuales y se han incorporado herramientas de evaluación orientadas a analizar la fidelidad al estilo pixel-art.

\medskip

Como resultado, se ha obtenido un sistema extensible y reproducible para entrenar, evaluar y mejorar LoRAs especializados en pixel-art, documentando el proceso completo y dejando una base sólida para futuras iteraciones del proyecto.

%%%%%%%%%%%%%%%%%%%%%%%Preliminares
\chapter*{Preámbulo}
\thispagestyle{empty}
La elección de este Trabajo de Fin de Grado surge principalmente de mi interés por comenzar a formarme en el sector de la inteligencia artificial, un ámbito que ha adquirido una gran relevancia en los últimos años y que está transformando muchas áreas relacionadas con la tecnología, la creación de contenidos y el diseño digital.

\medskip

Dentro de este campo, me resultó especialmente atractivo trabajar con inteligencia artificial generativa aplicada a imágenes, ya que conecta directamente con una parte más visual y creativa de la Ingeniería Multimedia. La posibilidad de experimentar con modelos capaces de generar contenido gráfico, adaptarlos a un estilo concreto y analizar cómo influyen los datos de entrenamiento en los resultados finales fue uno de los motivos principales para escoger este proyecto.

\medskip

El trabajo se centra en el desarrollo de un pipeline para la generación de imágenes en estilo pixel-art mediante modelos generativos y técnicas de ajuste como LoRA. A través de este proyecto se busca construir un flujo de trabajo completo que abarque la obtención de imágenes, la preparación del dataset, la generación de captions, el entrenamiento del modelo y la evaluación de los resultados obtenidos.

\medskip

La finalidad principal de este TFG es adquirir una visión práctica del proceso de adaptación de modelos generativos a un estilo visual específico, comprendiendo tanto las posibilidades que ofrecen estas técnicas como sus limitaciones técnicas, computacionales y metodológicas. Además, el proyecto permite aplicar conocimientos relacionados con programación, desarrollo web, tratamiento de imágenes, organización de datos y experimentación con modelos de inteligencia artificial.

\cleardoublepage %salta a nueva página impar
% Aquí va la dedicatoria si la hubiese. Si no, comentar la(s) linea(s) siguientes
\chapter*{Agradecimientos} 

\thispagestyle{empty} Para empezar, quiero agradecer a todos mis amigos dentro de la universidad, Fuentes, Juan, Marcos, Silva, Pablo, Antonio y Javi  por todos estos años en los que han estado a mi lado, apoyándome, ayudándome y, sobre todo, haciéndome reír. 

\medskip También quiero agradecer a mi grupo de amigos de Novelda por su apoyo y compañía durante toda mi vida. Gracias por estar siempre cuando más lo he necesitado, por preocuparos por mí y por formar parte de una etapa tan importante de mi vida. 

\medskip Además, quiero agradecer a Pedro Pernías y a Pablo Pernías su trabajo como tutores y supervisores de este Trabajo de Fin de Grado, así como la orientación recibida durante el desarrollo del proyecto. Quiero agradecerles especialmente las oportunidades profesionales que me han ofrecido dentro del sector de la inteligencia artificial generativa, ya que han supuesto una motivación muy importante para seguir formándome en este ámbito. 

\medskip Por último, quiero agradecer a mi familia todo el apoyo e interés que han mostrado hacia mi futuro. En especial, quiero dar las gracias a mi madre, por toda su ayuda, compañía y fortaleza durante los momentos más difíciles. Siempre será mi ejemplo a seguir.


\cleardoublepage %salta a nueva página impar

%%%%%%%%%%%%%%%%%%%%%%%Fin preliminares

%\input{capitulos/preliminares} 
%editar este texto (capitulos/preliminares.tex) para cambiar preámbulo, agradecimientos y dedicatorias
% En los preliminares podremos editar la siguiente información:
% [Preámbulo:] se describirán brevemente la motivación que ha originado la realización del TFG/TFM, así como de una breve descripción de los objetivos generales que se quieren alcanzar con el trabajo presentado.
% [Agradecimientos:] se podrá añadir las hojas necesarias para realizar los agradecimientos, a veces obligatorios, a las entidades y organismos colaboradores.
% [Dedicatoria:] se podrá añadir una única hoja con dedicatorias, su alineación será derecha y centradas de forma distribuida en la página.
% [Citas:] (frases célebres) se podrá añadir una única hoja con citas, su alineación será derecha y centradas de forma distribuida en la página.

%Con los siguientes comandos se genera el índice, indice de figuras, tablas y listados de código
\tableofcontents
\listoffigures
\listoftables

\mainmatter %entre frontmatter y mainmatter, la numeración es en romanos.



\chapter{Introducción}

La inteligencia artificial generativa se ha convertido en una herramienta cada vez más presente en procesos de creación visual. En particular, los modelos capaces de generar imágenes a partir de descripciones textuales han abierto nuevas posibilidades en ámbitos como el diseño gráfico, la ilustración, los videojuegos, la publicidad y la creación de recursos visuales \cite{amankwahamoah2024generativeai}.

\medskip

Sin embargo, que un modelo sea capaz de producir imágenes visualmente atractivas no significa necesariamente que pueda reproducir de forma consistente un estilo visual concreto. Este problema es especialmente relevante en el caso del pixel-art, un estilo caracterizado por el uso visible del píxel, la simplificación de formas, las paletas de color limitadas y una estética muy vinculada al videojuego clásico y a la cultura digital. Aunque muchos modelos generalistas pueden generar imágenes que recuerdan parcialmente a este estilo, los resultados suelen presentar problemas como exceso de suavizado, detalles demasiado complejos, pérdida de nitidez o falta de coherencia visual entre distintas generaciones \cite{hertz2024stylealigned}.

\medskip

Una posible forma de abordar esta limitación consiste en adaptar modelos generativos ya entrenados mediante técnicas de ajuste fino. En lugar de entrenar un modelo desde cero, lo cual requiere grandes cantidades de datos y recursos computacionales, se puede partir de un modelo base y especializarlo utilizando un conjunto de imágenes representativas del estilo deseado. En este contexto, técnicas como LoRA permiten realizar esta adaptación de forma más eficiente, reduciendo el número de parámetros entrenables y facilitando la experimentación con datasets de menor tamaño \cite{hu2022lora}.

\medskip

No obstante, el resultado de una adaptación de este tipo no depende únicamente de la técnica de entrenamiento utilizada. La calidad del dataset, la selección y preparación de las imágenes, las descripciones textuales asociadas, la configuración del entrenamiento y la revisión de los resultados influyen directamente en la capacidad del modelo para aprender un estilo visual. Por ello, el problema consiste en diseñar un proceso completo que permita obtener datos adecuados, organizarlos correctamente, entrenar de forma iterativa y analizar los resultados obtenidos.

\medskip

Este Trabajo de Fin de Grado propone el desarrollo de un pipeline para la generación de imágenes en estilo pixel-art mediante modelos de difusión y entrenamiento LoRA. El sistema desarrollado parte de \textit{ai-toolkit} y añade una capa propia para organizar datasets por iteraciones, preparar imágenes, generar y sincronizar captions, configurar entrenamientos LoRA, consultar resultados y generar imágenes mediante una interfaz web de apoyo.

\medskip

La aportación principal del trabajo es la construcción de un flujo reproducible que integra las distintas fases necesarias para adaptar un modelo generativo preentrenado al estilo pixel-art. El alcance del proyecto es diseñar y validar un pipeline práctico de adaptación estilística orientado al pixel-art.


\chapter{Estudio de viabilidad} \label{viabilidad}

El presente Trabajo de Fin de Grado se plantea  como un proyecto experimental y técnico orientado a diseñar, ejecutar y evaluar un pipeline para adaptar modelos generativos de imagen al estilo pixel-art. Por este motivo, no se incluye un modelo de negocio ni un Lean Canvas, ya que el objetivo principal es estudiar la viabilidad técnica y metodológica del sistema dentro del alcance académico del trabajo. 

\medskip 

El estudio de viabilidad se centra en valorar las condiciones técnicas, metodológicas y organizativas necesarias para llevar a cabo el proyecto. Para ello, se realiza un análisis DAFO adaptado al carácter técnico del trabajo y se identifican los principales riesgos que pueden afectar al desarrollo del pipeline, la construcción del dataset, la generación de captions, el entrenamiento de los modelos y la obtención de resultados.

\section{Análisis DAFO}

El análisis DAFO permite identificar los factores internos y externos que pueden influir en el desarrollo del proyecto. En este caso se plantea como una herramienta para valorar la viabilidad del pipeline propuesto para la generación de imágenes en estilo pixel-art mediante modelos de difusión y entrenamiento LoRA (ver figura \ref{fig:dafo}).

\begin{figure}
    \centering
    \includegraphics[width=0.85\textwidth]{imagenes/Gráfico Análisis DAFO Creativo Multicolor.png}
    \caption{Análisis DAFO del proyecto. Fuente: elaboración propia.}
    \label{fig:dafo}
\end{figure}

El análisis DAFO muestra que la viabilidad del proyecto se apoya principalmente en la reutilización de modelos preentrenados y en el uso de LoRA como técnica de adaptación eficiente. Estas fortalezas reducen la necesidad de entrenar un modelo desde cero y permiten centrar el trabajo en la construcción del dataset, la preparación de captions y la organización del pipeline por iteraciones.

\medskip

No obstante, estas fortalezas no eliminan ciertas debilidades relevantes. La calidad de los resultados depende directamente de la coherencia del dataset, de la configuración del entrenamiento y de la capacidad para evaluar la fidelidad al estilo pixel-art. Además, el proyecto requiere integrar distintas herramientas y procesos, como el repositorio base de entrenamiento, los scripts propios, la interfaz web, los modelos de captioning y los servicios de ejecución en la nube.

\medskip

En cuanto al contexto externo, el crecimiento de la inteligencia artificial generativa y su aplicación en videojuegos, recursos visuales y creación de contenidos representa una oportunidad clara para este tipo de sistemas. Sin embargo, también existen amenazas asociadas a la dependencia de herramientas externas, posibles restricciones de acceso a plataformas, cambios en dependencias técnicas y cuestiones legales relacionadas con la recopilación y uso de imágenes procedentes de internet.

\section{Análisis de riesgos}

El desarrollo del proyecto presenta distintos riesgos técnicos y metodológicos que deben tenerse en cuenta para asegurar su viabilidad. La Tabla~\ref{tab:riesgos} recoge los principales riesgos identificados, junto con su probabilidad, impacto y medidas de mitigación.

\begingroup
\estilotabla
\begin{longtable}{|L{0.25\textwidth}|C{0.12\textwidth}|C{0.12\textwidth}|L{0.39\textwidth}|}
\hline
\textbf{Riesgo} & \textbf{Prob.} & \textbf{Impacto} & \textbf{Medidas de mitigación} \\
\hline
\endfirsthead

\hline
\textbf{Riesgo} & \textbf{Prob.} & \textbf{Impacto} & \textbf{Medidas de mitigación} \\
\hline
\endhead

Limitaciones de hardware local &
Alta &
Alto &
Utilizar LoRA y técnicas de optimización de memoria. Recurrir a GPUs en la nube cuando el entrenamiento local no sea suficiente. \\
\hline

Errores de configuración del entrenamiento &
Media &
Alto &
Documentar configuraciones por iteración, conservar registros y realizar pruebas reducidas antes de entrenamientos largos. \\
\hline

Dataset insuficiente o poco coherente &
Alta &
Alto &
Ampliar progresivamente el dataset, aplicar filtros técnicos y revisar manualmente las imágenes aceptadas. \\
\hline

Captions incorrectas o demasiado genéricas &
Media &
Medio &
Generar captions de forma asistida, revisarlos posteriormente y mantener una estructura coherente con el contenido y el estilo. \\
\hline

Sobreajuste del modelo al dataset &
Media &
Alto &
Evitar datasets demasiado pequeños o repetitivos, comparar iteraciones y revisar si el modelo memoriza ejemplos concretos. \\
\hline

Coste de computación en la nube &
Baja &
Medio &
Controlar la duración de las ejecuciones, lanzar pruebas cortas y ajustar la configuración al presupuesto disponible. \\
\hline

Dependencia de herramientas externas &
Media &
Medio &
Mantener una capa propia del pipeline, documentar versiones utilizadas y evitar depender exclusivamente de una única plataforma. \\
\hline

Problemas legales o de licencia en las imágenes &
Media &
Alto &
Registrar fuentes, evitar redistribuir imágenes sin permiso y priorizar recursos con condiciones de uso claras cuando sea posible. \\
\hline

Resultados visuales insuficientes &
Media &
Medio &
Plantear el proyecto como un proceso experimental, documentando tanto las mejoras obtenidas como las limitaciones detectadas. \\
\hline

\caption{Principales riesgos identificados en el desarrollo del proyecto.}
\label{tab:riesgos}
\end{longtable}
\endgroup

El análisis de riesgos muestra que el proyecto es viable siempre que se aborde de forma iterativa, controlando la calidad del dataset, documentando las configuraciones de entrenamiento y manteniendo alternativas ante limitaciones de cómputo o dependencia de herramientas externas. Esta forma de trabajo permite realizar pruebas progresivas, conservar la trazabilidad de cada experimento y ajustar el pipeline a partir de los resultados obtenidos.



\chapter{Marco teórico y estado del arte} \label{marcoteorico}

Este capítulo presenta los conceptos necesarios para comprender el desarrollo del pipeline propuesto. Se abordan los fundamentos de la generación texto-imagen, los modelos de difusión, la adaptación de modelos preentrenados mediante fine-tuning y LoRA, y la preparación de datasets orientados a un estilo visual concreto. Además, se revisan las particularidades del pixel-art como dominio visual y se analiza de forma cualitativa el comportamiento de varias plataformas generalistas ante prompts de este estilo.

\section{Inteligencia artificial generativa}

La inteligencia artificial generativa es una rama de la inteligencia artificial centrada en la creación de contenido nuevo a partir de los patrones aprendidos durante el entrenamiento. A diferencia de otros sistemas de aprendizaje automático orientados principalmente a clasificar, predecir o reconocer información existente, los modelos generativos buscan aprender la distribución de los datos de entrenamiento para producir nuevas muestras con características similares. Este contenido puede adoptar distintas formas, como texto, imágenes, audio, vídeo, código fuente u otros tipos de datos digitales \cite{feuerriegel2024generative}.

\medskip

La generación automática de datos no es una idea nueva dentro de la inteligencia artificial. Sin embargo, el avance del aprendizaje profundo, el aumento de la capacidad de cómputo y la disponibilidad de grandes conjuntos de datos han permitido que estos sistemas alcancen resultados mucho más complejos. En el ámbito de la imagen, uno de los hitos más relevantes fueron las redes generativas adversarias, conocidas como GANs, introducidas por Goodfellow et al. Este enfoque se basa en el entrenamiento conjunto de una red generadora, que produce imágenes sintéticas, y una red discriminadora, que intenta distinguir entre imágenes reales y generadas \cite{goodfellow2014generative}.

\medskip

Posteriormente, otros enfoques han ganado protagonismo, especialmente los modelos de difusión, que se han convertido en una de las bases principales de muchos sistemas actuales de generación de imágenes. Estos modelos permiten generar imágenes de alta calidad y pueden condicionarse mediante texto, lo que los hace especialmente relevantes para sistemas de generación texto-imagen.

\section{Generación de imágenes a partir de texto}

La generación de imágenes a partir de texto, conocida habitualmente como \textit{text-to-image generation}, consiste en producir una imagen a partir de una descripción escrita o \textit{prompt}. En este tipo de sistemas, el usuario introduce una instrucción en lenguaje natural y el modelo genera una imagen que intenta representar los elementos indicados en dicha descripción.

\medskip

Este proceso combina dos capacidades. Por un lado, el modelo debe interpretar el contenido semántico del texto, identificando elementos como el sujeto principal, el estilo visual, la composición o la perspectiva. Por otro lado, debe transformar esa información en una imagen visualmente coherente. Por esta razón, los modelos texto-imagen se sitúan dentro del ámbito de los modelos multimodales, al relacionar información textual y visual durante el entrenamiento.

\medskip

Uno de los trabajos más representativos en este campo es DALL·E, presentado por Ramesh et al., que mostró la posibilidad de generar imágenes a partir de descripciones textuales mediante un modelo entrenado con grandes cantidades de pares texto-imagen. Este trabajo contribuyó a popularizar la generación de imágenes desde lenguaje natural y demostró que un modelo podía combinar conceptos, atributos y estilos visuales a partir de instrucciones escritas \cite{ramesh2021zeroshot}.

\medskip

Aun así, generar una imagen visualmente atractiva no implica que el resultado sea adecuado para cualquier objetivo. Los modelos pueden omitir detalles del \textit{prompt}, interpretar de forma ambigua ciertas instrucciones o producir resultados distintos ante descripciones similares. Estas limitaciones son especialmente relevantes cuando se busca controlar un estilo visual concreto de forma consistente.

\section{Modelos de difusión y difusión latente}

Los modelos de difusión son una familia de modelos generativos basados en un proceso progresivo de eliminación de ruido. Durante la generación, el modelo parte de una representación ruidosa y la transforma paso a paso hasta obtener una imagen coherente. Este enfoque se ha convertido en una de las bases principales de muchos sistemas actuales de generación visual, especialmente por su capacidad para producir imágenes de alta calidad y por su flexibilidad para ser condicionados mediante texto.

\subsection{Funcionamiento general de los modelos de difusión}

De forma general, el entrenamiento de un modelo de difusión puede entenderse a partir de dos procesos complementarios. En primer lugar, existe un proceso directo en el que se añade ruido de forma progresiva a una imagen real hasta que la información visual queda prácticamente destruida. En segundo lugar, el modelo aprende el proceso inverso: eliminar ese ruido paso a paso para reconstruir una imagen coherente.

\medskip

Durante la generación, el modelo parte de una muestra de ruido aleatorio en lugar de una imagen real. A partir de esa representación inicial, aplica sucesivos pasos de eliminación de ruido hasta obtener una imagen final. En los sistemas texto-imagen, este proceso se condiciona mediante el \textit{prompt}, de forma que la imagen generada debe ser coherente con la descripción textual introducida por el usuario.

\medskip

Este mecanismo permite generar imágenes diversas y de gran calidad, pero también implica un coste computacional elevado. A diferencia de otros modelos que generan una salida en una única pasada, los modelos de difusión suelen requerir múltiples pasos de inferencia, lo que aumenta el tiempo de generación y el uso de recursos.

\subsection{Difusión latente y reducción del coste computacional}

Una de las principales dificultades de aplicar modelos de difusión directamente sobre imágenes es el gran tamaño del espacio de píxeles. Una imagen de alta resolución contiene una cantidad elevada de información, por lo que entrenar y generar directamente en ese espacio puede requerir una gran cantidad de memoria y capacidad de cómputo.

\medskip

Los modelos de difusión latente abordan este problema trasladando el proceso de difusión desde el espacio de píxeles a un espacio latente comprimido. Para ello, se utiliza normalmente un autoencoder que codifica la imagen original en una representación más compacta. El modelo de difusión trabaja sobre esa representación latente y, una vez completado el proceso de generación, la imagen se reconstruye de nuevo en el espacio de píxeles mediante un decodificador.

\medskip

El trabajo de Rombach et al. sobre \textit{Latent Diffusion Models} fue especialmente relevante en este sentido, ya que mostró que era posible generar imágenes de alta resolución con menor coste computacional al realizar el proceso de difusión en el espacio latente de un autoencoder previamente entrenado \cite{rombach2022high}. Esta idea permite entender por qué muchos sistemas modernos parten de modelos preentrenados y posteriormente los adaptan a dominios o estilos específicos.

\section{Adaptación de modelos preentrenados}

El entrenamiento de modelos generativos desde cero requiere grandes cantidades de datos, una elevada capacidad computacional y tiempos de entrenamiento prolongados. En el caso de los modelos texto-imagen, este coste es todavía más elevado, ya que el sistema debe aprender representaciones visuales y también la relación entre imágenes y descripciones textuales. Por este motivo, una estrategia habitual consiste en partir de modelos preentrenados, que ya han aprendido representaciones generales a partir de grandes conjuntos de datos, y adaptarlos posteriormente a una tarea, dominio o estilo concreto.

\medskip

En generación de imágenes, esta adaptación permite aprovechar las capacidades generales del modelo base y especializarlo con un conjunto de imágenes más reducido. De esta forma, el modelo conserva gran parte del conocimiento aprendido previamente, pero ajusta su comportamiento para responder mejor a un objetivo específico, como puede ser la generación de imágenes con una estética determinada.

\subsection{Fine-tuning}

El \textit{fine-tuning} consiste en continuar el entrenamiento de un modelo preentrenado utilizando un nuevo conjunto de datos más específico. En lugar de aprender desde cero, el modelo parte de unos pesos ya entrenados y los modifica para adaptarse al nuevo dominio. Esta técnica permite especializar modelos generales en tareas concretas, reduciendo el coste respecto a un entrenamiento completo.

\medskip

En modelos texto-imagen, el \textit{fine-tuning} se ha utilizado para personalizar modelos de difusión a partir de un número limitado de imágenes. Un ejemplo representativo es DreamBooth, que adapta un modelo texto-imagen preentrenado para asociar un identificador textual a un sujeto concreto y generar nuevas imágenes de ese sujeto en distintos contextos \cite{ruiz2023dreambooth}. Aunque este enfoque está centrado principalmente en sujetos específicos, ilustra la utilidad del ajuste fino para especializar modelos generativos a partir de datos concretos.

\medskip

No obstante, ajustar directamente todos los parámetros de un modelo grande puede resultar costoso y poco práctico. Al modificar una gran cantidad de pesos, el entrenamiento requiere más memoria, más tiempo de cómputo y un control más cuidadoso del proceso. Además, cuando el conjunto de entrenamiento es reducido, existe el riesgo de que el modelo memorice los ejemplos disponibles y pierda capacidad de generalización. Por esta razón, han surgido técnicas de adaptación más eficientes que permiten modificar el comportamiento del modelo sin reentrenarlo por completo.

\subsection{LoRA}

LoRA, abreviatura de \textit{Low-Rank Adaptation}, es una técnica de adaptación eficiente que permite ajustar modelos preentrenados entrenando únicamente un conjunto reducido de parámetros adicionales. En lugar de modificar todos los pesos del modelo base, LoRA congela dichos pesos e introduce pequeñas matrices entrenables de bajo rango en determinadas capas del modelo \cite{hu2022lora}.

\medskip

La idea puede entenderse a partir de una transformación lineal sencilla. En una red neuronal, muchas capas aplican una operación del tipo:

\[
y = Wx
\]

donde \(x\) representa la entrada, \(W\) es la matriz de pesos de la capa e \(y\) es la salida. En un \textit{fine-tuning} completo, el entrenamiento modificaría directamente la matriz \(W\). LoRA propone una estrategia distinta: mantener \(W\) congelada y aprender únicamente una corrección adicional \(\Delta W\). De esta forma, la operación pasa a expresarse como:

\[
y = (W + \Delta W)x
\]

\medskip

La clave de LoRA es que esa corrección \(\Delta W\) se entrena como el producto de dos matrices más pequeñas, en lugar de utilizar una matriz completa del mismo tamaño que \(W\):

\[
\Delta W = BA
\]

Si la matriz original \(W\) tiene tamaño \(k \times d\), entrenar una actualización completa implicaría aprender \(k \cdot d\) parámetros. En cambio, LoRA descompone la actualización en dos matrices de menor tamaño: una matriz \(A\) de tamaño \(r \times d\) y una matriz \(B\) de tamaño \(k \times r\), donde \(r\) es el rango de la adaptación y normalmente cumple que \(r \ll \min(k,d)\). Por tanto, el número de parámetros entrenables pasa de ser \(k \cdot d\) a:

\[
r(k+d)
\]

\medskip

Esta reducción explica por qué LoRA resulta mucho más eficiente que un \textit{fine-tuning} completo. A modo ilustrativo, si una matriz \(W\) tuviera tamaño \(4096 \times 4096\), entrenarla por completo implicaría modificar más de dieciséis millones de parámetros. Con una adaptación LoRA de rango \(r = 16\), el número de parámetros entrenables sería \(16(4096+4096)\), es decir, 131072 parámetros. Aunque se trata de un ejemplo simplificado, permite observar la diferencia de escala entre ambos enfoques.

\medskip

En la práctica, LoRA suele incorporar también un factor de escalado que controla la intensidad de la adaptación aprendida:

\[
W' = W + \frac{\alpha}{r}BA
\]

donde \(\alpha\) regula la influencia de la actualización de bajo rango sobre el modelo base \cite{hu2022lora}. En este trabajo no se realiza un estudio específico de este parámetro, pero resulta útil mencionarlo para comprender que la adaptación LoRA añade una modificación controlada sobre determinadas capas.

\medskip

Las ventajas de LoRA son especialmente relevantes en proyectos con recursos limitados. Reduce el consumo de memoria y el coste computacional, genera archivos de adaptación independientes y facilita la experimentación por iteraciones sin duplicar el modelo generativo completo. En este proyecto, LoRA se utiliza como mecanismo de adaptación del modelo base al estilo pixel-art, activado mediante el uso coherente de una palabra clave o \textit{trigger word} en las captions y en los prompts de generación.

\section{El pixel-art como estilo visual específico}

El pixel-art es un estilo de arte digital basado en la construcción de imágenes mediante píxeles visibles, normalmente asociado a resoluciones reducidas, formas simplificadas y paletas de color limitadas. A diferencia de otros estilos digitales que buscan ocultar la estructura de la imagen mediante suavizado, antialiasing o alta resolución, en el pixel-art el píxel forma parte esencial de la expresión visual.

\medskip

Desde una perspectiva visual, el pixel-art no se define únicamente por presentar una imagen de baja resolución, ya que utiliza el píxel como unidad expresiva principal. Sus características suelen relacionarse con la visibilidad de la cuadrícula, el uso controlado de la paleta de color, la simplificación de formas, la construcción clara de sprites y la búsqueda de legibilidad en resoluciones reducidas. En el ámbito de los videojuegos independientes, estas decisiones visuales no responden solo a una cuestión técnica, ya que también funcionan como una elección estética capaz de definir la identidad gráfica de un juego \cite{kylmaaho2019pixelgraphics}.


\medskip

Saravanan y Guzdial señalan que gran parte de los trabajos de \textit{machine learning} aplicados a imágenes se han centrado en imágenes realistas, dejando en segundo plano estilos más específicos como el pixel-art. Además, indican que muchos modelos tradicionales no funcionan igual de bien en este dominio, ya que suelen trabajar con grupos de píxeles, mientras que en el pixel-art cada píxel individual puede tener importancia visual \cite{saravanan2022pixelvqvae}. Esta idea también aparece en trabajos de procesamiento gráfico centrados en pixel-art, donde se destaca la importancia de las características a escala de píxel y de los contornos definidos \cite{kopf2011depixelizing}.

\medskip

Por este motivo, generar imágenes en estilo pixel-art no consiste únicamente en obtener una imagen de apariencia retro o aplicar un efecto pixelado. Un resultado puede tener una cuadrícula visible y, aun así, no respetar correctamente el estilo si presenta degradados suaves, exceso de detalle, contornos poco definidos o una paleta incoherente. Para lograr resultados consistentes, es necesario trabajar con datos representativos del estilo y adaptar el modelo para que aprenda sus restricciones visuales.

\section{Datasets y captions para entrenamiento}

En los modelos texto-imagen, el entrenamiento se basa en la relación entre imágenes y descripciones textuales. Por ello, la calidad del dataset depende de la coherencia entre cada imagen y su caption. Un conjunto de datos poco representativo, con imágenes de baja calidad o descripciones incorrectas, puede hacer que el modelo aprenda asociaciones imprecisas y produzca resultados poco consistentes.

\subsection{Construcción de datasets para adaptación estilística}

La construcción de un dataset para adaptar un modelo generativo a un estilo visual concreto requiere equilibrar dos objetivos. Por un lado, las imágenes deben compartir rasgos comunes que permitan al modelo identificar el estilo que se desea aprender. Por otro lado, el conjunto debe contener suficiente diversidad de sujetos, composiciones y colores para evitar que el modelo memorice ejemplos concretos o produzca resultados repetitivos.

\medskip

En el caso del pixel-art, la selección del dataset resulta especialmente importante porque pequeñas diferencias visuales pueden afectar al aprendizaje del modelo. Imágenes con antialiasing, degradados suaves, exceso de detalle o estilos demasiado diferentes pueden introducir señales contradictorias durante el entrenamiento. Por este motivo, recopilar imágenes etiquetadas de forma genérica como pixel-art no resulta suficiente; también es necesario aplicar criterios de filtrado y revisión que aseguren la coherencia visual del conjunto.

\medskip

Entre los criterios más relevantes se encuentran la presencia de píxeles claramente definidos, bordes duros, formas reconocibles, paletas relativamente limitadas y sombreado por bloques. También resulta importante evitar imágenes corruptas, animaciones no adecuadas para el entrenamiento, resoluciones inconsistentes o ejemplos que mezclen pixel-art con ilustración digital suavizada.

\subsection{Captions y condicionamiento textual}

Las descripciones textuales, o \textit{captions}, actúan como vínculo entre el contenido visual y el lenguaje natural. Durante el entrenamiento, el modelo aprende a asociar los elementos presentes en la imagen con las palabras utilizadas para describirla. Por este motivo, captions demasiado genéricos, incompletos o incorrectos pueden reducir la capacidad del modelo para responder de forma precisa a los prompts.

\medskip

La generación automática de captions mediante modelos \textit{vision-language} permite acelerar la creación de datasets, especialmente cuando el número de imágenes aumenta. Sin embargo, estos captions también deben revisarse o filtrarse, ya que pueden contener errores, omitir elementos relevantes o describir la imagen con un nivel de detalle inadecuado. Trabajos como BLIP señalan la importancia de mejorar la calidad de los pares imagen-texto mediante generación y filtrado de captions \cite{li2022blip}.

\medskip

En un entrenamiento orientado a pixel-art, los captions pueden describir tanto el contenido de la imagen como rasgos visuales del estilo. Por ejemplo, pueden indicar si la imagen representa un personaje, un objeto o una escena, la perspectiva utilizada, la composición general y características como \textit{limited palette}, \textit{hard edges}, \textit{flat colors} o \textit{cluster shading}. El nivel de detalle debe ajustarse al tamaño y objetivo del dataset: si las descripciones son demasiado simples, pueden no aportar suficiente información; si son excesivamente específicas, pueden favorecer la memorización de ejemplos concretos.

\medskip

En entrenamientos LoRA orientados a un estilo concreto, también es habitual utilizar una palabra clave o \textit{trigger word} que aparece de forma consistente en las captions y posteriormente en los prompts de generación. En este proyecto se utiliza el término \textit{pixelart} como señal textual asociada al estilo objetivo. Su utilidad depende de que exista coherencia entre imagen, caption y estilo: si las imágenes no representan correctamente pixel-art o las descripciones contienen información contradictoria, la asociación aprendida por el modelo puede debilitarse.

\section{Estado del arte: generación de pixel-art con modelos actuales}

Para contextualizar el problema desde una perspectiva práctica, se realizó una prueba sencilla con plataformas actuales de generación de imágenes accesibles desde la web. La intención fue reproducir una situación habitual: un usuario que quiere generar rápidamente una imagen en estilo pixel-art utilizando herramientas comunes y un prompt directo, sin plantear un estudio comparativo exhaustivo ni evaluar internamente los modelos de cada plataforma.

\medskip

La prueba se realizó con Pollo.ai, Leonardo.Ai y Magnific, utilizando el siguiente prompt: \textit{pixelart, a knight sprite, front view, limited color palette, hard edges, no antialiasing, retro game style}. Se emplearon las opciones disponibles por defecto en cada plataforma, sin realizar un ajuste avanzado de parámetros ni generar múltiples variantes, debido a las limitaciones de uso gratuito de estas herramientas. Por tanto, los resultados deben entenderse como una observación cualitativa exploratoria y no como una evaluación sistemática.

\medskip

En la Figura~\ref{fig:estado_arte_modelos_generalistas} se muestran los resultados obtenidos. Las tres imágenes representan correctamente la idea general de un caballero con armadura, pero presentan diferencias importantes en la fidelidad al estilo pixel-art.

\begin{figure}[htbp]
    \centering
    \includegraphics[width=0.31\textwidth]{imagenes/pollo.jpg}
    \hfill
    \includegraphics[width=0.31\textwidth]{imagenes/leonardo.jpg}
    \hfill
    \includegraphics[width=0.31\textwidth]{imagenes/magnifik.jpeg}
    \caption{Resultados generados por plataformas generalistas ante un prompt orientado a pixel-art. De izquierda a derecha: Pollo.ai, Leonardo.Ai y Magnific. Fuente: generaciones propias obtenidas con Pollo.ai, Leonardo.Ai y Magnific a partir del prompt indicado.}
    \label{fig:estado_arte_modelos_generalistas}
\end{figure}

\subsection{Prueba con plataformas generalistas}

Los resultados muestran que las plataformas generalistas son capaces de interpretar parcialmente la petición de pixel-art. En las imágenes generadas aparecen rasgos relacionados con el estilo, como siluetas simplificadas, contornos marcados, estética retro y cierta estructura basada en bloques. Sin embargo, también se observan limitaciones comunes que alejan los resultados de un pixel-art más estricto.

\medskip

El principal problema es la falta de control estilístico. En algunos casos, el resultado se aproxima más a una ilustración digital pixelada que a un sprite diseñado bajo restricciones claras de pixel-art. Esto se aprecia en el exceso de detalle, la presencia de brillos complejos, el uso de demasiados tonos intermedios, la aparición de transiciones suaves y la falta de una paleta realmente limitada. También se observan diferencias notables entre plataformas: algunas salidas se acercan más a un sprite simplificado, mientras que otras generan una imagen más cercana a una ilustración fantástica suavizada.

\medskip

Por tanto, esta prueba muestra que el término \textit{pixelart} puede activar una estética retro o pixelada, pero no garantiza por sí solo una construcción visual basada en píxeles definidos, bordes duros, ausencia de antialiasing y sombreado por bloques. Las imágenes pueden resultar atractivas y representar correctamente el contenido del prompt, pero no siempre mantienen de forma consistente las restricciones visuales propias del estilo.

\subsection{Limitaciones observadas y necesidad de adaptación}

Las limitaciones observadas justifican la necesidad de un proceso de adaptación más controlado. El objetivo del proyecto es comprobar si un proceso específico de preparación de datos y entrenamiento puede mejorar la fidelidad al pixel-art, sin plantearlo como una competición directa con plataformas generalistas de generación de imágenes.

\medskip

La hipótesis del trabajo se apoya precisamente en esta idea: un modelo base puede generar imágenes atractivas, pero para obtener resultados más coherentes con un estilo visual específico es necesario adaptar el modelo mediante un pipeline diseñado para ese propósito. Por ello, el sistema desarrollado integra la selección y preparación de imágenes, la generación de captions, el uso de un \textit{trigger word} y el entrenamiento de una adaptación LoRA orientada al estilo pixel-art.



\chapter{Objetivos} \label{objetivos}

\section{Objetivo general}

El objetivo general de este Trabajo de Fin de Grado es diseñar, implementar y evaluar un pipeline reproducible para adaptar un modelo generativo de difusión preentrenado a la generación de imágenes en estilo pixel-art mediante entrenamiento LoRA, datasets específicos y captions coherentes.

\section{Objetivos específicos}

Para alcanzar este objetivo general, se definen los siguientes objetivos específicos:

\begin{enumerate}
    \item Analizar los fundamentos de la generación de imágenes a partir de texto, los modelos de difusión y las técnicas de adaptación de modelos preentrenados, con especial atención al uso de LoRA.

    \item Definir los criterios visuales y técnicos que debe cumplir un dataset orientado al entrenamiento de modelos generativos en estilo pixel-art.

    \item Diseñar e implementar un flujo de recopilación, filtrado, revisión y organización iterativa de imágenes para entrenamiento.

    \item Incorporar un sistema de generación y gestión de captions que permita asociar descripciones textuales coherentes al contenido visual y al estilo pixel-art.

    \item Configurar y ejecutar entrenamientos LoRA sobre un modelo generativo preentrenado, comparando distintas iteraciones de dataset, captions y configuración de entrenamiento.

    \item Desarrollar una interfaz web que centralice la gestión de datasets, captions, entrenamientos, iteraciones y generación de imágenes.

    \item Evaluar cualitativamente los resultados obtenidos, identificando mejoras, limitaciones y posibles líneas de evolución del pipeline desarrollado.
\end{enumerate}


\chapter{Análisis y especificación} \label{analisis_especificacion}

En este capítulo se definen los requisitos del sistema desarrollado, tomando como referencia una adaptación simplificada de la estructura de especificación de requisitos software propuesta por la norma IEEE 830. El objetivo es describir de forma clara qué debe permitir realizar el pipeline, qué restricciones debe cumplir y qué características no funcionales son necesarias para garantizar su utilidad dentro del contexto del Trabajo de Fin de Grado.

\medskip

El sistema propuesto integra varias fases relacionadas con la preparación del dataset, la generación y revisión de captions, la configuración de entrenamientos, la gestión de iteraciones y la generación de imágenes mediante los modelos obtenidos, evitando quedar reducido a la ejecución aislada de un entrenamiento LoRA. Por tanto, los requisitos se han definido teniendo en cuenta tanto la parte experimental del proyecto como la necesidad de mantener un flujo reproducible y organizado.


\section{Usuarios y contexto de uso}

El sistema está diseñado principalmente para un usuario técnico o investigador que necesita experimentar con la adaptación de modelos generativos a un estilo visual concreto. En el contexto de este TFG, el usuario principal es el propio desarrollador del proyecto, que utiliza el pipeline para construir datasets, generar captions, entrenar LoRAs y evaluar cualitativamente los resultados obtenidos.

\medskip

El usuario necesita conocimientos básicos sobre modelos generativos, gestión de datasets y configuración de entrenamientos, pero la interfaz web busca reducir la necesidad de modificar manualmente archivos, ejecutar comandos dispersos o localizar resultados en distintas carpetas del proyecto. Aunque el sistema no se plantea como una aplicación comercial final, su estructura permite que pueda ser ampliado en el futuro para otros estilos visuales, otros modelos base o nuevos flujos de entrenamiento.

\section{Requisitos funcionales}

Los requisitos funcionales describen las acciones que el sistema debe permitir realizar. En la Tabla~\ref{tab:requisitos_funcionales} se recogen los requisitos principales del pipeline desarrollado.

\begingroup
\estilotabla

\begin{longtable}{|C{1.1cm}|L{3.0cm}|L{7.0cm}|C{1.5cm}|}
\caption{Requisitos funcionales del sistema.}
\label{tab:requisitos_funcionales}\\
\hline
\textbf{ID} & \textbf{Requisito} & \textbf{Descripción} & \textbf{Prioridad} \\
\hline
\endfirsthead

\hline
\textbf{ID} & \textbf{Requisito} & \textbf{Descripción} & \textbf{Prioridad} \\
\hline
\endhead

RF-01 &
Gestión de datasets &
El sistema debe organizar los datasets utilizados en el proyecto, separando imágenes, captions y resultados por iteraciones de entrenamiento. &
Alta \\
\hline

RF-02 &
Importación de imágenes &
El sistema debe incorporar imágenes al flujo de trabajo desde carpetas locales o desde procesos de recopilación automatizada. &
Alta \\
\hline

RF-03 &
Filtrado y revisión de imágenes &
El sistema debe facilitar la revisión de imágenes recopiladas para descartar ejemplos corruptos, repetidos o poco adecuados para el estilo pixel-art. &
Alta \\
\hline

RF-04 &
Preprocesamiento de imágenes &
El sistema debe preparar las imágenes en un formato adecuado para el entrenamiento, manteniendo la estética pixel-art y evitando suavizados no deseados. &
Alta \\
\hline

RF-05 &
Gestión de captions &
El sistema debe asociar captions a las imágenes del dataset y mantener la correspondencia entre cada archivo de imagen y su descripción textual. &
Alta \\
\hline

RF-06 &
Generación automática de captions &
El sistema debe generar captions de forma automática o asistida mediante un modelo vision-language, con el objetivo de acelerar la preparación del dataset. &
Media \\
\hline

RF-07 &
Revisión y edición de captions &
El sistema debe permitir revisar, modificar y corregir los captions generados antes de utilizarlos en el entrenamiento. &
Alta \\
\hline

RF-08 &
Configuración de trigger word &
El sistema debe permitir configurar una palabra clave asociada al estilo y aplicarla de forma consistente en captions y prompts de generación. &
Alta \\
\hline

RF-09 &
Configuración de entrenamientos &
El sistema debe permitir definir los parámetros principales de entrenamiento, como dataset utilizado, número de pasos, modelo base, nombre de la iteración y configuración LoRA. &
Alta \\
\hline

RF-10 &
Ejecución de entrenamientos &
El sistema debe permitir lanzar entrenamientos LoRA sobre el modelo generativo preentrenado, tanto en entorno local como en la nube cuando sea necesario. &
Alta \\
\hline

RF-11 &
Conservación del seguimiento &
El sistema debe conservar logs, configuraciones y muestras generadas durante el entrenamiento para su revisión posterior. &
Media \\
\hline

RF-12 &
Almacenamiento de artefactos &
El sistema debe guardar los resultados generados durante el entrenamiento, incluyendo el LoRA final, checkpoints, configuraciones utilizadas, logs y muestras intermedias. &
Alta \\
\hline

RF-13 &
Gestión de iteraciones &
El sistema debe conservar la trazabilidad entre datasets, captions, configuraciones y modelos generados en cada iteración. &
Alta \\
\hline

RF-14 &
Generación de imágenes &
El sistema debe permitir generar imágenes utilizando el modelo base o las adaptaciones LoRA entrenadas, a partir de prompts definidos por el usuario. &
Alta \\
\hline

RF-15 &
Organización para comparación &
El sistema debe organizar las imágenes generadas por iteración y conservar los prompts utilizados para permitir una comparación cualitativa entre resultados. &
Media \\
\hline
RF-16 & Evaluación de imágenes generadas &
El sistema debe permitir analizar imágenes generadas mediante criterios o métricas orientadas a valorar su fidelidad al estilo pixel-art. &
Media \\
\hline
\end{longtable}
\endgroup

\section{Requisitos no funcionales}

Los requisitos no funcionales describen las propiedades de calidad que debe cumplir el sistema. En este proyecto son especialmente importantes la reproducibilidad, la organización de datos, la mantenibilidad y la viabilidad computacional.

La Tabla~\ref{tab:requisitos_no_funcionales} recoge los requisitos no funcionales definidos para el sistema.

\begingroup
\estilotabla

\begin{longtable}{|C{1.1cm}|L{3.0cm}|L{7.0cm}|C{1.5cm}|}
\caption{Requisitos no funcionales del sistema.}
\label{tab:requisitos_no_funcionales}\\
\hline
\textbf{ID} & \textbf{Requisito} & \textbf{Descripción} & \textbf{Prioridad} \\
\hline
\endfirsthead

\hline
\textbf{ID} & \textbf{Requisito} & \textbf{Descripción} & \textbf{Prioridad} \\
\hline
\endhead

RNF-01 &
Reproducibilidad &
El sistema debe conservar la información necesaria para repetir o revisar cada entrenamiento, incluyendo dataset, captions, configuración y resultados asociados. &
Alta \\
\hline

RNF-02 &
Trazabilidad &
Cada iteración debe estar vinculada a sus datos de entrada, parámetros de entrenamiento y artefactos generados. &
Alta \\
\hline

RNF-03 &
Modularidad &
El sistema debe organizarse en componentes diferenciados, separando la gestión de datos, captioning, entrenamiento, generación e interfaz web. &
Alta \\
\hline

RNF-04 &
Mantenibilidad &
El código y la estructura del proyecto deben permitir realizar cambios o ampliaciones sin rehacer por completo el pipeline. &
Media \\
\hline

RNF-05 &
Usabilidad &
La interfaz web debe facilitar el uso de las funciones principales del pipeline, reduciendo la necesidad de ejecutar manualmente scripts o modificar archivos de configuración. &
Media \\
\hline

RNF-06 &
Eficiencia computacional &
El sistema debe priorizar técnicas de entrenamiento eficientes, como LoRA y configuraciones adaptadas a los recursos disponibles. &
Alta \\
\hline

RNF-07 &
Escalabilidad &
El pipeline debe poder ampliarse progresivamente para trabajar con datasets de mayor tamaño y nuevas iteraciones de entrenamiento. &
Media \\
\hline

RNF-08 &
Compatibilidad &
El sistema debe ser compatible con el entorno de desarrollo utilizado, las dependencias de entrenamiento, el modelo base seleccionado y la ejecución en GPU local o en servicios de nube previstos. &
Alta \\
\hline

RNF-09 &
Integridad de los datos &
El sistema debe evitar la pérdida o sobrescritura accidental de imágenes, captions, configuraciones y resultados generados. &
Alta \\
\hline

RNF-10 &
Gestión de errores &
El sistema debe registrar errores relevantes durante los procesos de captioning, entrenamiento o generación para facilitar su diagnóstico y corrección. &
Media \\
\hline

RNF-11 &
Seguridad básica &
Las credenciales, tokens o claves necesarias para acceder a modelos o plataformas externas deben gestionarse fuera del código fuente principal. &
Media \\
\hline

RNF-12 &
Evaluabilidad de resultados &
El sistema debe conservar las imágenes generadas y la información asociada a cada iteración para permitir una evaluación cualitativa de la fidelidad al pixel-art, la limpieza visual y la coherencia con el prompt. &
Alta \\
\hline

\end{longtable}
\endgroup


\section{Trazabilidad entre objetivos y requisitos}

Para comprobar que los requisitos definidos están alineados con los objetivos del proyecto, la Tabla~\ref{tab:trazabilidad_requisitos} relaciona cada objetivo específico con los requisitos funcionales y no funcionales más relevantes.

\begin{center}
\estilotabla
\begin{tabular}{|C{2.4cm}|L{9.8cm}|}
\hline
\textbf{Objetivo} & \textbf{Requisitos relacionados} \\
\hline
Objetivo 1 & RF-09, RF-10, RNF-06, RNF-08. \\
\hline
Objetivo 2 & RF-01, RF-03, RF-04, RNF-01, RNF-02. \\
\hline
Objetivo 3 & RF-02, RF-03, RF-04, RF-13, RNF-02, RNF-09. \\
\hline
Objetivo 4 & RF-05, RF-06, RF-07, RF-08, RNF-10. \\
\hline
Objetivo 5 & RF-09, RF-10, RF-11, RF-12, RF-13, RNF-01, RNF-06. \\
\hline
Objetivo 6 & RF-01, RF-05, RF-09, RF-14, RNF-03, RNF-05. \\
\hline
Objetivo 7 & RF-11, RF-12, RF-15, RF-16, RNF-12. \\
\hline
\end{tabular}
\captionof{table}{Trazabilidad entre objetivos específicos y requisitos del sistema.}
\label{tab:trazabilidad_requisitos}
\end{center}



\chapter{Diseño} \label{disenyo}

El diseño del sistema se ha planteado con el objetivo de convertir un conjunto de procesos técnicos, inicialmente dependientes de scripts, configuraciones y carpetas, en un flujo más organizado y trazable. Para ello, se combinan una interfaz web adaptada, una base de datos ligera, el sistema de archivos, scripts de procesamiento y el motor de entrenamiento LoRA utilizado en el proyecto.

\section{Planteamiento general del diseño}

El sistema desarrollado parte de una base existente asociada a \texttt{ai-toolkit}, utilizada como entorno de entrenamiento y gestión de modelos. Sobre esta base se ha incorporado una capa propia de adaptación para cubrir las necesidades específicas de este Trabajo de Fin de Grado.

\medskip

La aportación principal del diseño se encuentra en la integración de las fases que componen el pipeline: recopilación de imágenes, revisión y filtrado, creación de grupos de captions, generación automática de descripciones, revisión humana, preparación de datasets, configuración de entrenamientos LoRA, ejecución local o remota, seguimiento de resultados y generación final de imágenes.

\medskip

El entrenamiento de un LoRA requiere mantener una relación coherente entre los datos de entrada, las captions, los parámetros de entrenamiento, los modelos generados y las imágenes producidas posteriormente. Por ello, el sistema se plantea como un pipeline en el que cada fase produce artefactos que pueden revisarse, almacenarse y reutilizarse en fases posteriores.

\medskip

A nivel estructural, la interfaz web actúa como punto de control para el usuario, mientras que el backend coordina la base de datos, el sistema de archivos y los procesos externos. La base de datos almacena información operativa, como trabajos, colas y ajustes globales. Los artefactos pesados del proyecto, como imágenes, captions, modelos, logs y resultados, se conservan en carpetas del sistema de archivos. En la práctica, esta separación evita almacenar imágenes o modelos pesados en la base de datos y permite inspeccionar manualmente los resultados de cada iteración.

\medskip

Otro aspecto relevante del diseño es el soporte para entrenamiento iterativo. El sistema permite utilizar un LoRA generado en una iteración anterior como punto de partida de un nuevo entrenamiento. Esta funcionalidad convierte las distintas pruebas en un proceso progresivo, en el que los resultados de una fase pueden reutilizarse y compararse con los de fases posteriores.

\medskip

En conjunto, el diseño facilita la experimentación con modelos generativos aplicados al pixel-art, reduciendo la dispersión de archivos y procesos manuales. La interfaz adaptada integra el motor de entrenamiento dentro de un flujo más amplio, orientado a mantener la trazabilidad entre datasets, captions, entrenamientos y resultados.

\section{Arquitectura seleccionada}

La arquitectura seleccionada organiza el sistema en bloques diferenciados: interfaz, backend, persistencia, servicios del pipeline, capa de entrenamiento y modelos generados. Esta separación permite coordinar las tareas principales del proyecto y asignar a cada componente una responsabilidad concreta dentro del flujo de trabajo.

\medskip

La interfaz web centraliza la interacción con el usuario. Desde ella se gestionan operaciones como la revisión de imágenes, la preparación de captions, la configuración de entrenamientos y la generación de imágenes. El backend actúa como núcleo de coordinación, conectando las acciones de la interfaz con la base de datos, el sistema de archivos y los procesos externos.

\medskip

El sistema combina dos mecanismos de persistencia. La información estructurada relacionada con trabajos, colas o configuraciones se gestiona mediante base de datos. Los datasets, captions, grupos de captions, modelos entrenados, imágenes generadas y artefactos de salida se conservan en el sistema de archivos. Esta decisión resulta adecuada para un proyecto experimental en el que gran parte de la información relevante se materializa como archivos.

\subsection{Arquitectura general del sistema}

La Figura~\ref{fig:arquitectura_general} muestra la arquitectura general del sistema desarrollado. La interacción parte del usuario, que accede a una interfaz web adaptada desde la que se gestionan las operaciones principales del proyecto. Esta interfaz permite centralizar tareas que, de otro modo, quedarían distribuidas entre scripts, carpetas y configuraciones manuales.

\begin{figure}[H]
\centering
\includegraphics[width=0.95\textwidth]{imagenes/Diagrama1.png}
\caption{Arquitectura general del sistema adaptado, mostrando la relación entre interfaz, backend, persistencia, servicios del pipeline, entrenamiento y modelos generados. Fuente: elaboración propia.}
\label{fig:arquitectura_general}
\end{figure}

\medskip

La capa de API y servicios de backend desempeña el papel de núcleo coordinador. Canaliza las peticiones de la interfaz y las conecta con los servicios internos, la persistencia y la capa de entrenamiento. Esta capa permite iniciar procesos, consultar estados, acceder a artefactos y mantener organizada la información asociada a cada fase del pipeline.

\medskip

En el bloque de persistencia se distinguen dos componentes. SQLite, gestionado mediante Prisma, almacena información estructurada como \textit{jobs}, colas y configuraciones. El sistema de archivos conserva datasets, captions, grupos de captions, \textit{outputs}, logs y modelos generados. Esta división permite combinar datos estructurados con artefactos inspeccionables directamente desde el proyecto.

\medskip

Los servicios del pipeline agrupan las operaciones especializadas relacionadas con la obtención de datos, el captioning automático mediante Qwen y la generación de imágenes. Estos servicios automatizan tareas repetitivas y conectan la preparación de datos con el entrenamiento y el uso posterior de los modelos.

\medskip

La capa de entrenamiento se apoya en \textit{ai-toolkit} como herramienta principal para la ejecución de entrenamientos LoRA. Esta capa contempla la ejecución local y la ejecución en la nube mediante Modal, según los recursos disponibles. El entrenamiento parte de un modelo base FLUX y produce adaptadores LoRA, que pueden utilizarse para generar imágenes o como punto de partida de nuevas iteraciones.

\medskip

La reutilización de LoRAs generados es una característica importante de la arquitectura. Un modelo obtenido en una iteración puede incorporarse de nuevo al flujo de trabajo como LoRA previo, permitiendo continuar el proceso de ajuste a partir de resultados anteriores.

\subsection{Flujo funcional del pipeline}

La Figura~\ref{fig:flujo_pipeline} representa el flujo funcional del sistema desde la preparación inicial de los datos hasta la generación final de imágenes. Este diagrama muestra el recorrido que siguen los datos y los modelos a lo largo del pipeline.

\begin{figure}[H]
\centering
\includegraphics[width=0.95\textwidth]{imagenes/Diagrama2.png}
\caption{Flujo funcional del pipeline desde la obtención de imágenes hasta la generación final mediante un LoRA entrenado. Fuente: elaboración propia.}
\label{fig:flujo_pipeline}
\end{figure}

\medskip

El flujo comienza con la preparación del dataset. En esta primera fase se obtienen imágenes mediante procesos de recopilación o scraping, se aplican filtros automáticos y se realiza una revisión manual para descartar ejemplos poco adecuados. Las imágenes aceptadas pasan posteriormente a un grupo de captions, donde se generan descripciones automáticas y se revisan o editan antes de construir el dataset final de entrenamiento.

\medskip

Una vez preparado el dataset, este se utiliza como entrada para la configuración del entrenamiento. En dicha configuración se definen los elementos necesarios para lanzar el proceso LoRA, como el dataset seleccionado, el modelo base, la palabra clave asociada al estilo, los parámetros principales del entrenamiento y, de forma opcional, un LoRA previo. Esta opción permite reutilizar una adaptación generada en una iteración anterior dentro de una nueva fase de entrenamiento.

\medskip

El entrenamiento puede ejecutarse en local o mediante Modal, dependiendo de los recursos disponibles y de las necesidades del experimento. Ambas opciones conducen a la ejecución del entrenamiento LoRA, que produce un nuevo modelo LoRA y un conjunto de artefactos asociados. Entre estos artefactos se encuentran los checkpoints, los logs, el archivo de configuración y las muestras generadas durante el propio entrenamiento. Estas muestras permiten revisar visualmente la evolución del proceso.

\medskip

El LoRA generado puede utilizarse posteriormente en la fase de generación de imágenes. En esta fase, el modelo seleccionado y el prompt introducido por el usuario se combinan para producir nuevas imágenes en estilo pixel-art. Las imágenes generadas se almacenan y organizan en función del modelo o LoRA utilizado, lo que facilita la comparación entre distintas iteraciones y mantiene una trazabilidad básica entre el modelo empleado y los resultados obtenidos.

\medskip

El flujo refleja el carácter iterativo del sistema. Un LoRA obtenido puede reutilizarse en nuevos entrenamientos, compararse con versiones anteriores o emplearse para generar imágenes adicionales. Esta organización permite trabajar de forma progresiva, ajustando datos, captions y configuraciones a partir de los resultados obtenidos en cada iteración.

\section{Tecnologías empleadas}
\label{sec:tecnologias_empleadas}

El sistema combina tecnologías web, herramientas de procesamiento de datos y librerías de inteligencia artificial debido a la diversidad de operaciones que forman el pipeline. La solución debe gestionar imágenes y metadatos, supervisar la preparación de datasets, generar captions, configurar entrenamientos y utilizar los modelos resultantes desde una misma interfaz. La aplicación web actúa como capa de coordinación sobre distintos servicios y scripts especializados.

\medskip

La Tabla~\ref{tab:tecnologias_diseno} recoge las principales tecnologías utilizadas y su función dentro del diseño. Se consideran tecnologías centrales aquellas que sostienen las capas principales del sistema, mientras que las tecnologías auxiliares proporcionan funciones concretas de procesamiento, optimización o serialización.
\begingroup
\estilotabla
\begin{longtable}{|L{0.18\textwidth}|L{0.27\textwidth}|L{0.43\textwidth}|}
\caption{Tecnologías empleadas en el sistema.}
\label{tab:tecnologias_diseno}\\
\hline
\textbf{Área} & \textbf{Tecnologías} & \textbf{Uso dentro del sistema} \\
\hline
\endfirsthead

\hline
\textbf{Área} & \textbf{Tecnologías} & \textbf{Uso dentro del sistema} \\
\hline
\endhead
Interfaz web &
Next.js, React, TypeScript y Tailwind CSS &
Construcción y adaptación de las pantallas destinadas a gestionar el scraping, los captions, los datasets, los entrenamientos y la generación de imágenes. \\
\hline
Backend e integración &
Node.js y API Routes de Next.js &
Exposición de operaciones internas y coordinación entre la interfaz, la base de datos, el sistema de archivos y los scripts de Python. \\
\hline
Persistencia &
Prisma y SQLite &
Almacenamiento estructurado de trabajos de entrenamiento, estado de las colas y ajustes generales de la aplicación. \\
\hline
Archivos y formatos &
Sistema de archivos, JSON, JSONL, CSV, TXT, YAML, PNG y safetensors &
Almacenamiento de imágenes, metadatos, captions, configuraciones, datasets, registros, modelos y artefactos de entrenamiento. \\
\hline
Obtención y preparación de datos &
Python, Scrapy, Pillow y Sharp &
Obtención automática de imágenes mediante topics, filtrado, análisis de propiedades y preparación de las imágenes para los datasets. \\
\hline
Captioning automático &
Qwen2.5-VL, Transformers, PyTorch, Pillow y bitsandbytes &
Carga del modelo multimodal, procesamiento de las imágenes y generación automática de captions con optimización del uso de memoria. \\
\hline
Entrenamiento LoRA &
ai-toolkit, PyTorch, Diffusers, PEFT, bitsandbytes y safetensors &
Integración del motor de entrenamiento, ejecución de operaciones sobre modelos de difusión y almacenamiento de los pesos LoRA resultantes. \\
\hline
Ejecución &
Python, CUDA y Modal &
Ejecución de los procesos de inteligencia artificial sobre GPU, tanto en el entorno local como en infraestructura cloud. \\
\hline
Modelos y generación &
FLUX.1-dev, LoRA, Diffusers y PyTorch &
Uso del modelo base y de los adaptadores entrenados para generar nuevas imágenes desde la aplicación. \\
\hline
\end{longtable}
\endgroup

\medskip

Como herramientas auxiliares de desarrollo se utilizaron Visual Studio Code y Git para la edición del código y el control de versiones, aunque no forman parte de la arquitectura funcional del pipeline.

\medskip

\textbf{Interfaz web.}
La interacción con el sistema se realiza mediante una \textbf{interfaz web adaptada} a partir de la base proporcionada por \texttt{ai-toolkit}. Esta base se ha ampliado con pantallas y componentes específicos para las distintas fases del pipeline. La interfaz está desarrollada con \textbf{Next.js} y \textbf{React}, mientras que \textbf{TypeScript} permite definir de forma explícita las estructuras intercambiadas entre los componentes y los servicios. \textbf{Tailwind CSS} se emplea como tecnología auxiliar para mantener una presentación coherente entre las vistas.

\medskip

\textbf{Backend y API.}
Next.js proporciona la capa de backend mediante sus \textbf{API Routes}, ejecutadas sobre \textbf{Node.js}. Estas rutas constituyen la capa propia de integración del sistema: reciben las acciones realizadas desde la interfaz, validan los datos y coordinan el acceso a la persistencia, los archivos y los procesos externos. Gracias a esta capa, la aplicación web puede iniciar un scraper, gestionar grupos de captions, preparar datasets, lanzar entrenamientos o solicitar la generación de imágenes desde un mismo entorno.

\medskip

\textbf{Persistencia.}
La información estructurada de la aplicación se almacena mediante \textbf{SQLite}, utilizando \textbf{Prisma} como capa de acceso. Esta persistencia contiene principalmente los trabajos de entrenamiento, su configuración y estado, las colas de ejecución y los ajustes generales. Prisma permite que la capa de servicios trabaje con estas entidades mediante una interfaz tipada y evita que la lógica de la aplicación dependa directamente de consultas SQL.

\medskip

\textbf{Sistema de archivos.}
La base de datos se complementa con el \textbf{sistema de archivos}, que desempeña un papel central debido al volumen y naturaleza de los artefactos gestionados. Las imágenes recopiladas, los grupos de captions, los datasets, los logs, las configuraciones, los checkpoints y los modelos entrenados se conservan en directorios organizados por fase e iteración.

Se utilizan formatos diferentes según el tipo de información. \textbf{JSON} y \textbf{JSONL} almacenan manifiestos y metadatos; \textbf{CSV} y archivos \textbf{TXT} representan los captions; \textbf{YAML} describe configuraciones de entrenamiento; \textbf{PNG} se emplea para las imágenes; y \textbf{safetensors} contiene los pesos de los adaptadores LoRA. Estos formatos facilitan tanto la lectura por parte de los scripts como la inspección y trazabilidad de los datos.

\medskip

\textbf{Scraping y preparación de datos.}
La obtención automática de imágenes se implementa en \textbf{Python} mediante \textbf{Scrapy}. El scraper consulta las galerías de Pixilart tomando como punto de partida distintos \textbf{topics}, permitiendo distribuir la recolección entre categorías y controlar el número de elementos obtenido de cada una.

Durante esta fase, \textbf{Pillow} se utiliza para abrir las imágenes descargadas y extraer propiedades necesarias para su validación. Posteriormente, la aplicación permite revisar las imágenes aceptadas y organizar su incorporación a los grupos de captions. En la preparación final de los datasets también se utiliza \textbf{Sharp} desde la capa Node.js, principalmente para transformar y escalar las imágenes conservando el tipo de interpolación adecuado para pixel-art. Scrapy constituye la tecnología central de recolección, mientras que Pillow y Sharp actúan como herramientas auxiliares de procesamiento.

\medskip

\textbf{Captioning automático.}
La generación automática de captions utiliza el modelo multimodal \textbf{Qwen2.5-VL-7B-Instruct}. El modelo recibe cada imagen junto con una instrucción diseñada para producir descripciones adecuadas para el entrenamiento LoRA. Su integración se realiza mediante \textbf{Transformers} y \textbf{PyTorch}, mientras que Pillow prepara las imágenes antes de introducirlas en el modelo.

La librería \textbf{bitsandbytes} permite cargar el modelo con cuantización, reduciendo el consumo de memoria durante el captioning. También se emplean utilidades específicas del ecosistema Qwen para adaptar la información visual al formato esperado por el modelo. Los captions resultantes se almacenan en CSV y quedan disponibles en la interfaz para su revisión y normalización antes de incorporarlos al dataset.

\medskip

\textbf{Entrenamiento LoRA.}
El entrenamiento se apoya en \texttt{ai-toolkit}, que actúa como \textbf{motor de entrenamiento integrado} en la aplicación. El sistema construye las configuraciones necesarias, inicia los trabajos y recupera los artefactos resultantes.

La ejecución se sustenta principalmente en \textbf{PyTorch} y \textbf{Diffusers}, responsables de las operaciones sobre el modelo de difusión. La técnica \textbf{LoRA} permite adaptar el modelo base entrenando un conjunto reducido de parámetros. Librerías como \textbf{bitsandbytes} proporcionan optimizadores de menor consumo de memoria, mientras que \textbf{PEFT} forma parte del soporte auxiliar para adaptadores eficientes. Los pesos resultantes se almacenan en formato \textbf{safetensors}.

\medskip

\textbf{Ejecución local y cloud.}
Los procesos de entrenamiento pueden ejecutarse utilizando una GPU local compatible con \textbf{CUDA} o trasladarse a la infraestructura cloud proporcionada por \textbf{Modal}. CUDA permite que PyTorch realice las operaciones intensivas sobre GPU, mientras que Modal encapsula el entorno de entrenamiento y proporciona acceso a recursos gráficos remotos.

Desde el punto de vista del diseño, Modal constituye una capa de ejecución alternativa integrada en el mismo pipeline. La configuración de entrenamiento se adapta al entorno remoto, los datasets y modelos necesarios se transfieren al contenedor y los resultados se recuperan posteriormente en la estructura local del proyecto.

\medskip

\textbf{Generación de imágenes.}
El modelo base utilizado por el sistema es \textbf{FLUX.1-dev}. Sobre este modelo pueden cargarse los adaptadores LoRA obtenidos durante los entrenamientos para incorporar el estilo aprendido. La generación se implementa mediante \textbf{Diffusers} y \textbf{PyTorch}, utilizando un pipeline específico de FLUX.

La capa de generación permite seleccionar el modelo base o uno de los LoRA disponibles, introducir un prompt y, opcionalmente, establecer una semilla. La imagen generada se almacena en el sistema de archivos y se devuelve a la interfaz web para su visualización. De esta manera, el sistema conecta la preparación de datos, el entrenamiento y el uso posterior de los modelos obtenidos dentro de un flujo común.

\section{Modelo de entidades y trazabilidad}

Además de definir el flujo funcional del pipeline, resulta necesario identificar qué elementos permiten conservar la relación entre los datos de entrada, los procesos ejecutados y los resultados obtenidos. En este proyecto, la trazabilidad se reparte entre la estructura de carpetas, los manifiestos generados, los archivos de captions, las configuraciones de entrenamiento y los artefactos producidos durante cada ejecución.

\medskip

El modelo de entidades representa los elementos principales del pipeline y las relaciones que permiten reconstruir la procedencia de los datos y resultados. Esta visión resulta adecuada para el sistema desarrollado, ya que muchas de sus entidades se materializan como archivos, directorios, registros JSONL, configuraciones YAML o modelos almacenados en formato \texttt{safetensors}.

\medskip

La trazabilidad del sistema se organiza en tres fases principales: la preparación del dataset, el entrenamiento del modelo LoRA y la generación de imágenes. Cada una de estas fases conserva información distinta y aporta un nivel diferente de seguimiento sobre el proceso completo. En los siguientes apartados se describe cómo se mantiene esta relación entre imágenes recopiladas, captions, datasets, entrenamientos, modelos y resultados generados.

\subsection{Trazabilidad en la preparación del dataset}

La Figura~\ref{fig:trazabilidad_dataset} muestra la trazabilidad asociada a la preparación del dataset. Esta fase comienza con una ejecución del scraper, representada mediante un manifiesto de ejecución que conserva información sobre el proceso de recopilación. A partir de esta ejecución se generan registros de imágenes recopiladas, donde se almacenan datos asociados a cada imagen, como su procedencia, estado de filtrado y ruta local.

\medskip

Cada imagen recopilada puede quedar registrada como rechazada o pasar al conjunto de imágenes aceptadas. Cuando una imagen se descarta, el sistema conserva información sobre el motivo del rechazo. Cuando se acepta, se copia a las carpetas correspondientes y puede incorporarse a grupos de captions. Esta separación permite conservar tanto los ejemplos utilizados como los descartados durante la preparación del dataset.

\medskip

Las imágenes aceptadas se incorporan posteriormente a un grupo de captions. Este grupo actúa como una unidad intermedia de trabajo, en la que se reúnen las imágenes seleccionadas y se preparan sus descripciones textuales. El grupo conserva su propio manifiesto y permite mantener la relación entre las imágenes revisadas, los captions generados automáticamente y las posibles modificaciones realizadas antes de la incorporación al dataset final.

\medskip

Durante esta fase también se utilizan metadatos de imagen, como el título, el autor, las etiquetas o el tamaño del lienzo. Esta información aporta contexto adicional para la generación automática de descripciones y permite trabajar con datos procedentes del proceso de scraping.

\medskip

El resultado de esta fase es un dataset organizado en una estructura compatible con el entrenamiento. Cada dataset contiene las imágenes aceptadas, los captions en formato \texttt{.txt} asociados a cada archivo y un fichero \texttt{captions.csv}. Esta organización facilita el acceso a los pares imagen-texto y conserva una relación trazable entre la imagen original recopilada, su caption y el dataset en el que finalmente se utiliza.

\clearpage

\begin{figure}[H]
\centering
\includegraphics[height=0.78\textheight]{imagenes/diagrama3-1.png}
\caption{Trazabilidad en la preparación del dataset desde la ejecución del scraper hasta la construcción del dataset final. Fuente: elaboración propia.}
\label{fig:trazabilidad_dataset}
\end{figure}

\clearpage

\subsection{Trazabilidad del entrenamiento}

La Figura~\ref{fig:trazabilidad_entrenamiento} representa la trazabilidad asociada a la fase de entrenamiento. Esta parte del sistema conecta el dataset preparado previamente con la configuración del entrenamiento, la ejecución concreta del proceso LoRA, el modelo resultante y los artefactos generados durante la ejecución.

\medskip

El punto de partida de esta fase es el dataset final, que contiene las imágenes y captions preparados para el entrenamiento. Este dataset se selecciona dentro de la configuración del trabajo de entrenamiento, quedando vinculado al proceso que se va a ejecutar. Esta relación permite identificar qué conjunto de imágenes y descripciones se ha utilizado para obtener un determinado LoRA.

\medskip

El entrenamiento se inicia a partir de un \textit{TrainingJob}, que representa la solicitud de entrenamiento dentro del sistema. Esta entidad almacena información estructurada sobre el trabajo, como su nombre, estado, posición en cola y configuración asociada. La configuración concreta del entrenamiento se conserva en el \textit{JobConfig}, donde se recogen elementos como la ruta del dataset, el modelo base, los parámetros principales y, en caso necesario, la ruta de un LoRA previo.

\medskip

La posibilidad de indicar un LoRA previo permite encadenar entrenamientos entre iteraciones. El sistema puede partir de una adaptación generada anteriormente y utilizarla como base para un nuevo entrenamiento. Esta relación queda reflejada mediante el parámetro correspondiente en la configuración del trabajo.

\medskip

En este diseño, \textit{TrainingRun} se utiliza como entidad lógica de trazabilidad. La ejecución concreta se reconstruye a partir del trabajo configurado, la carpeta de salida, los archivos de configuración y los logs generados. Esta representación permite distinguir la solicitud del entrenamiento, la configuración aplicada y los resultados producidos.

\medskip

El resultado principal de esta ejecución es un modelo LoRA almacenado como archivo \texttt{.safetensors}. Este archivo constituye la adaptación entrenada sobre el modelo base y puede utilizarse posteriormente para generar imágenes o como punto de partida de nuevas iteraciones. Además del LoRA final, el entrenamiento produce distintos artefactos de salida, como muestras intermedias, checkpoints, logs y archivos de configuración. Estos elementos permiten revisar el comportamiento del entrenamiento, analizar su evolución y conservar evidencias del proceso ejecutado.

\medskip

Esta estructura permite reconstruir la relación entre un LoRA generado, el entrenamiento que lo produjo, la configuración utilizada y el dataset de origen. Algunas relaciones se apoyan en la estructura de carpetas y archivos, por lo que la trazabilidad se basa en la combinación entre datos estructurados, configuraciones conservadas y artefactos almacenados.

\clearpage

\begin{figure}[H]
\centering
\includegraphics[height=0.78\textheight]{imagenes/diagrma3-2.png}
\caption{Trazabilidad del entrenamiento desde el dataset utilizado hasta el modelo LoRA generado y sus artefactos de salida. Fuente: elaboración propia.}
\label{fig:trazabilidad_entrenamiento}
\end{figure}

\clearpage

\subsection{Trazabilidad de la generación}

La Figura~\ref{fig:trazabilidad_generacion} representa la trazabilidad asociada a la fase de generación de imágenes. Esta parte del sistema conecta el modelo utilizado, ya sea el modelo base FLUX o un LoRA entrenado previamente, con la petición de generación y con la imagen final almacenada en el sistema de archivos.

\medskip

La generación se realiza a partir de un modelo seleccionado y de un prompt introducido por el usuario. El modelo puede ser el modelo base FLUX o un LoRA entrenado previamente. Cuando se utiliza un LoRA, el sistema carga la adaptación correspondiente sobre el modelo base para incorporar el estilo aprendido durante el entrenamiento. Además del prompt, la generación puede incorporar otros parámetros, como la semilla, que permite diferenciar ejecuciones y reproducir parcialmente determinados resultados si se conserva la configuración empleada.

\medskip

Una característica relevante del diseño es que las imágenes generadas se almacenan en una carpeta asociada al modelo utilizado. De esta forma, las generaciones realizadas con el modelo base se guardan bajo la carpeta correspondiente a \texttt{flux_base}, mientras que las imágenes producidas con un LoRA se organizan dentro del directorio asociado a su \texttt{modelId}. Esta estructura permite separar los resultados por modelo y facilita la comparación visual entre distintas iteraciones de entrenamiento.

\medskip

El nombre de cada imagen generada incorpora información como la semilla y una marca temporal. Esto permite distinguir diferentes generaciones, incluso cuando se realizan con el mismo modelo o con prompts similares. La estructura de carpetas y la nomenclatura de archivos proporcionan una trazabilidad básica entre el modelo usado y la imagen resultante.

\medskip

La trazabilidad de la generación presenta ciertas limitaciones. El prompt completo utilizado durante la generación no queda persistido como metadato estructurado junto a la imagen, por lo que la petición original no siempre puede reconstruirse únicamente a partir del archivo generado. La relación entre la imagen, el modelo y la semilla se conserva principalmente mediante la ruta de salida, el nombre del archivo y la información devuelta durante la ejecución.

\clearpage

\begin{figure}[H]
\centering
\includegraphics[height=0.72\textheight]{imagenes/diagrama3-3.png}
\caption{Trazabilidad de la generación de imágenes a partir del modelo base o de un LoRA entrenado. Fuente: elaboración propia.}
\label{fig:trazabilidad_generacion}
\end{figure}

\clearpage
\section{Diseño de la interfaz adaptada}

La interfaz web adaptada actúa como punto de control de las distintas fases del pipeline. Su objetivo es centralizar tareas que, inicialmente, quedaban repartidas entre carpetas, scripts y archivos de configuración. Desde esta capa se coordinan la obtención de imágenes, la revisión del material recopilado, la preparación de captions, la construcción de datasets, la configuración de entrenamientos, el seguimiento de ejecuciones y la generación de resultados.

\medskip

El diseño parte de la interfaz base de \texttt{ai-toolkit}, orientada principalmente a la gestión de entrenamientos, colas, datasets y monitorización de la GPU. Sobre esta base se añadieron secciones específicas para el flujo de trabajo del proyecto, como la generación de imágenes con modelos seleccionables, la gestión de entrenamientos previos, la preparación de captions, el scraper y las galerías de resultados organizadas por iteración.

\medskip

La Figura~\ref{fig:interfaz_base_adaptada} compara la barra lateral de la interfaz inicial con la versión adaptada. La primera recoge las funciones generales del entorno de entrenamiento. La segunda incorpora accesos específicos del pipeline desarrollado, lo que permite recorrer de forma más directa las fases de obtención de datos, captioning, entrenamiento, generación y revisión de resultados.

\begin{figure}[H]
\centering
\includegraphics[height=0.47\textheight]{imagenes/interfaz_aitoolkit_base.png}
\hspace{0.8cm}
\includegraphics[height=0.47\textheight]{imagenes/interfaz_pipeline_adaptada.png}
\caption{Comparación entre la barra lateral de la interfaz base de \texttt{ai-toolkit} y la barra lateral de la interfaz adaptada al pipeline del proyecto. Fuente: captura de la interfaz base de \texttt{ai-toolkit} y de la interfaz adaptada en el proyecto.}
\label{fig:interfaz_base_adaptada}
\end{figure}

\medskip

La interfaz se diseña como una capa de coordinación del pipeline. Cada grupo de pantallas responde a una fase concreta del flujo: obtención y revisión de imágenes, gestión de captions y preparación del dataset, configuración del entrenamiento, seguimiento de ejecuciones y generación de imágenes. Esta separación ayuda a mantener la trazabilidad entre los artefactos generados en cada fase y reduce la necesidad de localizar manualmente archivos en distintas carpetas del proyecto.

\subsection{Obtención y revisión de imágenes}

La obtención y revisión de imágenes constituye la primera fase de trabajo dentro de la interfaz. Esta sección resuelve un problema práctico del proyecto: recopilar imágenes de distintas categorías y revisarlas antes de que formen parte del dataset de entrenamiento. La pantalla de scraping integra esta tarea en la aplicación, evitando que la recopilación dependa únicamente de la ejecución manual de scripts externos.

\medskip

En la Figura~\ref{fig:interfaz_scraper_lanzamiento} se muestra la pantalla desde la que se inicia el proceso de scraping. La vista permite definir una carpeta de destino, seleccionar la cantidad de imágenes y escoger los \textit{topics} de Pixilart utilizados como punto de partida. Esta organización hace que cada recopilación quede asociada a una carpeta identificable, lo que facilita su revisión posterior y conserva la separación entre distintas búsquedas.

\begin{figure}[htbp]
\centering
\includegraphics[width=0.95\textwidth]{imagenes/interfaz_scraper_lanzamiento.png}
\caption{Interfaz para lanzar un proceso de scraping seleccionando carpeta, cantidad de imágenes y topics de Pixilart. Fuente: captura de la interfaz desarrollada en el proyecto.}
\label{fig:interfaz_scraper_lanzamiento}
\end{figure}

\medskip

Una vez finalizada la recopilación, la interfaz ofrece una vista de revisión organizada por carpetas. Cada carpeta muestra un resumen del número de imágenes aceptadas, curadas y rechazadas, de forma que el usuario puede conocer rápidamente el estado del conjunto recopilado antes de entrar en la revisión visual.

\medskip

La revisión se plantea mediante una galería de tarjetas porque el criterio de selección es principalmente visual. En la Figura~\ref{fig:interfaz_scraper_revision} se observa la separación entre imágenes rechazadas, aceptadas y aceptadas manualmente. Cada tarjeta incluye la imagen y metadatos básicos como nombre, resolución, tamaño y etiquetas asociadas.

\begin{figure}[htbp]
\centering
\includegraphics[width=0.95\textwidth]{imagenes/interfaz_scraper_revision.png}
\caption{Vista de revisión de imágenes recopiladas, organizada por estado de filtrado para separar imágenes rechazadas, aceptadas y aceptadas manualmente. Fuente: captura de la interfaz desarrollada en el proyecto.}
\label{fig:interfaz_scraper_revision}
\end{figure}

\medskip

Esta vista actúa como filtro visual previo a la construcción del dataset. Durante la revisión se pueden detectar imágenes corruptas, poco representativas, con exceso de suavizado, con una resolución poco adecuada o alejadas del estilo pixel-art buscado. La pestaña de aceptadas manualmente introduce una fase de control humano sobre los resultados del filtrado automático, evitando que todo el material recopilado pase directamente a las fases de captioning y entrenamiento.

\subsection{Gestión de grupos de captions y preparación del dataset}

La gestión de grupos de captions funciona como una fase intermedia entre las imágenes aceptadas y el dataset final de entrenamiento. Su objetivo es agrupar un subconjunto de imágenes, generar descripciones automáticas, revisarlas y convertirlas posteriormente en una estructura compatible con el entrenamiento LoRA.

\medskip

La creación de grupos separa la revisión de imágenes de la construcción definitiva del dataset. En la Figura~\ref{fig:interfaz_caption_group_creacion} se muestra la ventana desde la que se selecciona una carpeta procedente del scraper y se escogen las imágenes que formarán parte del grupo. Esta decisión de diseño permite trabajar con subconjuntos concretos según temática, calidad visual o utilidad para una iteración determinada.

\begin{figure}[H]
\centering
\includegraphics[width=0.95\textwidth]{imagenes/interfaz_caption_group_creacion.png}
\caption{Interfaz para crear un grupo de captions a partir de una carpeta del scraper y seleccionar las imágenes que lo compondrán. Fuente: captura de la interfaz desarrollada en el proyecto.}
\label{fig:interfaz_caption_group_creacion}
\end{figure}

\medskip

Una vez creado el grupo, la generación automática de captions se configura desde una ventana específica. En esta fase se selecciona el modelo de visión, el \textit{trigger word}, el nivel de detalle y el modo de sobrescritura. Estos parámetros están pensados para producir descripciones útiles para el entrenamiento, manteniendo una relación coherente entre imagen, caption y estilo visual.

\medskip

La Figura~\ref{fig:interfaz_caption_group_generacion} recoge esta configuración. El \textit{trigger word} \texttt{pixelart} se utiliza como señal textual asociada al estilo y se añade al inicio de los captions generados. El nivel de detalle controla la extensión de las descripciones, mientras que el modo de sobrescritura evita reemplazar captions ya revisados cuando solo se quieren completar imágenes sin descripción.

\begin{figure}[H]
\centering
\includegraphics[width=0.34\textwidth]{imagenes/interfaz_caption_group_generacion.png}
\caption{Configuración de la generación automática de captions, incluyendo modelo de visión, trigger word, nivel de detalle y modo de sobrescritura. Fuente: captura de la interfaz desarrollada en el proyecto.}
\label{fig:interfaz_caption_group_generacion}
\end{figure}

\medskip

La revisión posterior de captions combina automatización y control humano. El sistema genera una primera versión de las descripciones, pero el usuario conserva la posibilidad de inspeccionar y editar el resultado antes de usarlo en el entrenamiento. Esta decisión es importante porque un caption incorrecto puede introducir asociaciones poco útiles entre la imagen y el texto.

\medskip

En la Figura~\ref{fig:interfaz_caption_group_resultado} se muestra un grupo después de generar las descripciones automáticas. Cada tarjeta reúne la imagen y su caption asociado, lo que permite comprobar si la descripción representa correctamente el contenido visual y si mantiene un nivel de detalle adecuado. Esta vista evita trabajar únicamente sobre archivos CSV o TXT y hace más directa la detección de captions demasiado genéricos, incorrectos o poco coherentes.

\begin{figure}[H]
\centering
\includegraphics[width=0.65\textwidth]{imagenes/interfaz_caption_group_resultado.png}
\caption{Vista de revisión de un grupo de captions tras la generación automática de descripciones. Fuente: captura de la interfaz desarrollada en el proyecto.}
\label{fig:interfaz_caption_group_resultado}
\end{figure}

\medskip

Cuando el grupo ya está revisado, la interfaz permite moverlo al dataset de entrenamiento seleccionado, como se muestra en la Figura~\ref{fig:interfaz_mover_dataset}. Esta operación cierra la fase de preparación de captions y convierte el grupo en datos listos para entrenamiento. En este paso se indica de nuevo el \textit{trigger word} y se selecciona un dataset de destino creado previamente desde la sección de datasets.

\medskip

Durante esta operación, las imágenes del grupo se copian al dataset seleccionado, se preparan con el tamaño requerido para el entrenamiento y se generan los archivos \texttt{.txt} asociados a cada imagen a partir de sus captions. El \textit{trigger word} vuelve a aparecer en esta fase para asegurar que los archivos finales de caption conservan la señal textual utilizada durante el entrenamiento. Así, el grupo de captions queda transformado en una estructura de pares imagen-texto compatible con el entrenamiento LoRA.

\begin{figure}[H]
\centering
\includegraphics[width=0.30\textwidth]{imagenes/interfaz_mover_dataset.png}
\caption{Ventana para mover un grupo de captions al dataset de entrenamiento seleccionado. Fuente: captura de la interfaz desarrollada en el proyecto.}
\label{fig:interfaz_mover_dataset}
\end{figure}

\subsection{Configuración de entrenamientos}

La configuración de entrenamientos conecta el dataset preparado con la ejecución del proceso LoRA. Esta pantalla concentra los parámetros necesarios para lanzar una iteración reproducible: nombre del entrenamiento, GPU utilizada, modelo base, palabra clave asociada al estilo, carpeta de salida, dataset seleccionado, número de pasos, resolución y opciones de guardado.

\medskip

La Figura~\ref{fig:interfaz_training_configuracion} presenta la pantalla de creación de un nuevo entrenamiento. Aunque la interfaz contiene un número amplio de campos, el flujo del proyecto se centra en los elementos que afectan directamente a la trazabilidad de la iteración. El nombre del entrenamiento identifica la prueba, la carpeta de salida agrupa sus artefactos, el dataset vincula el entrenamiento con los datos preparados y el modelo base define el punto de partida del proceso.

\begin{figure}[H]
\centering
\includegraphics[width=0.85\textwidth]{imagenes/interfaz_training_configuracion.png}
\caption{Pantalla de configuración de un nuevo entrenamiento LoRA, incluyendo modelo, GPU, dataset, parámetros principales y carpeta de salida. Fuente: captura de la interfaz desarrollada en el proyecto.}
\label{fig:interfaz_training_configuracion}
\end{figure}

\medskip

La selección del dataset es uno de los elementos principales de esta pantalla. El entrenamiento queda asociado a un conjunto concreto de imágenes y archivos \texttt{.txt}, lo que permite reconstruir qué datos se utilizaron para obtener cada LoRA. Esta relación es fundamental para comparar iteraciones y detectar si una mejora o un problema procede del dataset, de los captions o de la configuración de entrenamiento.

\medskip

El campo \textit{Pretrained LoRA Path} incorpora la estrategia iterativa directamente en la interfaz. Al seleccionar un LoRA previo, el usuario puede iniciar un nuevo entrenamiento a partir de una adaptación ya generada. Esta opción evita modificar manualmente configuraciones para encadenar entrenamientos y mantiene el vínculo entre versiones del modelo dentro del propio flujo de trabajo.

\medskip

La configuración también incluye opciones de guardado y control del entrenamiento. La frecuencia de guardado conserva checkpoints o versiones intermedias, mientras que la carpeta de salida agrupa el modelo resultante, los logs, las configuraciones y los artefactos asociados. Al conservar dataset, modelo base, \textit{trigger word}, pasos, resolución, salida y LoRA previo, cada entrenamiento queda asociado a una configuración concreta que puede revisarse posteriormente.

\medskip

La generación de muestras durante el entrenamiento se configura en la sección mostrada en la Figura~\ref{fig:interfaz_training_samples}. Esta parte define la frecuencia de muestreo, la resolución, la semilla y los prompts utilizados para observar la evolución visual del LoRA durante el proceso.

\begin{figure}[H]
\centering
\includegraphics[width=0.95\textwidth]{imagenes/interfaz_training_samples.png}
\caption{Configuración de muestras generadas durante el entrenamiento, incluyendo frecuencia, resolución, semilla y prompts de evaluación visual. Fuente: captura de la interfaz desarrollada en el proyecto.}
\label{fig:interfaz_training_samples}
\end{figure}

\medskip

Los prompts de muestra cumplen una función de control visual. En un proyecto centrado en estilo, revisar imágenes intermedias ayuda a detectar problemas de configuración o aprendizaje antes de analizar únicamente el resultado final. Esta vista convierte el muestreo en parte del diseño experimental del pipeline, ya que permite observar si el modelo empieza a asociar el \textit{trigger word} con rasgos propios del pixel-art.

\subsection{Seguimiento de entrenamientos}

El seguimiento de entrenamientos cubre dos necesidades distintas: controlar el estado operativo de los trabajos y revisar los artefactos generados durante la ejecución. La interfaz reúne estas dos dimensiones mediante vistas de cola, pantallas de detalle con logs y galerías de samples.

\medskip

La Figura~\ref{fig:interfaz_training_queue} muestra la vista general de la cola de entrenamientos. En ella se distinguen las colas asociadas a la ejecución local y a Modal, junto con el estado de los trabajos registrados. Esta separación resulta útil porque el proyecto combina entrenamiento local y entrenamiento en la nube, y cada entorno puede tener tiempos, recursos y estados diferentes.

\begin{figure}[H]
\centering
\includegraphics[width=0.95\textwidth]{imagenes/interfaz_training_queue.png}
\caption{Vista general de la cola de entrenamientos, con trabajos asociados a ejecución local y Modal. Fuente: captura de la interfaz desarrollada en el proyecto.}
\label{fig:interfaz_training_queue}
\end{figure}

\medskip

Cada entrenamiento dispone además de una vista de detalle. Esta pantalla agrupa el nombre del trabajo, la GPU asignada, el estado de ejecución, el progreso y los mensajes generados durante el proceso. En ejecuciones mediante Modal, los logs también recogen información sobre la ejecución remota y la descarga de artefactos.

\medskip

La Figura~\ref{fig:interfaz_training_logs} muestra un entrenamiento completado. El área de logs funciona como evidencia del proceso ejecutado y ayuda a revisar errores de configuración, incidencias de ejecución o problemas relacionados con la recuperación de resultados. Esta vista conecta directamente con los requisitos de seguimiento, gestión de errores y conservación de artefactos.

\begin{figure}[H]
\centering
\includegraphics[width=0.95\textwidth]{imagenes/interfaz_training_logs.png}
\caption{Vista de detalle de un entrenamiento, con estado de ejecución, GPU asignada y registros del proceso. Fuente: captura de la interfaz desarrollada en el proyecto.}
\label{fig:interfaz_training_logs}
\end{figure}

\medskip

El seguimiento visual se realiza mediante los samples generados durante el entrenamiento. Estas imágenes se organizan por prompt y por paso, lo que permite comparar la evolución del modelo a lo largo de la ejecución. Esta vista resulta especialmente útil cuando se busca aprender un estilo visual, porque permite comprobar si los resultados intermedios se acercan progresivamente a la estética pixel-art esperada.

\medskip

La Figura~\ref{fig:interfaz_training_samples_resultado} muestra la visualización de samples de un entrenamiento. Cada fila corresponde a un prompt de prueba y cada columna representa un paso concreto. Esta organización facilita detectar cambios en coherencia visual, composición, detalle y fidelidad al estilo antes de pasar al análisis de resultados finales.

\begin{figure}[H]
\centering
\includegraphics[width=1\textwidth]{imagenes/interfaz_training_samples_resultado.png}
\caption{Vista de samples generados durante el entrenamiento, organizados por prompt y paso de entrenamiento. Fuente: captura de la interfaz desarrollada en el proyecto.}
\label{fig:interfaz_training_samples_resultado}
\end{figure}

\medskip

La combinación de cola, logs y samples aporta una trazabilidad operativa y visual del entrenamiento. La cola informa sobre el estado de los trabajos, los logs conservan el desarrollo de la ejecución y los samples permiten revisar la evolución del modelo. Esta separación evita depender únicamente del archivo final \texttt{.safetensors} para comprender qué ocurrió durante una iteración.

\subsection{Generación de imágenes}

La generación de imágenes cierra el flujo de trabajo de la interfaz adaptada. Una vez entrenados los LoRA, el sistema permite utilizarlos desde una vista de inferencia en la que el usuario selecciona el modelo, introduce un \textit{prompt} y obtiene nuevas imágenes en estilo pixel-art.

\medskip

La selección de modelo es una decisión clave de esta pantalla. El usuario puede generar con el modelo base FLUX o con cualquiera de los LoRA producidos durante las iteraciones del proyecto. La Figura~\ref{fig:interfaz_generacion_selector} muestra el selector disponible en la interfaz, donde cada entrada identifica el entrenamiento asociado. Esta organización facilita comparar resultados entre iteraciones sin buscar manualmente los modelos en el sistema de archivos.

Además de los LoRAs catalogados en las iteraciones principales, el selector se actualiza a partir de las carpetas disponibles en \texttt{outputs/}. De esta forma, cuando un nuevo entrenamiento genera una carpeta de salida con un archivo \texttt{.safetensors}, el modelo puede aparecer en la interfaz sin tener que registrarlo manualmente en el código.

\begin{figure}[H]
\centering
\includegraphics[width=0.36\textwidth]{imagenes/interfaz_generacion_selector.png}
\caption{Selector del modelo de generación, incluyendo el modelo base FLUX y los distintos LoRA entrenados. Fuente: captura de la interfaz desarrollada en el proyecto.}
\label{fig:interfaz_generacion_selector}
\end{figure}

\medskip

Tras seleccionar el modelo, el usuario introduce un \textit{prompt} y lanza la generación. La imagen resultante aparece en la misma vista junto con información básica sobre el modelo utilizado y la semilla. La Figura~\ref{fig:interfaz_generacion_resultado} muestra una generación realizada con el modelo \textit{Iteracion 3 - Entrenamiento 5 (LoRA v4)} a partir de un prompt orientado a un perro en estilo pixel-art.

\begin{figure}[H]
\centering
\includegraphics[width=0.95\textwidth]{imagenes/interfaz_generacion_resultado.png}
\caption{Interfaz de generación de imágenes con un prompt de ejemplo y la imagen resultante obtenida con un LoRA entrenado. Fuente: captura de la interfaz desarrollada y generación propia mediante el LoRA seleccionado.}
\label{fig:interfaz_generacion_resultado}
\end{figure}

\medskip

La aplicación conserva las imágenes generadas dentro de una galería asociada al modelo seleccionado. Esta decisión responde a una necesidad de comparación: las salidas deben poder revisarse por modelo o iteración sin depender de una búsqueda manual en carpetas. La Figura~\ref{fig:interfaz_generacion_galeria} muestra una galería correspondiente a un entrenamiento concreto, donde se agrupan varias imágenes generadas desde la interfaz.

\begin{figure}[H]
\centering
\includegraphics[width=0.95\textwidth]{imagenes/interfaz_generacion_galeria.png}
\caption{Galería de imágenes generadas asociada a un entrenamiento concreto. Fuente: captura de la interfaz desarrollada en el proyecto.}
\label{fig:interfaz_generacion_galeria}
\end{figure}

\medskip

La galería incorpora una capa de evaluación visual orientada al análisis de resultados. Al seleccionar una imagen, la interfaz abre una vista ampliada junto con un panel lateral de métricas y puntuaciones. Esta evaluación se plantea como una herramienta auxiliar y automática de apoyo, no como una valoración definitiva de la calidad artística. Su objetivo es ofrecer indicadores que ayuden a revisar aspectos relevantes del pixel-art, como bordes sin suavizado, coherencia de la figura, ruido visual o compacidad de la paleta.

\medskip

La Figura~\ref{fig:interfaz_generacion_evaluacion} muestra una imagen generada junto con su panel de evaluación. Esta vista permite relacionar el resultado visual con puntuaciones concretas y detectar de forma más sistemática posibles problemas, por ejemplo exceso de suavizado, baja legibilidad o presencia de ruido que dificulte la lectura de la figura.

\begin{figure}[H]
\centering
\includegraphics[width=0.95\textwidth]{imagenes/interfaz_generacion_evaluacion.png}
\caption{Vista ampliada de una imagen generada junto con su panel de evaluación y análisis visual. Fuente: captura de la interfaz desarrollada en el proyecto.}
\label{fig:interfaz_generacion_evaluacion}
\end{figure}

\medskip

El diseño de la generación conecta el uso práctico de los modelos con su revisión posterior. El selector de modelo, el prompt, la imagen generada, la galería por entrenamiento y el panel de evaluación forman un flujo orientado a experimentar con los LoRA entrenados y analizar sus resultados sin romper la trazabilidad con la iteración de origen.

\chapter{Implementación} \label{implementacion}

Este capítulo describe el proceso de implementación del sistema desarrollado, explicando cómo se construyó el pipeline de generación de imágenes pixel-art a partir de una evolución por iteraciones. Mientras que el capítulo anterior se centraba en el diseño general del sistema, en este apartado se detallan las decisiones técnicas, los cambios introducidos en cada fase, los scripts utilizados, las configuraciones de entrenamiento, los problemas encontrados y las soluciones aplicadas.

\medskip

La implementación del proyecto se realizó como un proceso progresivo, en lugar de plantearse como un desarrollo cerrado desde el inicio. En primer lugar, se validó el entrenamiento LoRA con un dataset reducido y un flujo local. Después, se mejoró la calidad de los datos y se introdujo el encadenamiento entre LoRAs. Finalmente, se amplió el sistema hasta convertirlo en un pipeline más completo, incorporando interfaz web, scraping, generación automática de captions, entrenamiento en la nube, generación de imágenes y evaluación visual.

\medskip

Esta organización permite explicar la evolución real del proyecto y justificar por qué cada cambio fue necesario. Cada iteración parte de las limitaciones detectadas en la anterior y añade nuevas funcionalidades orientadas a mejorar la calidad del dataset, la reproducibilidad del entrenamiento y la trazabilidad de los resultados.

\section{Enfoque de implementación}

El enfoque de implementación adoptado se basa en una estrategia incremental. En lugar de construir todo el sistema completo desde el principio, se desarrollaron primero los elementos mínimos necesarios para comprobar la viabilidad técnica del entrenamiento LoRA y, a partir de esa base, se añadieron mejoras de forma progresiva.

\medskip

Este planteamiento resultaba adecuado para el proyecto por varios motivos. En primer lugar, el entrenamiento de modelos generativos depende de muchos factores difíciles de ajustar al inicio, como la calidad del dataset, los captions, la configuración LoRA, los recursos de GPU y la compatibilidad entre entrenamiento e inferencia. Trabajar por iteraciones permitió validar cada parte del pipeline antes de ampliar su complejidad.

\medskip

En segundo lugar, el proyecto combina componentes muy distintos: preparación de imágenes, gestión de captions, configuración de entrenamientos, ejecución local y en la nube, generación de imágenes y evaluación de resultados. La implementación progresiva permitió integrar estos bloques de forma controlada, manteniendo siempre una versión funcional del sistema.

\medskip

Por último, esta metodología facilitó la documentación del proceso. Cada iteración generó datasets, configuraciones, modelos LoRA, samples, logs e imágenes de prueba que podían compararse con los resultados anteriores. De este modo, la implementación produjo un sistema final y una evolución trazable del trabajo realizado.

\subsection{Metodología incremental}

La implementación del proyecto se organizó mediante una metodología incremental, dividiendo el desarrollo en varias iteraciones con objetivos concretos. Esta forma de trabajo permitió construir primero una versión mínima funcional del pipeline y, posteriormente, ampliar sus capacidades a partir de las limitaciones detectadas en cada fase.

\medskip

Este enfoque resulta adecuado para un proyecto experimental basado en entrenamiento de modelos generativos, ya que no todos los problemas pueden preverse al inicio. La calidad del dataset, el comportamiento de los captions, la configuración LoRA, el consumo de memoria, la compatibilidad entre entrenamiento e inferencia y la calidad visual de los resultados son aspectos que se van ajustando a medida que se realizan pruebas reales.

\medskip

Cada iteración se planteó con un objetivo principal y una serie de pasos orientados a alcanzarlo. La primera iteración permitió validar el funcionamiento básico del entrenamiento LoRA. La segunda se centró en mejorar el modelo y comprobar el encadenamiento entre entrenamientos. La tercera amplió el alcance del sistema mediante automatización, interfaz web, entrenamiento en la nube y evaluación de imágenes generadas.

\medskip

La Figura~\ref{fig:diagrama_iteraciones_implementacion} resume esta evolución. En ella se muestra cómo cada iteración parte de un objetivo concreto y sigue una serie de pasos que permiten avanzar desde una validación inicial hasta un pipeline más completo y automatizado.

\begin{figure}
\centering
\includegraphics[width=0.95\textwidth]{imagenes/diagrama_iteraciones_implementacion.png}
\caption{Resumen de la implementación por iteraciones, indicando el objetivo y los pasos principales de cada fase. Fuente: elaboración propia.}
\label{fig:diagrama_iteraciones_implementacion}
\end{figure}

\medskip

Esta estructura incremental también facilitó la toma de decisiones durante el desarrollo. En lugar de modificar todo el sistema a la vez, cada iteración introdujo cambios controlados sobre una base ya validada. De esta manera fue posible comparar resultados, detectar errores de configuración, conservar artefactos intermedios y justificar la evolución técnica del pipeline.


\subsection{Organización del trabajo por iteraciones}

Para mantener la trazabilidad del proyecto, el trabajo se organizó separando datasets, configuraciones, modelos y resultados por iteraciones. Esta decisión fue importante porque cada entrenamiento dependía de un conjunto concreto de imágenes, captions y parámetros. Sin una organización clara, habría sido difícil saber qué datos habían producido cada LoRA o qué configuración estaba asociada a cada resultado.

\medskip

En las primeras fases, esta organización se apoyaba principalmente en carpetas del repositorio. Cada iteración mantenía sus imágenes, captions, archivos de configuración y salidas de entrenamiento. Posteriormente, con la incorporación de la interfaz web, esta estructura se reforzó mediante vistas específicas para datasets, grupos de captions, entrenamientos, colas, galerías de imágenes generadas y evaluación de resultados.

\medskip

La organización por iteraciones también facilitó la comparación entre modelos. El LoRA de la primera iteración sirvió como punto de partida para la segunda, y los modelos generados en fases posteriores pudieron utilizarse desde la interfaz para generar nuevas imágenes y revisar diferencias visuales. Además, los artefactos asociados a cada entrenamiento, como checkpoints, logs, samples y configuraciones, permitieron reconstruir el proceso seguido en cada caso.

\medskip

En conjunto, esta forma de trabajo permitió convertir un conjunto de pruebas experimentales en un pipeline progresivo y documentado. Las siguientes secciones describen con más detalle cómo se implementó cada iteración, qué cambios se introdujeron y qué papel tuvo cada una en la evolución final del sistema.


\section{Iteración 1: validación inicial del entrenamiento LoRA}
\label{sec:iteracion_1}

La primera iteración tuvo como objetivo comprobar la viabilidad técnica de adaptar el modelo \textbf{FLUX.1-dev} al estilo pixel-art mediante un entrenamiento LoRA. Antes de ampliar el dataset o desarrollar mecanismos de automatización, era necesario validar que el pipeline básico podía ejecutarse utilizando los recursos disponibles y producir un adaptador funcional.

\medskip

Para ello, se preparó un entorno de entrenamiento local, se construyó un dataset reducido de imágenes seleccionadas manualmente y se definió una configuración inicial de 500 pasos. Esta iteración permitió comprobar la carga del modelo base, la lectura de pares imagen-caption, el entrenamiento sobre GPU y la generación de checkpoints y muestras intermedias.

\medskip

El alcance de esta fase fue deliberadamente limitado. Se pretendía establecer una primera línea base sobre la que desarrollar las siguientes iteraciones del proyecto, sin buscar todavía un modelo definitivo.

\subsection{Preparación del entorno de desarrollo}
\label{subsec:entorno_iteracion_1}

La primera tarea consistió en preparar un entorno local capaz de ejecutar \texttt{ai-toolkit} y cargar el modelo FLUX.1-dev. El desarrollo se realizó sobre Windows, utilizando PowerShell para ejecutar los scripts y administrar el entorno de Python. En esta primera iteración, tanto la preparación de los datos como el entrenamiento se realizaron directamente en el equipo local.

\medskip

Para aislar las dependencias del proyecto se creó un \textbf{entorno virtual de Python} en la carpeta \texttt{venv/}. Este entorno permitió mantener separadas las librerías requeridas por el proyecto de aquellas instaladas globalmente en el equipo, evitando posibles conflictos de compatibilidad. El archivo \texttt{venv/pyvenv.cfg} indica que el entorno fue creado utilizando \textbf{Python 3.10.9}.

\medskip

La activación se realizaba desde PowerShell mediante el comando \texttt{.\textbackslash venv\textbackslash Scripts\textbackslash Activate.ps1}. Como se observa en la Figura~\ref{fig:entorno_virtual_iteracion_1}, la aparición del identificador \texttt{(venv)} en la terminal confirmaba que los comandos posteriores utilizarían el entorno aislado del proyecto.

\begin{figure}[H]
    \centering
    \includegraphics[width=0.8\textwidth]{imagenes/iteracion1_entorno_virtual.png}
    \caption{Activación del entorno virtual de Python utilizado en la primera iteración. Fuente: captura del entorno de desarrollo local utilizado en el proyecto.}
    \label{fig:entorno_virtual_iteracion_1}
\end{figure}

\medskip

El motor de entrenamiento empleado fue \texttt{ai-toolkit}, utilizado como base para configurar y ejecutar el entrenamiento LoRA. Las dependencias necesarias se gestionaron mediante el archivo \texttt{requirements.txt}. La Tabla~\ref{tab:dependencias_iteracion_1} recoge las principales librerías relacionadas con esta primera ejecución.

\begin{table}[htbp]
\centering
\estilotabla
\begin{tabular}{|L{0.30\textwidth}|L{0.58\textwidth}|}
\hline
\textbf{Dependencia} & \textbf{Utilidad en el proyecto} \\
\hline
Python & Ejecución de \texttt{ai-toolkit} y de los scripts auxiliares. \\
\hline
PyTorch & Entrenamiento del modelo y utilización de la GPU mediante CUDA. \\
\hline
Diffusers & Carga y gestión de los componentes del modelo FLUX. \\
\hline
Transformers & Procesamiento de los componentes textuales del modelo. \\
\hline
Accelerate & Gestión del dispositivo y apoyo a la ejecución del entrenamiento. \\
\hline
bitsandbytes & Utilización del optimizador AdamW de 8 bits. \\
\hline
safetensors & Almacenamiento de los pesos del LoRA y sus checkpoints. \\
\hline
PyYAML & Lectura de la configuración del entrenamiento. \\
\hline
python-dotenv & Carga de variables privadas del entorno. \\
\hline
\end{tabular}
\caption{Dependencias principales utilizadas en la primera iteración.}
\label{tab:dependencias_iteracion_1}
\end{table}

El modelo base seleccionado fue \texttt{black-forest-labs/FLUX.1-dev}. Para acceder a él fue necesario aceptar sus condiciones en Hugging Face y proporcionar un token de lectura. Este token se cargaba mediante una variable de entorno y no se almacenaba en el repositorio.

\medskip

También se adaptó el proceso de descarga al entorno Windows. En \texttt{run.py} se deshabilitó la transferencia acelerada de Hugging Face mediante la variable \texttt{HF\_HUB\_ENABLE\_HF\_TRANSFER}, evitando los problemas de bloqueo de archivos encontrados durante la preparación del entorno.

\medskip

El entrenamiento se ejecutó mediante CUDA sobre una GPU NVIDIA local con 16 GB de VRAM. La ejecución local permitió comprobar directamente la compatibilidad de las dependencias, el consumo de memoria y el funcionamiento del pipeline antes de incorporar una infraestructura de entrenamiento remoto.

\subsection{Construcción del dataset inicial}
\label{subsec:dataset_iteracion_1}

El dataset inicial se construyó mediante una selección manual de una cantidad reducida de imágenes representativas de pixel-art. Se escogieron principalmente personajes, animales, vehículos, objetos e iconos con una resolución visual baja, píxeles claramente distinguibles y una apariencia próxima a los sprites empleados en videojuegos clásicos.

\medskip

La selección priorizó imágenes con formas reconocibles, colores relativamente simples y bordes definidos. El reducido tamaño del conjunto respondía al carácter exploratorio de la iteración: se buscaba verificar el entrenamiento con el menor coste computacional posible antes de invertir esfuerzo en la construcción de un dataset de mayor escala.

\medskip

Originalmente, las imágenes se almacenaron en \texttt{datasets/pixel\_art/}. Durante una reorganización posterior del repositorio, el conjunto se trasladó a \texttt{datasets/iter\_01/images/}, conservando las imágenes y captions empleados en el entrenamiento.

\medskip

La Figura~\ref{fig:dataset_inicial_ejemplos} muestra dos ejemplos del conjunto inicial. Ambas imágenes presentan sujetos aislados, formas reconocibles y píxeles claramente visibles, características representativas del material utilizado en esta primera prueba.

\begin{figure}[H]
    \centering
    \begin{minipage}{0.50\textwidth}
        \centering
        \includegraphics[width=0.42\textwidth]{imagenes/Bunny.png}
        \hspace{0.08\textwidth}
        \includegraphics[width=0.42\textwidth]{imagenes/Lion.png}
    \end{minipage}
    \caption{Ejemplos de imágenes pertenecientes al dataset inicial. De izquierda a derecha: Bunny y Lion. Fuente: imágenes recopiladas manualmente para el dataset inicial del proyecto.}
    \label{fig:dataset_inicial_ejemplos}
\end{figure}
\subsection{Preparación de imágenes y captions}
\label{subsec:preparacion_iteracion_1}

Las imágenes seleccionadas ya presentaban las características necesarias para el primer entrenamiento, por lo que su preparación se centró principalmente en asociar cada una con una descripción textual compatible con el formato esperado por \texttt{ai-toolkit}.

\medskip

Cada imagen se acompañó de un archivo de texto con el mismo nombre base. Por ejemplo, \texttt{Bunny.png} se asociaba con \texttt{Bunny.txt}. La opción \texttt{caption\_ext: txt} del archivo de configuración indicaba al motor que debía utilizar estos archivos como captions durante el entrenamiento.

\medskip

Los captions se redactaron en inglés con ayuda de \textbf{ChatGPT} y se almacenaron manualmente en los archivos de texto correspondientes. Se utilizaron descripciones breves que identificaban el sujeto y algunas características generales de la imagen. Un ejemplo es el siguiente:

\begin{quote}
\texttt{pixelart, blocky bunny sprite, side view, simple colors, retro game style}
\end{quote}

Todos los captions comenzaban con la palabra \texttt{pixelart}, utilizada como \textbf{trigger word}. Este término permitía relacionar el estilo visual compartido por las imágenes con una palabra concreta que posteriormente podía incluirse en los prompts de generación.

\subsection{Configuración del primer entrenamiento LoRA}
\label{subsec:configuracion_iteracion_1}

El primer entrenamiento se definió mediante el archivo \texttt{config/iter\_01\_train.yaml}, derivado de la configuración inicial \texttt{train\_lora\_pixelart\_tfg.yaml}. El modelo base seleccionado fue \texttt{black-forest-labs/FLUX.1-dev} y el adaptador se entrenó desde cero, sin utilizar un LoRA procedente de un entrenamiento anterior.

\medskip

La Tabla~\ref{tab:parametros_iteracion_1} recoge los principales parámetros seleccionados para adaptar el entrenamiento a los recursos disponibles.

\begin{table}[htbp]
\centering
\estilotabla
\begin{tabular}{|L{0.31\textwidth}|L{0.24\textwidth}|L{0.33\textwidth}|}
\hline
\textbf{Parámetro} & \textbf{Valor} & \textbf{Finalidad} \\
\hline
Modelo base & FLUX.1-dev & Modelo adaptado mediante LoRA \\
\hline
Trigger word & \texttt{pixelart} & Activación del estilo aprendido \\
\hline
Pasos de entrenamiento & 500 & Duración de la prueba inicial \\
\hline
Resolución de entrenamiento & 512$\times$512 & Resolución procesada por el modelo. \\
\hline
Batch size & 1 & Reducción del consumo de VRAM \\
\hline
Rank y alpha LoRA & 16 / 16 & Capacidad del adaptador \\
\hline
Learning rate & $1 \times 10^{-4}$ & Tasa de actualización de los pesos \\
\hline
Optimizador & AdamW de 8 bits & Reducción del consumo de memoria \\
\hline
Precisión & BF16 & Menor uso de memoria durante el entrenamiento \\
\hline
Noise scheduler & FlowMatch & Planificador empleado con FLUX \\
\hline
Guardado & Cada 100 pasos & Creación de checkpoints recuperables \\
\hline
Sampling & Cada 100 pasos & Seguimiento visual del entrenamiento \\
\hline
\end{tabular}
\caption{Parámetros principales del primer entrenamiento LoRA.}
\label{tab:parametros_iteracion_1}
\end{table}

\subsection{Ejecución local y artefactos generados}
\label{subsec:ejecucion_iteracion_1}

Una vez activado el entorno virtual, el entrenamiento se inició desde la raíz del repositorio mediante el comando \texttt{python run.py config/iter\_01\_train.yaml}. El archivo \texttt{run.py} cargaba la configuración y delegaba la ejecución en el motor de entrenamiento de \texttt{ai-toolkit}.

\medskip

Los resultados se almacenaron originalmente en \texttt{output/pixel\_art\_lora\_v1/}. Tras la reorganización del proyecto, los artefactos se conservaron en \texttt{outputs/iter\_01/pixel\_art\_lora\_v1/}.

La Tabla~\ref{tab:artefactos_iteracion_1} resume los principales artefactos generados durante esta primera ejecución.

\begin{table}[htbp]
\centering
\estilotabla
\begin{tabular}{|L{0.42\textwidth}|L{0.46\textwidth}|}
\hline
\textbf{Artefacto} & \textbf{Utilidad} \\
\hline
\texttt{pixel\_art\_lora\_v1.safetensors} & Adaptador LoRA final utilizado posteriormente para inferencia. \\
\hline
\texttt{pixel\_art\_lora\_v1\_000000NNN.safetensors} & Checkpoints correspondientes a los pasos 100, 200, 300 y 400. \\
\hline
\texttt{optimizer.pt} & Estado del optimizador necesario para reanudar el entrenamiento. \\
\hline
\texttt{config.yaml} & Copia de la configuración efectiva utilizada durante la ejecución. \\
\hline
\texttt{samples/} & Imágenes generadas periódicamente para observar la evolución del entrenamiento. \\
\hline
\end{tabular}
\caption{Artefactos generados durante el primer entrenamiento.}
\label{tab:artefactos_iteracion_1}
\end{table}

Durante una de las ejecuciones, el equipo se apagó automáticamente debido a la configuración de apagado del sistema operativo. La interrupción impidió que el entrenamiento finalizara en esa sesión, pero permitió comprobar el funcionamiento del mecanismo de recuperación de \texttt{ai-toolkit}.

\medskip

Al volver a iniciar el mismo entrenamiento, el sistema detectó el último checkpoint disponible y restauró los pesos del LoRA y el estado del optimizador. De esta forma, la ejecución continuó desde la última copia guardada en lugar de comenzar nuevamente desde el paso cero.

\medskip

La carpeta \texttt{samples/} almacenó imágenes generadas en diferentes puntos del entrenamiento. La Figura~\ref{fig:evolucion_samples_iteracion_1} presenta una selección organizada en una cuadrícula de 2$\times$2, permitiendo observar los cambios producidos durante el aprendizaje.
\begin{figure}[H]
    \centering
    \setlength{\tabcolsep}{8pt}
    \renewcommand{\arraystretch}{1.1}

    \begin{tabular}{cc}
        \includegraphics[width=0.27\textwidth]{imagenes/iteracion1_sample_paso_0.jpg} &
        \includegraphics[width=0.27\textwidth]{imagenes/iteracion1_sample_paso_100.jpg} \\
        \small Paso 0 & \small Paso 100 \\[0.25cm]

        \includegraphics[width=0.27\textwidth]{imagenes/iteracion1_sample_paso_300.jpg} &
        \includegraphics[width=0.27\textwidth]{imagenes/iteracion1_sample_paso_500.jpg} \\
        \small Paso 300 & \small Paso 500 \\
    \end{tabular}

    \caption{Evolución de una muestra durante el primer entrenamiento LoRA. Fuente: generación propia obtenida durante el entrenamiento LoRA.}
    \label{fig:evolucion_samples_iteracion_1}
\end{figure}

La ejecución confirmó que el entorno preparado era capaz de cargar FLUX.1-dev, entrenar un adaptador LoRA, generar muestras periódicas, conservar checkpoints y reanudar el proceso después de una interrupción. De este modo, la primera iteración cumplió su función como validación técnica del pipeline.

\medskip

Una vez comprobado su funcionamiento, la segunda iteración se orientó a ampliar el material de entrenamiento y continuar el aprendizaje a partir del LoRA obtenido. El primer entrenamiento se convirtió así en la base del enfoque iterativo seguido durante el resto del proyecto.
\section{Iteración 2: mejora incremental del modelo}
\label{sec:iteracion_2}

La segunda iteración se planteó como una evolución directa de la primera. Una vez comprobado que el entrenamiento LoRA sobre FLUX.1-dev podía ejecutarse correctamente en el entorno local, el objetivo pasó a ser mejorar la calidad del material de entrenamiento, organizar el proyecto de forma más trazable y validar si era posible continuar el aprendizaje a partir del LoRA ya obtenido.

\medskip

A diferencia de la primera iteración, esta fase no partió desde cero. Se reutilizó el entorno local ya validado y se introdujo el uso de \texttt{pretrained\_lora\_path}, permitiendo que el segundo entrenamiento comenzara desde el LoRA generado en la Iteración 1. De esta forma, la iteración se diseñó como una mejora incremental del modelo anterior, en lugar de una repetición independiente.

\subsection{Reorganización del proyecto por iteraciones}
\label{subsec:reorganizacion_iteracion_2}

Durante esta fase se reorganizó el repositorio para separar de forma explícita los datasets, configuraciones y resultados de cada entrenamiento. Esta decisión resolvía una limitación práctica de la primera iteración: al trabajar con una estructura menos diferenciada, resultaba más difícil distinguir qué imágenes, captions, configuración y salidas pertenecían a cada prueba.

\medskip

La organización por iteraciones permitió conservar una relación directa entre los datos utilizados, el archivo de configuración y los artefactos producidos por el entrenamiento. Así, la segunda iteración quedó asociada a \texttt{datasets/iter\_02/}, \texttt{config/iter\_02\_train.yaml} y \texttt{outputs/iter\_02/pixel\_art\_lora\_v2/}. Esta separación facilita la comparación entre entrenamientos y evita sobrescribir resultados anteriores.

La Figura~\ref{fig:estructura_iteracion_2} muestra la estructura de carpetas utilizada para separar el material de esta segunda iteración.

\begin{figure}[H]
\centering
\includegraphics[width=0.38\textwidth]{imagenes/iteracion2_estructura_carpetas.png}
\caption{Estructura del dataset de la segunda iteración. Fuente: captura de la estructura de carpetas del proyecto.}
\label{fig:estructura_iteracion_2}
\end{figure}

La Tabla~\ref{tab:organizacion_iteracion_2} resume la función de los principales archivos y directorios asociados a la segunda iteración.

\begin{table}[htbp]
\centering
\estilotabla
\begin{tabular}{|L{0.32\textwidth}|L{0.58\textwidth}|}
\hline
\textbf{Elemento} & \textbf{Función dentro de la iteración} \\
\hline
\texttt{datasets/iter\_02/raw/} & Imágenes originales antes del procesamiento. \\
\hline
\texttt{datasets/iter\_02/images/} & Imágenes finales procesadas a 512x512, junto con sus captions \texttt{.txt}. \\
\hline
\texttt{datasets/iter\_02/captions.csv} & Archivo principal desde el que se generaban los captions individuales por imagen. \\
\hline
\texttt{config/iter\_02\_train.yaml} & Configuración del segundo entrenamiento LoRA. \\
\hline
\texttt{outputs/iter\_02/pixel\_art\_lora\_v2/} & Carpeta de salida con el LoRA generado, checkpoints, configuración copiada y samples. \\
\hline
\end{tabular}
\caption{Organización de archivos en la segunda iteración.}
\label{tab:organizacion_iteracion_2}
\end{table}

\subsection{Mejora del dataset}
\label{subsec:dataset_iteracion_2}

El dataset de la segunda iteración se diseñó como una versión más cuidada y controlada que el conjunto inicial. En esta fase se introdujo una separación clara entre las imágenes originales y las imágenes preparadas para entrenamiento. Los archivos originales se conservaron en \texttt{datasets/iter\_02/raw/}, mientras que las versiones finales utilizadas por el entrenamiento se almacenaron en \texttt{datasets/iter\_02/images/}.

\medskip

Para preparar las imágenes se utilizó el script \texttt{scripts/tfg/scale\_pixelart.py}. Este script toma una carpeta de entrada y una carpeta de salida, carga las imágenes originales y genera versiones en formato PNG adaptadas al tamaño esperado por el entrenamiento. En el caso de esta iteración, el flujo aplicado fue equivalente a procesar las imágenes de \texttt{datasets/iter\_02/raw/} y guardar el resultado en \texttt{datasets/iter\_02/images/}.

\medskip

El script utiliza la biblioteca Pillow para abrir las imágenes, convertirlas a RGB y escalarlas hasta una resolución de 512x512 píxeles. Para mantener la estética pixel-art, el reescalado se realiza mediante nearest-neighbor, evitando interpolaciones suaves que podrían difuminar los bordes o generar colores intermedios. Cuando la imagen original no encaja exactamente en el formato cuadrado, se coloca sobre un lienzo de 512x512, manteniendo la proporción y rellenando el espacio restante.

\medskip

Este procesamiento permitió unificar el formato de entrada del dataset antes del entrenamiento. A diferencia de la primera iteración, donde las imágenes ya se encontraban preparadas para el flujo inicial, en esta segunda fase se incorporó explícitamente una etapa de preparación de imágenes. Esto hacía que el dataset fuese más homogéneo y reducía problemas derivados de trabajar con formatos o tamaños distintos.

\medskip

El conjunto final quedó formado por un número reducido de imágenes, pero con mayor variedad visual que en la iteración anterior. Incluía personajes, criaturas, iconos, objetos y escenas sencillas, como un mago, un slime, una consola portátil, una isla flotante o composiciones de atardecer. La selección seguía siendo manual, pero buscaba ejemplos más representativos del estilo que se quería reforzar en el modelo.

La Figura~\ref{fig:dataset_iteracion_2} muestra dos ejemplos de imágenes incluidas en el dataset de esta segunda iteración.

\begin{figure}[H]
\centering
\includegraphics[width=0.34\textwidth]{imagenes/circularsunset.png}
\hspace{0.06\textwidth}
\includegraphics[width=0.34\textwidth]{imagenes/hero.png}
\caption{Ejemplos de imágenes utilizadas en el dataset de la segunda iteración. Fuente: imágenes recopiladas manualmente para el dataset de la segunda iteración.}
\label{fig:dataset_iteracion_2}
\end{figure}

\subsection{Creación de captions descriptivos}
\label{subsec:captions_iteracion_2}

También se mejoró la forma de preparar los captions del dataset. En esta iteración se utilizó \texttt{datasets/iter\_02/captions.csv} como archivo principal de referencia. Este fichero contenía dos columnas, \texttt{file} y \texttt{caption}, que relacionaban cada imagen con su descripción textual.

\medskip

A partir de ese CSV se generaban los archivos \texttt{.txt} individuales mediante el script \texttt{scripts/tfg/sync\_captions.py}. Este script recorre las imágenes presentes en \texttt{datasets/iter\_02/images/}, busca su entrada correspondiente en \texttt{captions.csv} y escribe un archivo de texto con el mismo nombre base que la imagen. Por ejemplo, la entrada asociada a \texttt{Slime.png} genera el archivo \texttt{Slime.txt}.

\medskip

Este paso era necesario porque la configuración de entrenamiento usa \texttt{caption\_ext: txt}. Es decir, el motor de entrenamiento no lee directamente el CSV, ya que espera encontrar un archivo de texto junto a cada imagen. El CSV actúa como fuente centralizada de edición y los \texttt{.txt} funcionan como el formato final consumido por el entrenamiento.

\medskip

El script \texttt{sync\_captions.py} también normaliza la palabra de activación. Si una caption no comienza con \texttt{pixelart}, el script la añade automáticamente al principio. De esta manera se mantenía una estructura homogénea en todos los captions y se reforzaba la asociación entre el estilo aprendido y el término utilizado después en los prompts de generación.

La Tabla~\ref{tab:scripts_dataset_iteracion_2} resume los scripts utilizados para preparar las imágenes y sincronizar los captions en esta iteración.

\begin{table}[htbp]
\centering
\estilotabla
\begin{tabular}{|L{0.30\textwidth}|L{0.58\textwidth}|}
\hline
\textbf{Script} & \textbf{Función en la segunda iteración} \\
\hline
\texttt{scale\_pixelart.py} & Procesa las imágenes originales y genera versiones PNG de 512x512 píxeles preservando el aspecto pixel-art. \\
\hline
\texttt{sync\_captions.py} & Genera los captions \texttt{.txt} individuales a partir de \texttt{captions.csv}. \\
\hline
\end{tabular}
\caption{Scripts utilizados en la preparación del dataset de la segunda iteración.}
\label{tab:scripts_dataset_iteracion_2}
\end{table}

Los captions de esta segunda iteración eran más descriptivos que los iniciales. Además del sujeto principal, incorporaban información sobre la perspectiva, el tipo de imagen, los colores, el fondo y rasgos propios del pixel-art, como \textit{limited palette}, \textit{hard edges}, \textit{cluster shading} y \textit{retro game style}. Con ello se buscaba que el modelo recibiera señales textuales más precisas y no aprendiera el estilo únicamente a partir de descripciones demasiado genéricas.

\subsection{Encadenamiento mediante pretrained\_lora\_path}
\label{subsec:pretrained_lora_iteracion_2}

El cambio más importante de esta iteración fue la incorporación de \texttt{pretrained\_lora\_path}. Este campo indica al entrenamiento que debe cargar un LoRA existente como punto de partida antes de continuar el aprendizaje. En este caso, se utilizó el LoRA final generado en la Iteración 1, almacenado en \texttt{outputs/iter\_01/pixel\_art\_lora\_v1/pixel\_art\_lora\_v1.safetensors}.

La Figura~\ref{fig:pretrained_lora_iteracion_2} muestra la configuración utilizada para cargar el LoRA de la iteración anterior como punto de partida.

\begin{figure}[H]
\centering
\includegraphics[width=0.82\textwidth]{imagenes/iteracion2_pretrained_lora_path.png}
\caption{Configuración del segundo entrenamiento con carga del LoRA previo mediante \texttt{pretrained\_lora\_path}. Fuente: captura de la configuración utilizada en el proyecto.}
\label{fig:pretrained_lora_iteracion_2}
\end{figure}

Gracias a este mecanismo, el segundo entrenamiento empezó desde el adaptador ya entrenado previamente, en lugar de hacerlo desde pesos LoRA aleatorios. Esto convirtió la iteración en una continuación del proceso anterior y permitió comprobar si el entrenamiento podía funcionar de manera acumulativa.

\medskip

Desde el punto de vista de la trazabilidad, esta decisión es relevante porque la conexión entre ambas iteraciones queda registrada en el archivo YAML y en la copia de configuración almacenada en la carpeta de salida. El LoRA \texttt{pixel\_art\_lora\_v2} puede interpretarse, por tanto, como resultado de entrenar sobre el modelo base FLUX.1-dev más el aprendizaje heredado de \texttt{pixel\_art\_lora\_v1}.

\subsection{Entrenamiento local de la segunda iteración}
\label{subsec:entrenamiento_iteracion_2}

El entrenamiento se ejecutó de nuevo en local, reutilizando el entorno validado en la primera iteración. El archivo empleado fue \texttt{config/iter\_02\_train.yaml}, ejecutado mediante \texttt{python run.py config/iter\_02\_train.yaml}. La salida se almacenó en \texttt{outputs/iter\_02/pixel\_art\_lora\_v2/}.

La Tabla~\ref{tab:parametros_iteracion_2} recoge los parámetros principales utilizados en este segundo entrenamiento.

\begin{table}[htbp]
\centering
\estilotabla
\begin{tabular}{|L{0.28\textwidth}|L{0.22\textwidth}|L{0.40\textwidth}|}
\hline
\textbf{Parámetro} & \textbf{Valor} & \textbf{Finalidad} \\
\hline
Modelo base & FLUX.1-dev & Modelo sobre el que se aplica el adaptador LoRA. \\
\hline
LoRA previo & \texttt{pixel\_art\_lora\_v1} & Punto de partida de la mejora incremental. \\
\hline
Dataset & \texttt{datasets/iter\_02/images} & Imágenes y captions de la segunda iteración. \\
\hline
Pasos & 500 & Duración del segundo entrenamiento. \\
\hline
Rango y alpha LoRA & 16 / 16 & Misma capacidad del adaptador que en la iteración anterior. \\
\hline
Learning rate & $1 \times 10^{-4}$ & Tasa de aprendizaje mantenida para comparar con la primera iteración. \\
\hline
Sampling & Cada 100 pasos & Generación de muestras intermedias para observar la evolución. \\
\hline
\end{tabular}
\caption{Parámetros principales de la segunda iteración.}
\label{tab:parametros_iteracion_2}
\end{table}

El entrenamiento generó el archivo final \texttt{pixel\_art\_lora\_v2.safetensors}, varios checkpoints intermedios y una carpeta de samples con imágenes generadas cada 100 pasos. Estos artefactos permitieron comprobar que el encadenamiento desde el LoRA anterior funcionaba y que la nueva organización por iteraciones conservaba los resultados de forma separada.

La Tabla~\ref{tab:artefactos_iteracion_2} resume los artefactos generados durante el segundo entrenamiento.

\begin{table}[htbp]
\centering
\estilotabla
\begin{tabular}{|L{0.34\textwidth}|L{0.52\textwidth}|}
\hline
\textbf{Artefacto} & \textbf{Función} \\
\hline
\texttt{pixel\_art\_lora\_v2.safetensors} & LoRA final de la segunda iteración. \\
\hline
\texttt{pixel\_art\_lora\_v2\_000000NNN.safetensors} & Checkpoints intermedios guardados durante el entrenamiento. \\
\hline
\texttt{optimizer.pt} & Estado del optimizador, útil para reanudar el entrenamiento. \\
\hline
\texttt{config.yaml} & Copia de la configuración usada, necesaria para la trazabilidad. \\
\hline
\texttt{samples/} & Muestras generadas automáticamente durante el entrenamiento. \\
\hline
\end{tabular}
\caption{Artefactos generados por el segundo entrenamiento.}
\label{tab:artefactos_iteracion_2}
\end{table}

La Figura~\ref{fig:samples_iteracion_2} muestra la evolución de una muestra generada durante el segundo entrenamiento.

\begin{figure}[H]
\centering
\includegraphics[width=0.23\textwidth]{imagenes/iteracion2_sample_forest_step_0.jpg}
\hfill
\includegraphics[width=0.23\textwidth]{imagenes/iteracion2_sample_forest_step_100.jpg}
\hfill
\includegraphics[width=0.23\textwidth]{imagenes/iteracion2_sample_forest_step_300.jpg}
\hfill
\includegraphics[width=0.23\textwidth]{imagenes/iteracion2_sample_forest_step_500.jpg}
\caption{Evolución de una muestra generada durante el segundo entrenamiento. Fuente: generación propia obtenida durante el segundo entrenamiento LoRA.}
\label{fig:samples_iteracion_2}
\end{figure}

La segunda iteración permitió validar dos aspectos importantes: por un lado, que el proyecto podía organizarse de forma más clara por iteraciones; por otro, que era posible continuar el entrenamiento desde un LoRA anterior mediante \texttt{pretrained\_lora\_path}. Sin embargo, el dataset seguía siendo reducido y preparado manualmente, por lo que la siguiente fase debía centrarse en ampliar el volumen de datos, automatizar parte del proceso y mejorar la escalabilidad del entrenamiento.

\section{Iteración 3: automatización y escalado del pipeline}
\label{sec:iteracion_3}

\subsection{Objetivos de la tercera iteración}
\label{subsec:objetivos_iteracion_3}

La tercera iteración supuso el paso de un flujo de entrenamiento local y parcialmente manual a un pipeline más completo, automatizado y escalable. Las dos primeras iteraciones habían permitido validar el entrenamiento LoRA, la preparación básica de datasets y el encadenamiento entre modelos mediante \texttt{pretrained\_lora\_path}. Sin embargo, todavía quedaban limitaciones importantes: la obtención de imágenes dependía de selección manual, los captions requerían mucho trabajo de edición, el entrenamiento local limitaba la escala de las pruebas y la revisión de resultados estaba poco integrada.

\medskip

El objetivo de esta fase fue ampliar el sistema para que no se limitara a ejecutar entrenamientos aislados. Para ello se adaptó la interfaz web del proyecto base, se integró un scraper de Pixilart, se añadió generación automática de captions con Qwen2.5-VL, se conectó el entrenamiento con Modal y se incorporó una vista de generación y evaluación de imágenes. Esta iteración convirtió el proyecto en un flujo más trazable: desde la imagen recopilada hasta el dataset final, el entrenamiento LoRA y las imágenes generadas.

\subsection{Adaptación de la interfaz web}
\label{subsec:interfaz_iteracion_3}

La interfaz web existente de \texttt{ai-toolkit} se amplió para actuar como punto de control del pipeline. La aplicación está implementada con Next.js, React, TypeScript y Tailwind, y utiliza Prisma con SQLite para registrar trabajos, colas y ajustes. Sobre esta base se añadieron o adaptaron secciones para scraping, grupos de captions, gestión de datasets, entrenamientos, visualización de \textit{samples}, generación de imágenes y evaluación.

\medskip

La interfaz actúa como coordinadora de los scripts internos. Las acciones del usuario se traducen en llamadas a rutas API de Next.js, que a su vez lanzan scripts de Python, leen o escriben archivos del sistema y actualizan estados persistentes. De esta forma, procesos que antes requerían ejecutar comandos separados pasan a estar conectados dentro de un mismo flujo.

La Tabla~\ref{tab:interfaz_iteracion_3} resume las principales secciones adaptadas en la interfaz durante esta tercera iteración.

\begin{table}[htbp]
\centering
\estilotabla
\begin{tabular}{|L{0.30\textwidth}|L{0.60\textwidth}|}
\hline
\textbf{Módulo de la interfaz} & \textbf{Función implementada} \\
\hline
\texttt{/scraper} & Lanzamiento del scraper por topics y revisión de imágenes aceptadas o rechazadas. \\
\hline
\texttt{/captions} & Creación de grupos, generación automática, revisión y exportación al dataset. \\
\hline
\texttt{/datasets} & Consulta y organización de datasets de entrenamiento. \\
\hline
\texttt{/jobs} & Creación, cola y seguimiento de entrenamientos locales o en Modal. \\
\hline
\texttt{/trainings} & Visualización de \textit{samples} agrupados por iteración y entrenamiento. \\
\hline
\texttt{/generate} & Generación de imágenes con FLUX base o con LoRAs entrenados. \\
\hline
\end{tabular}
\caption{Secciones principales adaptadas en la interfaz durante la tercera iteración.}
\label{tab:interfaz_iteracion_3}
\end{table}

\subsection{Integración con scripts del pipeline}

La interfaz web actúa como una capa de integración sobre los scripts principales del proyecto. Las rutas API y la lógica de servidor permiten lanzar procesos de scraping, captioning, generación, evaluación y entrenamiento sin ejecutar manualmente cada herramienta por separado.

La Tabla~\ref{tab:scripts_iteracion3} recoge los scripts principales integrados en la interfaz y su función dentro del pipeline.

\begin{table}[htbp]
\centering
\estilotabla
\begin{tabular}{|L{0.28\textwidth}|L{0.27\textwidth}|L{0.35\textwidth}|}
\hline
\textbf{Script} & \textbf{Lanzado desde} & \textbf{Función dentro del sistema} \\
\hline
\texttt{run\_spider.py} & \texttt{ui/src/server/scraper.ts} & Ejecuta el scraper de Pixilart a partir de topics, límites de imágenes y filtros de calidad. \\
\hline
\texttt{generate\_captions.py} & \texttt{ui/src/server/captions.ts} & Genera captions automáticos con Qwen2.5-VL para las imágenes revisadas. \\
\hline
\texttt{generate\_api.py} & Rutas de generación & Permite generar imágenes usando FLUX.1-dev y, opcionalmente, un LoRA entrenado. \\
\hline
\texttt{evaluate\_style.py} & Servicios de evaluación & Calcula métricas técnicas para analizar características visuales de las imágenes generadas. \\
\hline
\texttt{run\_modal.py} & Sistema de entrenamiento & Adapta y lanza entrenamientos en Modal, preparando rutas, configuración y transferencia de artefactos. \\
\hline
\end{tabular}
\caption{Scripts principales integrados en la interfaz.}
\label{tab:scripts_iteracion3}
\end{table}

\subsection{Scraping, filtrado y revisión de imágenes}

El scraper se diseñó para automatizar la obtención de imágenes de Pixilart evitando depender de una descarga manual. En esta iteración se trabajó con topics en lugar de tags, ya que los topics permitían buscar imágenes dentro de categorías más amplias y más cercanas al tipo de contenido que se quería incorporar al dataset.


El proceso comienza desde la interfaz, donde se define la ejecución del scraper. A partir de esa petición, el backend lanza \texttt{run\_spider.py}, que invoca el spider de Scrapy encargado de recorrer las páginas de Pixilart, extraer las imágenes candidatas y asociarles metadatos como título, autor, topic, tags, número de likes, fecha de publicación y tamaño del canvas.

Después de la extracción se aplican varios filtros automáticos antes de aceptar una imagen. Esta fase evita que el dataset se llene de imágenes corruptas, animaciones, obras demasiado grandes para el objetivo de pixel-art de baja resolución o contenidos marcados con tags problemáticas.


Las primeras pruebas mostraron que obtener imágenes únicamente por temática no era suficiente. Aunque muchas imágenes pertenecían al estilo pixel-art, algunas tenían poca calidad visual, estaban poco terminadas o representaban figuras difíciles de reconocer. Para mejorar la consistencia del conjunto obtenido se añadió un filtro basado en el número mínimo de likes. La idea fue utilizar la valoración de la comunidad como una señal adicional de calidad: si una imagen tenía una recepción mínima, era más probable que estuviera terminada, fuese reconocible y tuviera un mayor valor para el entrenamiento.

La Tabla~\ref{tab:filtros_scraper_iteracion3} resume los criterios de filtrado aplicados durante el proceso de scraping.

\begin{table}[htbp]
\centering
\estilotabla
\begin{tabular}{|L{0.30\textwidth}|L{0.55\textwidth}|}
\hline
\textbf{Criterio} & \textbf{Tratamiento aplicado} \\
\hline
GIF animado & Rechazo automático, ya que el entrenamiento trabaja con imágenes estáticas. \\
\hline
Imagen corrupta & Rechazo automático si Pillow no puede abrir o validar el archivo descargado. \\
\hline
Canvas demasiado grande & Descarte de imágenes alejadas del pixel-art de baja resolución buscado. \\
\hline
Tags problemáticas & Rechazo automático de imágenes marcadas como \texttt{notmine}, \texttt{stolen}, \texttt{repost} o \texttt{traced}. \\
\hline
Número mínimo de likes & Filtro añadido tras las primeras pruebas para reducir imágenes poco reconocibles o de baja calidad visual. \\
\hline
\end{tabular}
\caption{Criterios de filtrado aplicados durante el scraping.}
\label{tab:filtros_scraper_iteracion3}
\end{table}

Cada ejecución del scraper genera una carpeta propia dentro de \texttt{datasets/pixilart/}, con imágenes aceptadas, rechazadas, pendientes de revisión y ficheros de metadatos. Esta organización permite conservar la trazabilidad de cada ejecución sin mezclar resultados de distintas búsquedas.

La Figura~\ref{fig:iteracion3_scraper_rejected} muestra la vista de revisión de imágenes rechazadas por el scraper.

\begin{figure}[H]
\centering
\includegraphics[width=0.90\textwidth]{imagenes/iteracion3_scraper_rejected.png}
\caption{Interfaz de revisión de imágenes rechazadas por el scraper. Fuente: captura de la interfaz desarrollada en el proyecto.}
\label{fig:iteracion3_scraper_rejected}
\end{figure}


\subsection{Captioning automático con Qwen2.5-VL}

Una vez filtradas y revisadas las imágenes, el siguiente paso consistió en generar captions más consistentes. Para ello se integró el modelo Qwen2.5-VL, un modelo visión-lenguaje capaz de recibir una imagen junto con una instrucción textual y devolver una descripción asociada al contenido visual.

En el sistema, esta funcionalidad se implementa mediante el script \texttt{generate\_captions.py}. Desde la interfaz se lanza a través de la ruta \texttt{/api/captions/groups/\{groupId\}/generate}, que permite seleccionar un grupo de imágenes y generar captions automáticos para ellas. El modelo recibe la imagen, la palabra de activación \texttt{pixelart} y el nivel de detalle solicitado.

La función de este bloque es transformar un conjunto de imágenes revisadas en un conjunto entrenable, donde cada imagen queda asociada a una descripción textual coherente. Esto es importante porque el entrenamiento LoRA aprende de las imágenes y de la relación entre cada imagen y su caption.

El script permite generar captions con tres niveles de detalle. Estos niveles cambian la longitud del texto y también el tipo de información que se fuerza a incluir.

La Tabla~\ref{tab:niveles_captioning_iteracion3} resume los niveles de detalle disponibles para generar captions automáticos.

\begin{table}[htbp]
\centering
\estilotabla
\begin{tabular}{|L{0.18\textwidth}|L{0.67\textwidth}|}
\hline
\textbf{Nivel} & \textbf{Función dentro del dataset} \\
\hline
\texttt{Corto} & Genera captions muy breves, centrados en el sujeto principal y algunos rasgos generales del estilo. \\
\hline
\texttt{Medio} & Genera captions intermedios con categoría, perspectiva, detalles del sujeto y rasgos comunes del pixel-art. \\
\hline
\texttt{Detallado} & Genera descripciones más largas, con información adicional sobre composición, escena, colores y elementos visuales. \\
\hline
\end{tabular}
\caption{Niveles de detalle disponibles en la generación automática de captions.}
\label{tab:niveles_captioning_iteracion3}
\end{table}

Durante la generación se utiliza una configuración determinista, sin muestreo aleatorio, para que el resultado sea más estable y fácil de revisar. Además, cuando existen metadatos procedentes del scraper, como título o tags, estos se utilizan como apoyo contextual, aunque el caption final se genera a partir de la interpretación visual de la imagen.

Para ilustrar la diferencia entre niveles, la Figura~\ref{fig:iteracion3_caption_ejemplo} muestra una imagen de ejemplo y la Tabla~\ref{tab:ejemplo_captions_iteracion3} compara los captions generados para esa misma imagen.

\begin{figure}[H]
\centering
\includegraphics[width=0.30\textwidth]{imagenes/iteracion3_caption_ejemplo_imagen.png}
\caption{Imagen utilizada como ejemplo para comparar niveles de captioning. Fuente: imagen recopilada para el dataset mediante el pipeline del proyecto.}
\label{fig:iteracion3_caption_ejemplo}
\end{figure}

\begingroup
\estilotabla
\begin{longtable}{|L{0.18\textwidth}|L{0.70\textwidth}|}
\caption{Comparación de captions generados con distintos niveles de detalle.}
\label{tab:ejemplo_captions_iteracion3}\\
\hline
\textbf{Nivel} & \textbf{Ejemplo de caption} \\
\hline
\endfirsthead

\hline
\textbf{Nivel} & \textbf{Ejemplo de caption} \\
\hline
\endhead
Corto & \texttt{pixelart, pixel art, cute dog, limited palette} \\
\hline
Medio & \texttt{pixelart, creature, front view, cute dog with heart, light purple and white colors, limited palette, hard edges, cluster shading, retro game style} \\
\hline
Detallado & \texttt{pixelart, This pixel art image features a cute dog as the main subject, rendered in a limited palette with a retro, 8-bit aesthetic. The dog has a white face with light blue eyes, dark purple ears, and a small red heart on its left side. The background is a simple green ground with small green plants and a single red flower. The overall mood is cheerful and playful. The image uses outlined sprites for the dog's features, with pixel shading adding depth. The limited color palette and outlined edges give it a classic pixel art look, suitable for a 64x64 canvas.} \\
\hline
\end{longtable}
\endgroup

Tras la generación automática, los captions podían revisarse desde la interfaz antes de mover las imágenes al dataset final. Esta revisión permitía corregir descripciones demasiado largas, poco precisas o mal adaptadas al estilo buscado.

\subsection{Preparación final del dataset}

Una vez revisados los captions, las imágenes se movían al dataset final desde la propia interfaz. Este paso se implementó dentro de la función \texttt{moveToDataset} de \texttt{ui/src/server/captions.ts}, sin utilizar un script externo.

La función \texttt{moveToDataset} actúa como puente entre la fase de captioning y la fase de entrenamiento. Su objetivo es convertir un grupo de imágenes revisadas en una estructura compatible con \texttt{ai-toolkit}. Para ello copia las imágenes seleccionadas, las redimensiona a \texttt{512x512} con \texttt{sharp}, usando \texttt{kernel.nearest}, y las guarda en formato PNG.

El uso de \texttt{kernel.nearest} es importante en este proyecto porque mantiene los bordes duros y evita el suavizado propio de otros métodos de interpolación. De esta forma, aunque la imagen se adapte al tamaño requerido por el entrenamiento, se conserva mejor la estética pixel-art.

Además de preparar las imágenes, la función actualiza \texttt{captions.csv} y genera un archivo \texttt{.txt} por cada imagen. Estos archivos de texto son necesarios porque la configuración de entrenamiento utiliza \texttt{caption\_ext: txt}, por lo que \texttt{ai-toolkit} espera encontrar un caption individual con el mismo nombre base que la imagen.

Después de varias pruebas de scraping, filtrado y revisión, el dataset utilizado en esta iteración quedó formado por 130 imágenes. 

\subsection{Entrenamiento en Modal}

La tercera iteración incorporó Modal para ejecutar entrenamientos en GPU cloud. Modal permite lanzar código Python en contenedores remotos con GPU, evitando depender únicamente de los recursos disponibles en el equipo local. En este proyecto se utilizó para reducir las limitaciones de VRAM, mejorar la estabilidad de ejecuciones largas y poder repetir entrenamientos sin bloquear el ordenador principal.

Desde la interfaz se añadió la posibilidad de crear un entrenamiento seleccionando el dataset, el LoRA previo, los parámetros principales y la GPU de ejecución. En lugar de tratar Modal como una herramienta externa aislada, se integró dentro del flujo del sistema: el usuario configura el entrenamiento desde la interfaz y el backend prepara la ejecución remota.

La Figura~\ref{fig:iteracion3_modal_gpu_selector} muestra el selector utilizado para elegir entre ejecución local o ejecución en Modal.

\begin{figure}[H]
\centering
\includegraphics[width=0.35\textwidth]{imagenes/iteracion3_modal_gpu_selector.png}
\caption{Selector de GPU para entrenamientos locales o en Modal. Fuente: captura de la interfaz desarrollada en el proyecto.}
\label{fig:iteracion3_modal_gpu_selector}
\end{figure}

La adaptación a Modal se concentró principalmente en \texttt{run\_modal.py}. Este archivo define la imagen del contenedor, instala las dependencias necesarias, monta un volumen persistente para conservar resultados y transforma la configuración local en una configuración válida dentro del entorno remoto.

Uno de los cambios más importantes fue la conversión de rutas. Las rutas locales de Windows, como las utilizadas para datasets u outputs, no existen dentro del contenedor Linux de Modal. Por ello, antes de iniciar el entrenamiento se reescriben las rutas para que apunten a las carpetas disponibles en el entorno remoto. Además, los resultados generados se guardan en un volumen de Modal y posteriormente se descargan al proyecto local.

\subsection{Transferencia del LoRA previo}

El sistema también permite seleccionar un LoRA anterior como punto de partida del nuevo entrenamiento. Esta funcionalidad es importante porque mantiene el enfoque incremental del proyecto: cada iteración puede partir de un modelo ya adaptado y refinarlo con nuevos datos o una configuración corregida.

La Figura~\ref{fig:iteracion3_selector_lora_previo} muestra la selección de un LoRA previo desde la configuración del entrenamiento.

\begin{figure}[H]
\centering
\includegraphics[width=0.35\textwidth]{imagenes/iteracion3_selector_lora_previo.png}
\caption{Selección de un LoRA previo desde la interfaz de configuración del entrenamiento. Fuente: captura de la interfaz desarrollada en el proyecto.}
\label{fig:iteracion3_selector_lora_previo}
\end{figure}

Para que esto funcione en Modal, \texttt{run\_modal.py} comprueba si la configuración contiene \texttt{pretrained\_lora\_path}. Si el archivo existe en local, se lee el \texttt{.safetensors}, se codifica en base64 y se añade a la configuración enviada al contenedor. Una vez dentro de Modal, el archivo se reconstruye antes de iniciar el entrenamiento.

Este mecanismo permite reutilizar LoRAs generados en iteraciones anteriores sin copiar manualmente archivos entre máquinas ni depender de que las rutas locales coincidan con las rutas del contenedor remoto.
\subsection{Seguimiento de la ejecución}

Durante la ejecución de un entrenamiento en Modal, el sistema permite realizar seguimiento desde dos niveles complementarios: la propia interfaz de \texttt{ai-toolkit} y el panel web de Modal. Esta doble visualización facilita comprobar que el entrenamiento se ha lanzado correctamente, identificar en qué fase se encuentra y detectar posibles errores durante la ejecución.

Desde la interfaz de \texttt{ai-toolkit} se muestra el estado general del job, el nombre del entrenamiento, la GPU asignada, el progreso en pasos y la salida de consola generada por el proceso. Esta vista permite seguir la evolución del entrenamiento sin abandonar la aplicación principal del proyecto. También resulta útil para comprobar si el modelo se está cargando correctamente, si el entrenamiento ha comenzado, si se han generado checkpoints o si aparece algún error durante la ejecución.

La Figura~\ref{fig:iteracion3_aitoolkit_seguimiento} muestra el seguimiento de un entrenamiento desde la interfaz de \texttt{ai-toolkit}.

\begin{figure}[H]
\centering
\includegraphics[width=0.95\textwidth]{imagenes/iteracion3_aitoolkit_seguimiento.png}
\caption{Seguimiento de un entrenamiento desde la interfaz de ai-toolkit. Fuente: captura de la interfaz de \texttt{ai-toolkit} utilizada en el proyecto.}
\label{fig:iteracion3_aitoolkit_seguimiento}
\end{figure}

Por otra parte, Modal proporciona una vista propia de la ejecución remota. Desde su panel se pueden consultar los logs de la aplicación, comprobar que la función de entrenamiento está activa, revisar la salida producida dentro del contenedor y verificar que se está utilizando el entorno remoto correctamente. Esta vista es especialmente útil para diagnosticar problemas relacionados con la ejecución cloud, como errores de carga, fallos del contenedor, problemas de dependencias o interrupciones en el proceso.

De esta forma, la interfaz de \texttt{ai-toolkit} ofrece una visión integrada del entrenamiento dentro del flujo del proyecto, mientras que Modal permite revisar con más detalle la ejecución remota. La combinación de ambas vistas fue importante durante esta iteración, ya que permitió comprobar el avance de los entrenamientos y analizar con mayor precisión los errores o comportamientos inesperados.

La Figura~\ref{fig:iteracion3_modal_ejecucion} muestra la vista de seguimiento del entrenamiento desde el panel web de Modal.

\begin{figure}[H]
\centering
\includegraphics[width=0.95\textwidth]{imagenes/iteracion3_modal_ejecucion.png}
\caption{Seguimiento de un entrenamiento desde el panel de Modal. Fuente: captura del panel web de Modal durante la ejecución del entrenamiento.}
\label{fig:iteracion3_modal_ejecucion}
\end{figure}



\subsection{Entrenamientos realizados en Modal}

En esta iteración se realizaron varios entrenamientos. El aspecto más relevante fue ejecutar el pipeline en cloud y detectar por qué una primera configuración producía resultados alejados del objetivo, para corregirla mediante una nueva ejecución.

\subsubsection{Primer entrenamiento en Modal}

El primer entrenamiento de esta iteración generó el LoRA \texttt{pixel\_art\_lora\_v3}, almacenado en \texttt{outputs/iter\_03/pixel\_art\_lora\_v3/}. El entrenamiento se completó y produjo checkpoints, samples y un archivo final \texttt{.safetensors}. No obstante, al comparar las muestras generadas durante la ejecución se observó que la evolución visual no avanzaba hacia el estilo buscado. Aunque el sample inicial presentaba una estética cercana al pixel-art, el resultado final mostraba una imagen más oscura, menos limpia y con menor claridad visual.

La Tabla~\ref{tab:config_entrenamiento3_modal} resume la configuración principal utilizada en el primer entrenamiento ejecutado en Modal.

\begingroup
\estilotabla
\begin{longtable}{|L{0.30\textwidth}|L{0.58\textwidth}|}
\caption{Configuración principal del primer entrenamiento ejecutado en Modal.}
\label{tab:config_entrenamiento3_modal}\\
\hline
\textbf{Parámetro} & \textbf{Valor en el primer entrenamiento Modal} \\
\hline
\endfirsthead

\hline
\textbf{Parámetro} & \textbf{Valor en el primer entrenamiento Modal} \\
\hline
\endhead
Nombre del LoRA & \texttt{pixel\_art\_lora\_v3} \\
\hline
Dataset & \texttt{datasets/iter\_03/images/} \\
\hline
Tipo de captions & Captions detallados \\
\hline
LoRA previo & \texttt{pixel\_art\_lora\_v2.safetensors} \\
\hline
Pasos de entrenamiento & 1500 \\
\hline
Batch size & 1 \\
\hline
Learning rate & \texttt{0.0001} \\
\hline
Rango LoRA & \texttt{linear: 16}, \texttt{conv: 16} \\
\hline
\texttt{bypass\_guidance\_embedding} & \texttt{true} \\
\hline
\texttt{quantize\_te} & \texttt{true} \\
\hline
\texttt{qtype\_te} & \texttt{qfloat8} \\
\hline
\texttt{timestep\_type} & \texttt{sigmoid} \\
\hline
Sampling & Cada 150 pasos, con prompts fijos \\
\hline
\end{longtable}
\endgroup

La Figura~\ref{fig:iteracion3_entrenamiento3_comparacion} compara el sample inicial y el sample final de este primer entrenamiento en Modal.

\begin{figure}[H]
\centering
\begin{minipage}{0.38\textwidth}
    \centering
    \includegraphics[width=\textwidth]{imagenes/iteracion3_entrenamiento3_step0.jpg}

    \small Sample inicial, paso 0.
\end{minipage}
\hspace{0.06\textwidth}
\begin{minipage}{0.38\textwidth}
    \centering
    \includegraphics[width=\textwidth]{imagenes/iteracion3_entrenamiento3_step1500.jpg}

    \small Sample final, paso 1500.
\end{minipage}
\caption{Comparación entre el sample inicial y el sample final del primer entrenamiento en Modal. Fuente: generaciones propias obtenidas durante el entrenamiento en Modal.}
\label{fig:iteracion3_entrenamiento3_comparacion}
\end{figure}

\subsubsection{Diagnóstico de la configuración}

El análisis de \texttt{config.yaml} permitió identificar varios parámetros que podían explicar la diferencia entre el entrenamiento y los resultados esperados. El problema principal fue que la configuración utilizada para entrenar no estaba alineada con la forma en la que después se generaban imágenes con FLUX.1-dev.

La Tabla~\ref{tab:diagnostico_parametros_modal} recoge los parámetros revisados y los cambios aplicados tras el diagnóstico del primer entrenamiento en Modal.

\begingroup
\estilotabla
\begin{longtable}{|L{0.23\textwidth}|L{0.22\textwidth}|L{0.43\textwidth}|}
\caption{Diagnóstico de parámetros y cambios aplicados tras el primer entrenamiento Modal.}
\label{tab:diagnostico_parametros_modal}\\
\hline
\textbf{Parámetro} & \textbf{Cambio aplicado} & \textbf{Consecuencia del cambio} \\
\hline
\endfirsthead

\hline
\textbf{Parámetro} & \textbf{Cambio aplicado} & \textbf{Consecuencia del cambio} \\
\hline
\endhead
\texttt{bypass\_guidance\_embedding} & De \texttt{true} a \texttt{false} & El entrenamiento vuelve a usar el condicionamiento de guidance de forma coherente con la inferencia. Esto evita que el LoRA aprenda bajo unas condiciones distintas a las usadas después al generar imágenes. \\
\hline
\texttt{quantize\_te} & De \texttt{true} a \texttt{false} & Los text encoders dejan de cuantizarse. Esto aumenta el consumo de memoria, pero mantiene representaciones textuales más precisas, algo importante cuando los captions contienen información descriptiva. \\
\hline
\texttt{qtype\_te} & Se elimina al desactivar \texttt{quantize\_te} & Al no cuantizar el text encoder en \texttt{qfloat8}, se reduce la pérdida de información en la interpretación del prompt y de los captions. \\
\hline
\texttt{timestep\_type} & Se elimina \texttt{sigmoid} & El entrenamiento vuelve a una configuración más próxima a la línea base que había funcionado en la segunda iteración, evitando introducir un cambio adicional en la distribución de timesteps. \\
\hline
Número de pasos & De 1500 a 500 en el entrenamiento corregido & Se reduce la exposición a una configuración experimental y se facilita comprobar si la corrección de parámetros mejora el comportamiento del LoRA. \\
\hline
Captions & De captions detallados a una revisión posterior con captions medios & Los captions medios reducen ruido textual y concentran la descripción en elementos más útiles para el aprendizaje del estilo. \\
\hline
\end{longtable}
\endgroup

En concreto, \texttt{bypass\_guidance\_embedding: true} hacía que el entrenamiento ignorase un condicionamiento que FLUX.1-dev sí utiliza durante la generación. Como consecuencia, el LoRA se optimizaba en un escenario distinto al de uso real. Al cambiarlo a \texttt{false}, el entrenamiento y la inferencia pasan a trabajar de forma más coherente.

También se modificó \texttt{quantize\_te}. En el primer entrenamiento, los text encoders estaban cuantizados con \texttt{qfloat8}. Esto reducía memoria, pero podía degradar la representación de captions largos. Al desactivar esta cuantización, se priorizó la estabilidad semántica frente al ahorro de VRAM.

Por último, se eliminó \texttt{timestep\_type: sigmoid}. Este parámetro alteraba la forma en la que se seleccionaban los timesteps durante el entrenamiento. Como la segunda iteración había funcionado con una configuración más sencilla, se decidió volver a una base más comparable antes de introducir nuevos cambios.

\subsubsection{Entrenamiento corregido}

El cuarto entrenamiento fue el más importante de esta iteración, ya que aplicó las correcciones detectadas tras el diagnóstico anterior. Este entrenamiento generó \texttt{pixel\_art\_lora\_v3\_Arreglo}, conservado en \texttt{outputs/iter\_03\_Arreglo/pixel\_art\_lora\_v3\_Arreglo/}.

A diferencia del primer entrenamiento Modal, no se partió del LoRA problemático. Se volvió a utilizar como punto de partida el LoRA válido de la segunda iteración, \texttt{pixel\_art\_lora\_v2}, y se aplicó una configuración más cercana a la que ya había dado resultados estables.

La Tabla~\ref{tab:comparacion_entrenamiento3_4} resume las principales diferencias entre el primer entrenamiento en Modal y el entrenamiento corregido.

\begingroup
\estilotabla
\begin{longtable}{|L{0.26\textwidth}|L{0.28\textwidth}|L{0.34\textwidth}|}
\caption{Comparación entre el primer entrenamiento Modal y el entrenamiento corregido.}
\label{tab:comparacion_entrenamiento3_4}\\
\hline
\textbf{Aspecto} & \textbf{Entrenamiento 3} & \textbf{Entrenamiento 4 corregido} \\
\hline
\endfirsthead

\hline
\textbf{Aspecto} & \textbf{Entrenamiento 3} & \textbf{Entrenamiento 4 corregido} \\
\hline
\endhead
LoRA generado & \texttt{pixel\_art\_lora\_v3} & \texttt{pixel\_art\_lora\_v3\_Arreglo} \\
\hline
LoRA previo & Iteración 2 & Iteración 2 \\
\hline
Captions & Detallados & Revisión con configuración más estable \\
\hline
Pasos & 1500 & 500 \\
\hline
\texttt{bypass\_guidance\_embedding} & \texttt{true} & \texttt{false} \\
\hline
\texttt{quantize\_te} & \texttt{true} & \texttt{false} \\
\hline
\texttt{timestep\_type} & \texttt{sigmoid} & No aplicado \\
\hline
Resultado visual & Alejado del pixel-art buscado & Mucho más cercano al estilo objetivo \\
\hline
\end{longtable}
\endgroup

La Figura~\ref{fig:iteracion3_muestra_corregida} muestra una muestra generada en el paso 500 del entrenamiento corregido.

\begin{figure}[H]
\centering
\includegraphics[width=0.28\textwidth]{imagenes/iteracion3_entrenamiento4_muestra_corregida.jpg}
\caption{Muestra generada en el paso 500 del entrenamiento corregido, con un estilo pixel-art más definido. Fuente: generación propia obtenida durante el entrenamiento corregido.}
\label{fig:iteracion3_muestra_corregida}
\end{figure}

Este entrenamiento confirmó que el problema se encontraba en la configuración concreta utilizada en el primer intento, y no en Modal ni en la transferencia de archivos. La conservación de \texttt{config.yaml}, samples y artefactos permitió comparar ambas ejecuciones y justificar los cambios aplicados.

\subsubsection{Entrenamiento posterior con captions medios}

Después del entrenamiento corregido se realizó una ejecución adicional para comprobar si partir del LoRA arreglado y utilizar captions de detalle medio podía mejorar la consistencia del resultado. Este entrenamiento generó \texttt{pixel\_art\_lora\_v4}, utilizando como punto de partida \texttt{pixel\_art\_lora\_v3\_Arreglo}.

La Tabla~\ref{tab:entrenamiento5_captions_medios} resume la configuración de este entrenamiento posterior con captions medios.

\begin{table}[htbp]
\centering
\estilotabla
\begin{tabular}{|L{0.30\textwidth}|L{0.50\textwidth}|}
\hline
\textbf{Aspecto} & \textbf{Valor} \\
\hline
LoRA generado & \texttt{pixel\_art\_lora\_v4} \\
\hline
LoRA previo & \texttt{pixel\_art\_lora\_v3\_Arreglo} \\
\hline
Dataset & \texttt{datasets/iter\_03\_captions\_medios/images/} \\
\hline
Número de imágenes & 130 \\
\hline
Tipo de captions & Medium \\
\hline
Pasos de entrenamiento & 600 \\
\hline
Objetivo & Comprobar si captions más compactos mejoraban la estabilidad del estilo aprendido. \\
\hline
\end{tabular}
\caption{Configuración resumida del entrenamiento posterior con captions medios.}
\label{tab:entrenamiento5_captions_medios}
\end{table}

\medskip

En conjunto, la tercera iteración permitió transformar el pipeline inicial en un sistema más completo y escalable. La incorporación de la interfaz, el scraper, el captioning automático, la ejecución en Modal y la evaluación de imágenes permitió reducir la dependencia de procesos manuales y mantener una trazabilidad más clara entre datos, entrenamientos y resultados generados.

\medskip

A diferencia de las iteraciones anteriores, esta fase se centró en consolidar el flujo completo de trabajo, además de obtener un nuevo LoRA. El resultado fue una base funcional desde la que generar datasets más amplios, lanzar entrenamientos en la nube, comparar modelos y revisar imágenes generadas de forma integrada.

\subsection{Generación de imágenes desde la interfaz}

Además del entrenamiento, en esta iteración se incorporó una vista para generar imágenes desde la propia interfaz del sistema. El objetivo era poder probar los LoRAs obtenidos sin recurrir a scripts manuales, seleccionando directamente el modelo que se quería utilizar y escribiendo el prompt de generación.

La generación se realiza a través de la ruta \texttt{/api/generate}, que valida la petición recibida desde la interfaz y lanza el script \texttt{scripts/tfg/generate\_api.py}. Este script permite generar imágenes usando el modelo base \texttt{FLUX.1-dev}, los LoRAs catalogados del proyecto o nuevos LoRAs detectados en carpetas de salida dentro de \texttt{outputs/}. La API valida que el modelo seleccionado exista y, si no pertenece al catálogo estático, comprueba que la carpeta correspondiente contenga un archivo \texttt{.safetensors} válido.

La Tabla~\ref{tab:generacion_imagenes_iteracion3} resume los elementos principales que intervienen en la generación de imágenes desde la interfaz.

\begin{table}[htbp]
\centering
\estilotabla
\begin{tabular}{|L{0.28\textwidth}|L{0.56\textwidth}|}
\hline
\textbf{Elemento} & \textbf{Función en la generación} \\
\hline
Prompt & Texto introducido por el usuario para describir la imagen que se quiere generar. \\
\hline
Modelo seleccionado & Permite elegir entre el modelo base FLUX, los LoRAs catalogados de las iteraciones y nuevos LoRAs detectados en \texttt{outputs/}. \\
\hline
Seed & Valor opcional que permite repetir una generación concreta y comparar modelos bajo condiciones similares. \\
\hline
\texttt{generate\_api.py} & Script encargado de cargar FLUX, aplicar el LoRA seleccionado y guardar la imagen generada. \\
\hline
Carpeta de salida & Las imágenes se almacenan en \texttt{outputs/<modelo>/generated/}, manteniendo separadas las generaciones de cada modelo. \\
\hline
\end{tabular}
\caption{Elementos principales del sistema de generación de imágenes.}
\label{tab:generacion_imagenes_iteracion3}
\end{table}

El script normaliza el prompt para asegurar que incluya la palabra de activación \texttt{pixelart}. Si el prompt no empieza por \texttt{pixelart}, se añade automáticamente al comienzo. Esto ayuda a mantener la coherencia con la forma en que fueron entrenados los LoRAs, ya que todos los datasets utilizaban esa palabra como trigger word.

La generación utiliza una resolución de \texttt{512x512}, 20 pasos de inferencia y un \textit{guidance scale} de 4. Cuando el usuario no indica una seed, el sistema genera una automáticamente. Al finalizar, el script devuelve una respuesta en formato JSON con la ruta de la imagen guardada y la seed utilizada, de forma que la interfaz puede mostrar el resultado y conservar la referencia para revisiones posteriores.

La Figura~\ref{fig:iteracion3_generacion_interfaz} muestra la vista de generación integrada en la interfaz.

\begin{figure}[H]
\centering
\includegraphics[width=0.90\textwidth]{imagenes/generador.png}
\caption{Interfaz de generación de imágenes con selección de modelo. Fuente: captura de la interfaz desarrollada en el proyecto.}
\label{fig:iteracion3_generacion_interfaz}
\end{figure}

Esta integración permitió comparar de forma más rápida el comportamiento del modelo base y de los distintos LoRAs entrenados. En lugar de evaluar únicamente los samples generados durante el entrenamiento, se podían crear nuevas imágenes con prompts específicos y comprobar si el estilo aprendido se aplicaba correctamente.
\subsection{Validación técnica de imágenes generadas}

Para complementar la revisión visual, se añadió un sistema de evaluación técnica de imágenes generadas. Esta funcionalidad no pretende sustituir la valoración manual, ya que aspectos como la estética, la composición o la fidelidad semántica siguen requiriendo interpretación humana. Su objetivo es aportar métricas objetivas relacionadas con rasgos característicos del pixel-art.

La validación se ejecuta desde la interfaz mediante la ruta \texttt{/api/generated-images/evaluate}, que llama al script \texttt{scripts/tfg/evaluate\_style.py}. Este script analiza la imagen seleccionada y devuelve métricas en bruto. Posteriormente, la API transforma esas métricas en apartados más fáciles de interpretar desde la interfaz.

El sistema separa la evaluación en tres bloques: construcción pixel-art, lectura visual experimental y paleta compacta. Esta separación es importante porque una imagen puede tener una figura reconocible, pero no presentar una estructura clara de pixel-art; o al contrario, puede tener bordes duros y paleta limitada, pero resultar difícil de interpretar visualmente.

La Tabla~\ref{tab:metricas_validacion_imagenes_iteracion3} recoge las métricas utilizadas para la validación técnica de imágenes generadas.

\begin{table}[htbp]
\centering
\estilotabla
\begin{tabular}{|L{0.22\textwidth}|L{0.24\textwidth}|L{0.44\textwidth}|}
\hline
\textbf{Bloque} & \textbf{Métrica} & \textbf{Qué analiza} \\
\hline
Construcción pixel-art & Tamaño y cuadrícula & Busca si existe un bloque base repetido y si los cambios de color coinciden con una cuadrícula regular. Permite detectar si el píxel aparente mantiene un tamaño estable. \\
\hline
Construcción pixel-art & Curvas escalonadas & Comprueba si curvas y diagonales están formadas por escalones regulares, combinando la alineación con la cuadrícula, la dureza de bordes y la estabilidad de zonas planas. \\
\hline
Construcción pixel-art & Bordes sin suavizado & Mide si los cambios de color son bruscos. Un valor alto indica pocos degradados o antialias, aunque no determina por sí solo si la cuadrícula es correcta. \\
\hline
Lectura visual experimental & Coherencia de la figura & Separa aproximadamente fondo y primer plano, y mide si la mayor parte del sujeto pertenece a una única forma conectada. Ayuda a estimar si la silueta principal se entiende. \\
\hline
Lectura visual experimental & Ruido que dificulta la lectura & Mide cuántos píxeles cambian respecto a sus vecinos. Muchos cambios locales generan textura o ruido visual que puede competir con la figura principal. \\
\hline
Dato adicional & Paleta compacta & Indica cuánto detalle conserva la imagen al reducirse a 16 colores. Es una métrica informativa sobre la simplicidad cromática, pero no modifica las dos notas principales. \\
\hline
\end{tabular}
\caption{Métricas utilizadas para la validación técnica de imágenes generadas.}
\label{tab:metricas_validacion_imagenes_iteracion3}
\end{table}

La primera nota, construcción pixel-art, agrupa las métricas relacionadas con la estructura visual propia del estilo. Para obtener una puntuación alta no basta con que la imagen tenga aspecto retro; también debe presentar bloques de píxel consistentes, bordes duros y diagonales o curvas construidas mediante escalones. Por este motivo, una imagen generada puede ser visualmente atractiva y, aun así, recibir una puntuación baja si presenta suavizado, variaciones de tamaño en el píxel o contornos poco alineados.

La segunda nota, lectura visual experimental, mide si existe una figura principal clara. Esta evaluación analiza si el primer plano forma una silueta continua y si hay demasiado ruido visual alrededor. Se considera experimental porque no identifica qué objeto representa la imagen ni evalúa si el prompt se ha cumplido correctamente; únicamente estima si la figura se puede leer con facilidad.

La paleta compacta se muestra como dato adicional. Esta métrica evalúa la fidelidad de la imagen al reducirse a 16 colores. Resulta útil para detectar si una generación mantiene una paleta relativamente simple o si contiene demasiadas variaciones cromáticas, algo habitual cuando el resultado se aleja del pixel-art más limpio.

\begin{figure}[H]
\centering
\includegraphics[width=0.90\textwidth]{imagenes/iteracion3_validacion_imagen.png}
\caption{Validación técnica de una imagen generada desde la interfaz. Fuente: captura de la interfaz desarrollada en el proyecto.}
\label{fig:iteracion3_validacion_imagen}
\end{figure}

En la Figura~\ref{fig:iteracion3_validacion_imagen} se observa un caso donde la imagen generada obtiene una valoración alta tanto en construcción pixel-art como en lectura visual experimental. La figura principal resulta reconocible y mantiene una silueta clara, mientras que la evaluación técnica detecta una estructura de píxel consistente, curvas escalonadas y bordes con poco suavizado. Esto permite comprobar cómo la herramienta separa distintos aspectos del resultado: por un lado, analiza si la imagen conserva rasgos propios del pixel-art y, por otro, si la figura principal se interpreta con facilidad. Gracias a esta validación, la comparación entre modelos puede apoyarse en métricas técnicas coherentes con el estilo buscado, además de la observación visual de las imágenes generadas.


\section{Integración final del sistema}

Tras completar las tres iteraciones, los distintos componentes desarrollados quedaron integrados dentro de un mismo flujo de trabajo. El sistema final ejecuta entrenamientos LoRA y conecta la obtención de imágenes, la revisión del dataset, la generación de captions, la preparación de datos, la configuración de entrenamientos, la ejecución local o remota, la generación de imágenes y la evaluación de resultados.

\medskip

La interfaz web actúa como punto de entrada principal del pipeline. Desde ella se pueden lanzar procesos de scraping, revisar imágenes aceptadas o rechazadas, crear grupos de captions, generar descripciones automáticas, mover el material revisado a un dataset de entrenamiento y configurar nuevos trabajos LoRA. Estas acciones se conectan con scripts internos y con el sistema de archivos, manteniendo separados los datos, configuraciones y resultados de cada iteración.

\medskip

La integración con Modal permitió ampliar el sistema más allá del entorno local. Los entrenamientos pueden configurarse desde la interfaz y ejecutarse en una GPU remota, conservando los artefactos generados y manteniendo la estructura de salidas del proyecto. Además, la posibilidad de seleccionar un LoRA previo mantiene el enfoque incremental desarrollado desde la segunda iteración.

\medskip

El sistema final también incorpora la fase posterior al entrenamiento. Los LoRA generados pueden seleccionarse desde la interfaz para producir nuevas imágenes, que quedan organizadas según el modelo utilizado. La evaluación visual y las métricas de apoyo añaden una capa adicional de revisión, permitiendo analizar los resultados generados sin romper la trazabilidad entre dataset, entrenamiento y modelo.

\medskip

En conjunto, la integración final convierte el proyecto en un pipeline reproducible y ampliable. Cada fase produce artefactos que pueden revisarse posteriormente y cada iteración conserva su relación con los datos, configuraciones y modelos utilizados.

\section{Síntesis de la implementación}

La implementación se desarrolló de forma incremental para reducir riesgos y validar progresivamente cada parte del sistema. La primera iteración permitió comprobar que el entorno local era capaz de cargar FLUX.1-dev, entrenar un adaptador LoRA y generar artefactos básicos como checkpoints, configuración y samples.

\medskip

La segunda iteración mejoró la organización del proyecto, introdujo un dataset más cuidado, captions más descriptivos y el uso de \texttt{pretrained_lora_path} para continuar el entrenamiento desde un LoRA anterior. Esta fase consolidó el enfoque iterativo del proyecto y permitió enlazar los resultados de una iteración con la siguiente.

\medskip

La tercera iteración amplió el alcance del sistema mediante automatización y escalado. Se incorporaron la interfaz web adaptada, el scraper de Pixilart, el captioning automático con Qwen2.5-VL, la ejecución de entrenamientos en Modal y la generación y evaluación de imágenes desde la aplicación. Además, permitió detectar y corregir errores reales de configuración, reforzando la importancia de conservar configuraciones, logs y samples durante cada ejecución.
\chapter{Resultados}
\label{resultados}

\section{Enfoque de evaluación}

El análisis de resultados se plantea mediante una comparación controlada entre modelos. Para ello, se genera una imagen con el mismo \textit{prompt}, la misma resolución y los mismos parámetros de inferencia, cambiando únicamente el modelo utilizado. La \textit{seed} queda registrada en el nombre de cada archivo generado, aunque en esta comparación las imágenes no proceden de una misma semilla fija.
\medskip

Este enfoque permite observar de forma directa qué aporta cada LoRA respecto al modelo base \texttt{FLUX.1-dev} y cómo evoluciona el sistema a lo largo de las distintas iteraciones. La comparación incluye el modelo base sin adaptación adicional, los LoRA generados en las primeras iteraciones, el entrenamiento problemático ejecutado en Modal y los modelos corregidos posteriores.

\medskip

El \textit{prompt} elegido no se plantea como una instrucción mínima, ya que el objetivo del trabajo no es comprobar únicamente si el modelo genera una figura reconocible. La evaluación busca analizar si los modelos son capaces de responder a una petición razonable de pixel-art, incorporando tanto el sujeto principal como rasgos visuales propios del estilo. Por este motivo, se utiliza un \textit{prompt} de dificultad media, con un sujeto sencillo y una serie de restricciones estilísticas explícitas.

\begin{table}[htbp]
\centering
\estilotabla
\begin{tabular}{|L{0.23\textwidth}|L{0.67\textwidth}|}
\hline
\textbf{Condición} & \textbf{Valor empleado} \\
\hline
\textit{Prompt} & {\small\ttfamily pixelart, small red fox sitting on a mossy stone, fluffy tail curled around its paws, bright curious eyes, forest background, limited color palette, hard edges, clean silhouette, retro game sprite style} \\
Resolución & 512$\times$512 píxeles. \\
\hline
Pasos de inferencia & 20 pasos. \\
\hline
\textit{Guidance scale} & 4. \\
\hline
Número de imágenes & Seis imágenes, una por cada modelo comparado. \\
\hline
Modelos comparados & \texttt{FLUX.1-dev}, LoRA v1, LoRA v2, LoRA v3, LoRA v3 Arreglo y LoRA v4. \\
\hline
\end{tabular}
\caption{Condiciones de generación empleadas en la comparación.}
\label{tab:condiciones_comparacion_resultados}
\end{table}

\section{Resultado principal: comparación entre modelos}

La comparación visual se organiza en seis salidas generadas bajo las condiciones indicadas en la Tabla~\ref{tab:condiciones_comparacion_resultados}. Esta disposición permite revisar de forma conjunta la respuesta del modelo base y de los cinco LoRA entrenados durante el proyecto.

\begin{figure}[H]
\centering

\begin{minipage}[t]{0.31\textwidth}
\centering
\includegraphics[width=\linewidth]{imagenes/resultados_flux_base.png}
\par\smallskip
{\small FLUX.1-dev}
\end{minipage}
\hspace{0.08\textwidth}
\begin{minipage}[t]{0.31\textwidth}
\centering
\includegraphics[width=\linewidth]{imagenes/resultados_lora_v1.png}
\par\smallskip
{\small LoRA v1}
\end{minipage}

\vspace{0.45cm}

\begin{minipage}[t]{0.31\textwidth}
\centering
\includegraphics[width=\linewidth]{imagenes/resultados_lora_v2.png}
\par\smallskip
{\small LoRA v2}
\end{minipage}
\hspace{0.08\textwidth}
\begin{minipage}[t]{0.31\textwidth}
\centering
\includegraphics[width=\linewidth]{imagenes/resultados_lora_v3.png}
\par\smallskip
{\small LoRA v3}
\end{minipage}

\vspace{0.45cm}

\begin{minipage}[t]{0.31\textwidth}
\centering
\includegraphics[width=\linewidth]{imagenes/resultados_lora_v3_arreglo.png}
\par\smallskip
{\small LoRA v3 Arreglo}
\end{minipage}
\hspace{0.08\textwidth}
\begin{minipage}[t]{0.31\textwidth}
\centering
\includegraphics[width=\linewidth]{imagenes/resultados_lora_v4.png}
\par\smallskip
{\small LoRA v4}
\end{minipage}

\caption{Comparación visual entre el modelo base y los LoRA entrenados. Fuente: generaciones propias obtenidas con el modelo base y los LoRA entrenados a partir del prompt indicado.}
\label{fig:resultados_comparacion_modelos}
\end{figure}



El modelo base \texttt{FLUX.1-dev}, al no estar adaptado al dominio concreto del proyecto, tiende a producir una imagen más cercana a una ilustración digital que a un sprite pixel-art. Puede representar parte del contenido indicado en el \textit{prompt}, pero suele presentar suavizado, transiciones graduales y una estructura de píxel poco definida.

\medskip

El primer LoRA, \texttt{pixel\_art\_lora\_v1}, muestra que el entrenamiento inicial fue capaz de introducir rasgos propios del pixel-art, especialmente bordes más duros, formas simplificadas y una estética más cercana a videojuegos retro. Este resultado confirmó la viabilidad del enfoque, aunque seguía estando limitado por el tamaño reducido del dataset inicial.

\medskip

Con \texttt{pixel\_art\_lora\_v2} se observa una mejora incremental. El modelo mantiene la adaptación al estilo y permite comprobar que el encadenamiento entre LoRA funcionaba correctamente. Aun así, algunas generaciones todavía podían mostrar artefactos, zonas poco limpias o elementos poco definidos.

\medskip

El caso de \texttt{pixel\_art\_lora\_v3} es especialmente relevante. Aunque el entrenamiento se completó y generó sus artefactos, la imagen obtenida se aleja del objetivo visual. Esta diferencia permite comprobar que un entrenamiento puede finalizar correctamente a nivel técnico y producir un modelo poco útil para la generación final.

\medskip

El modelo \texttt{pixel\_art\_lora\_v3\_Arreglo} muestra la recuperación del estilo tras corregir la configuración del entrenamiento. En comparación con \texttt{pixel\_art\_lora\_v3}, se aprecia una imagen más limpia, con mayor coherencia visual y una estructura más cercana al pixel-art buscado.

\medskip

Finalmente, \texttt{pixel\_art\_lora\_v4} representa una mejora posterior partiendo del modelo corregido y utilizando captions de detalle medio. Este resultado permite valorar si captions más compactos ayudan a estabilizar el estilo aprendido. Aunque todavía pueden aparecer limitaciones como ruido, texto residual o zonas suavizadas, el modelo refleja mejor el objetivo global del proyecto que el modelo base y que el entrenamiento problemático.
\section{Comparación de modelos generados}

La Tabla~\ref{tab:comparacion_modelos_resultados} resume la evolución de los modelos utilizados en la comparación principal. Su finalidad es relacionar cada modelo con su origen y con la observación visual principal, evitando repetir todos los parámetros internos del entrenamiento, ya explicados en el capítulo de implementación.

\begin{table}[H]
\centering
\estilotabla
\begin{tabular}{|L{0.22\textwidth}|L{0.16\textwidth}|L{0.54\textwidth}|}
\hline
\textbf{Modelo} & \textbf{Origen} & \textbf{Resultado observado} \tabularnewline
\hline
\texttt{FLUX.1-dev} & Modelo base & Imagen más ilustrativa, con poca estructura pixel-art. \tabularnewline
\hline
LoRA v1 & Iteración 1 & Primera adaptación válida al estilo pixel-art. \tabularnewline
\hline
LoRA v2 & Iteración 2 & Mejora incremental y validación del encadenamiento entre LoRA. \tabularnewline
\hline
LoRA v3 & Iteración 3 & Entrenamiento completado, con resultado visual alejado del objetivo. \tabularnewline
\hline
LoRA v3 Arreglo & Iteración 3 & Recuperación del estilo tras corregir la configuración. \tabularnewline
\hline
LoRA v4 & Iteración 3 &Entrenamiento posterior con captions medios; mantiene buena lectura visual, aunque no mejora todas las métricas. \tabularnewline
\hline
\end{tabular}
\caption{Resumen de los modelos comparados.}
\label{tab:comparacion_modelos_resultados}
\end{table}

\FloatBarrier

\section{Resultados del pipeline desarrollado}

Además de los modelos obtenidos, uno de los resultados principales del proyecto es el propio pipeline desarrollado. El sistema conecta la obtención de imágenes, la preparación del dataset, el entrenamiento, la generación y la evaluación dentro de un flujo común. Esto permite repetir el proceso con nuevas imágenes, nuevos captions o nuevas configuraciones de entrenamiento.

\medskip

La organización por iteraciones facilita comparar modelos sin sobrescribir resultados anteriores. Cada entrenamiento conserva sus configuraciones, salidas, \textit{samples}, checkpoints, logs y modelos finales, lo que permite reconstruir el proceso seguido en cada fase. La Tabla~\ref{tab:resultados_pipeline_funcional} resume los resultados funcionales obtenidos en las principales partes del pipeline.

\medskip

En conjunto, el sistema cumple el objetivo de ofrecer una base reproducible para entrenar y comparar modelos LoRA orientados a pixel-art. La conservación de datasets, configuraciones, checkpoints, \textit{samples}, logs y modelos finales permite revisar el proceso seguido en cada iteración y relacionar los resultados visuales con las decisiones tomadas durante la implementación.

\begin{table}[H]
\centering
\estilotabla
\begin{tabular}{|L{0.22\textwidth}|L{0.64\textwidth}|}
\hline
\textbf{Parte del sistema} & \textbf{Resultado obtenido} \tabularnewline
\hline
Scraper & Obtención automática de imágenes de Pixilart mediante topics, filtros automáticos y revisión manual. \tabularnewline
\hline
Captioning & Generación automática de captions con Qwen2.5-VL y revisión posterior desde la interfaz. \tabularnewline
\hline
Datasets & Organización por iteraciones con imágenes, captions CSV y archivos \texttt{.txt} compatibles con entrenamiento. \tabularnewline
\hline
Entrenamiento local & Validación inicial del entrenamiento LoRA y del encadenamiento entre modelos. \tabularnewline
\hline
Entrenamiento en Modal & Ejecución de entrenamientos en GPU cloud, conservación de logs y descarga de artefactos. \tabularnewline
\hline
Generación & Generación de imágenes desde la interfaz seleccionando modelo, \textit{prompt} y \textit{seed}. \tabularnewline
\hline
Evaluación & Análisis técnico de imágenes generadas mediante métricas de estructura pixel-art, lectura visual y paleta. \tabularnewline
\hline
\end{tabular}
\caption{Resultados funcionales del pipeline.}
\label{tab:resultados_pipeline_funcional}
\end{table}


\FloatBarrier

\section{Evaluación automática como apoyo}

Además de la comparación visual, el sistema incorpora una evaluación automática de imágenes generadas. Esta evaluación se ejecuta mediante \texttt{scripts/tfg/evaluate_style.py} y se integra en la interfaz mediante la ruta \texttt{/api/generated-images/evaluate}.

\medskip

La evaluación automática se utiliza como apoyo, no como valoración definitiva. Sus métricas permiten analizar rasgos técnicos del pixel-art, como la estructura del píxel, la dureza de los bordes, la lectura de la figura y la compactación de la paleta. En cambio, no evalúa la belleza de la imagen, la composición artística ni si el contenido responde perfectamente al \textit{prompt}.

\medskip

Para mantener la coherencia del análisis, las métricas se han calculado sobre las mismas seis imágenes incluidas en la Figura~\ref{fig:resultados_comparacion_modelos}. De esta forma, la evaluación automática refuerza la comparación principal y evita mezclar resultados procedentes de generaciones distintas.

La Tabla~\ref{tab:evaluacion_automatica_resultados} recoge las métricas automáticas obtenidas para las seis imágenes comparadas.

\begin{table}[H]
\centering
\scriptsize
\setlength{\tabcolsep}{3pt}
\renewcommand{\arraystretch}{1.08}
\begin{tabular}{|L{0.17\textwidth}|C{0.13\textwidth}|C{0.13\textwidth}|C{0.10\textwidth}|L{0.35\textwidth}|}
\hline
\textbf{Modelo} & \textbf{Constr. pixel-art} & \textbf{Lectura visual} & \textbf{Paleta} & \textbf{Comentario} \tabularnewline
\hline
\texttt{FLUX.1-dev} & 1.0/10 & 8.7/10 & 2.6/10 & Figura reconocible, pero con suavizado y estructura de píxel irregular. \tabularnewline
\hline
LoRA v1 & 7.7/10 & 9.7/10 & 2.4/10 & Buena cuadrícula y lectura visual; paleta poco compacta. \tabularnewline
\hline
LoRA v2 & 8.1/10 & 10/10 & 6.9/10 & Mejor equilibrio global en las métricas automáticas. \tabularnewline
\hline
LoRA v3 & 4.5/10 & 4.7/10 & 6.1/10 & Caída clara en lectura visual y regularidad de píxel. \tabularnewline
\hline
LoRA v3 Arreglo & 7.7/10 & 10/10 & 5.0/10 & Recupera estructura y claridad tras la corrección. \tabularnewline
\hline
LoRA v4 & 1.8/10 & 9.9/10 & 3.0/10 & Figura clara, pero construcción pixel-art baja. \tabularnewline
\hline
\end{tabular}
\caption{Métricas automáticas obtenidas en la comparación controlada.}
\label{tab:evaluacion_automatica_resultados}
\end{table}

Los resultados automáticos muestran que el modelo base consigue una lectura visual alta, pero no construye una estructura pixel-art sólida. Esto confirma que una imagen puede ser reconocible y, al mismo tiempo, alejarse del estilo objetivo por suavizado, ausencia de cuadrícula o falta de bordes definidos.

\medskip

Los modelos LoRA v1 y LoRA v2 mejoran claramente la construcción pixel-art respecto al modelo base. En especial, LoRA v2 obtiene el mejor equilibrio global en la evaluación automática, con buena estructura de píxel, lectura visual máxima y una paleta más compacta que la del resto de modelos iniciales.

\medskip

El resultado de LoRA v3 coincide con el problema detectado durante la implementación. Aunque el entrenamiento se completó correctamente a nivel técnico, la evaluación refleja una caída notable en lectura visual y una construcción pixel-art menos estable. Esto refuerza el diagnóstico realizado sobre la configuración del entrenamiento en Modal.

\medskip

LoRA v3 Arreglo muestra una recuperación importante tras aplicar las correcciones. La construcción pixel-art vuelve a situarse en un valor alto y la lectura visual alcanza la puntuación máxima, lo que respalda la importancia de haber conservado configuraciones, logs y \textit{samples} durante el proceso.

\medskip

Por último, LoRA v4 mantiene una lectura visual muy alta, pero obtiene una puntuación baja en construcción pixel-art. Este resultado indica que el entrenamiento posterior con captions medios no mejoró todas las dimensiones evaluadas. La imagen conserva una figura reconocible, aunque pierde regularidad de píxel según las métricas automáticas.


\section{Limitaciones observadas}

Aunque los resultados muestran una evolución clara, el sistema mantiene varias limitaciones. En primer lugar, el dataset final sigue siendo reducido para entrenar un estilo completamente robusto. El conjunto de 130 imágenes permitió mejorar la escala del entrenamiento, pero todavía depende mucho de la calidad de las imágenes seleccionadas y de la diversidad real del dataset.

\medskip

También se observaron limitaciones visuales en algunas generaciones. Algunos modelos pueden producir texto residual, zonas suavizadas, ruido local o estructuras de píxel poco consistentes. Estos problemas muestran que el estilo pixel-art no depende únicamente de usar la palabra \texttt{pixelart} en el \textit{prompt}; también requiere datos adecuados, captions equilibrados y una configuración de entrenamiento estable.

\medskip

Otra limitación está relacionada con la trazabilidad de la generación. El sistema conserva el modelo usado, la \textit{seed} y la ruta de salida, pero en generaciones previas el \textit{prompt} completo no siempre queda persistido como metadato estructurado. Esto hace recomendable que las comparaciones finales del capítulo se generen de forma controlada y se documenten explícitamente en la memoria.

\medskip

Por último, la evaluación automática debe interpretarse como una ayuda técnica. Sus métricas permiten detectar rasgos como cuadrícula, bordes o ruido visual, pero no sustituyen el análisis visual ni la valoración del cumplimiento semántico del \textit{prompt}.

\section{Conclusión de resultados}

Los resultados obtenidos muestran una evolución progresiva desde el modelo base hasta los LoRA finales. La comparación controlada permite observar que el modelo base no está especializado en pixel-art, que los primeros LoRA introducen rasgos del estilo, que el entrenamiento \texttt{pixel\_art\_lora\_v3} evidencia una configuración problemática y que los modelos corregidos recuperan una dirección más adecuada.

\medskip

El principal resultado del proyecto es haber construido un pipeline capaz de producir, comparar y analizar estas diferencias de forma ordenada. El sistema desarrollado permite preparar datasets, generar captions, entrenar en local o en Modal, conservar artefactos, generar imágenes y evaluarlas. Esta base deja el proyecto preparado para nuevas iteraciones con datasets más amplios, mejores captions y comparaciones más controladas.


\chapter{Conclusiones}

\section{Conclusiones}
\label{conclusiones}

El Trabajo de Fin de Grado ha permitido diseñar, implementar y evaluar un pipeline completo para la adaptación de un modelo generativo de difusión al estilo pixel-art mediante entrenamiento LoRA. El proyecto partía de un problema concreto: los modelos generativos generales pueden producir imágenes visualmente atractivas, pero no siempre respetan de forma consistente las restricciones propias de un estilo visual específico. En el caso del pixel-art, esto se aprecia en la presencia de suavizado, degradados, exceso de detalle, pérdida de bordes definidos o falta de una estructura de píxel regular.

\medskip

A partir de este planteamiento, se ha construido un sistema que integra las fases principales necesarias para trabajar con este tipo de adaptación: obtención de imágenes, filtrado, revisión manual, generación de captions, preparación del dataset, configuración de entrenamientos, ejecución local o remota, generación de imágenes y evaluación de resultados. Esta integración constituye una de las principales aportaciones del trabajo, ya que el proyecto no se ha limitado a entrenar un único modelo LoRA, sino que ha desarrollado un flujo reproducible y ampliable para realizar nuevas iteraciones.

\medskip

La implementación incremental ha sido una decisión acertada para el tipo de proyecto desarrollado. La primera iteración permitió validar el entrenamiento local con FLUX.1-dev y comprobar que el entorno era capaz de generar un adaptador LoRA funcional. La segunda iteración mejoró la organización del repositorio, incorporó captions más descriptivos y validó el uso de \texttt{pretrained_lora_path} para continuar el entrenamiento desde un LoRA anterior. La tercera iteración amplió el sistema mediante interfaz web, scraping, captioning automático, entrenamiento en Modal, generación de imágenes y evaluación técnica. Esta evolución permitió reducir riesgos y justificar cada ampliación a partir de las limitaciones detectadas en la fase anterior.

\medskip

También se ha comprobado la importancia de la trazabilidad en proyectos basados en experimentación con modelos generativos. Conservar datasets, captions, configuraciones, checkpoints, logs, \textit{samples} y modelos finales ha permitido comparar iteraciones, detectar errores y reconstruir el proceso seguido en cada entrenamiento. Esta trazabilidad fue especialmente relevante durante el diagnóstico del entrenamiento problemático en Modal, donde la conservación de \texttt{config.yaml} y de las muestras generadas permitió identificar diferencias de configuración que afectaban directamente a los resultados.

\medskip

Los resultados muestran que la adaptación mediante LoRA puede introducir rasgos propios del pixel-art en el modelo base. En la comparación controlada, el modelo \texttt{FLUX.1-dev} generó una imagen reconocible, pero con una construcción pixel-art muy débil. En cambio, los primeros adaptadores LoRA mejoraron claramente la estructura de píxel y la lectura visual. En particular, LoRA v2 obtuvo el mejor equilibrio en la evaluación automática, combinando buena construcción pixel-art, lectura visual alta y una paleta más compacta.

\medskip

El caso de LoRA v3 permitió extraer una conclusión importante: que un entrenamiento finalice correctamente a nivel técnico no garantiza que el modelo resultante sea útil para el objetivo visual buscado. El entrenamiento generó sus artefactos, checkpoints y modelo final, pero la comparación visual y las métricas mostraron una caída clara en la calidad del resultado. La posterior corrección de la configuración y el entrenamiento de LoRA v3 Arreglo demostraron que el problema no estaba en la ejecución en Modal ni en la transferencia de archivos, sino en una configuración no alineada con el comportamiento esperado durante la inferencia.

\medskip

La evaluación automática incorporada al sistema ha resultado útil como herramienta de apoyo. Al separar construcción pixel-art, lectura visual y paleta compacta, permite observar diferencias que no siempre son evidentes con una revisión visual rápida. Por ejemplo, una imagen puede ser muy legible y representar correctamente el sujeto, pero obtener una puntuación baja en construcción pixel-art si presenta suavizado, cuadrícula irregular o bordes poco definidos. Esta separación ayuda a valorar los resultados de forma más técnica y coherente con el objetivo del trabajo.

\medskip

A pesar de los avances obtenidos, el proyecto mantiene limitaciones. El dataset final, aunque más amplio que en las primeras iteraciones, sigue siendo reducido para entrenar un estilo completamente robusto. Además, la calidad de los resultados depende de la coherencia de las imágenes seleccionadas, del equilibrio de los captions y de la configuración concreta del entrenamiento. También se ha observado que añadir más entrenamiento o utilizar captions distintos no siempre mejora todas las dimensiones evaluadas, como refleja el comportamiento de LoRA v4.

\medskip

En relación con los objetivos planteados, el trabajo ha cumplido su propósito principal. Se han estudiado los fundamentos necesarios para comprender el uso de modelos de difusión y LoRA, se han definido criterios para construir datasets orientados a pixel-art, se ha implementado un flujo de recopilación y preparación de imágenes, se ha incorporado generación automática de captions, se han ejecutado entrenamientos locales y en la nube, se ha desarrollado una interfaz web de apoyo y se han analizado los resultados obtenidos mediante comparación visual y métricas técnicas.

\medskip

Como conclusión general, el proyecto demuestra que la generación de imágenes en un estilo específico no depende únicamente del modelo base ni del \textit{prompt} utilizado. Requiere un proceso completo de preparación de datos, control de captions, configuración de entrenamiento, seguimiento de artefactos y evaluación posterior. El sistema desarrollado proporciona una base funcional para seguir experimentando con la adaptación de modelos generativos al estilo pixel-art y deja preparado un flujo de trabajo reutilizable para futuras mejoras.

\section{Líneas de trabajo futuras}
\label{futuro}

A partir del trabajo realizado, existen varias líneas de evolución que permitirían mejorar la calidad de los modelos obtenidos y ampliar las capacidades del pipeline desarrollado.

\medskip

La primera línea de trabajo futura sería ampliar y depurar el dataset. Aunque el conjunto final permitió realizar entrenamientos más completos que en las primeras iteraciones, sigue siendo limitado para aprender un estilo visual con suficiente robustez. Sería recomendable aumentar el número de imágenes aceptadas, revisar con mayor detalle la coherencia entre ejemplos y separar subconjuntos por tipo de contenido, como personajes, animales, objetos, escenarios o iconos. Esta división permitiría entrenar adaptadores más especializados o analizar qué tipos de imagen aprende mejor el modelo.

\medskip

También sería interesante mejorar el sistema de filtrado automático. Actualmente se aplican criterios como el descarte de imágenes corruptas, animaciones, canvas demasiado grandes, tags problemáticas o imágenes con baja valoración. En futuras versiones se podrían añadir filtros visuales más avanzados para detectar antialiasing, exceso de degradados, baja legibilidad, resoluciones poco adecuadas o imágenes que mezclen pixel-art con ilustración digital suavizada. De esta forma, el dataset final dependería menos de una revisión manual extensa.

\medskip

Otra línea importante sería profundizar en la generación y revisión de captions. El uso de Qwen2.5-VL permitió automatizar una parte importante del proceso, pero los resultados muestran que el nivel de detalle de los captions influye en el comportamiento del entrenamiento. Sería conveniente comparar de forma más sistemática captions cortos, medios y detallados, así como estudiar qué términos ayudan realmente al aprendizaje del estilo. También podría incorporarse una fase de validación automática de captions para detectar descripciones demasiado largas, repetitivas o poco relacionadas con la imagen.

\medskip

En cuanto al entrenamiento, sería útil realizar experimentos más controlados modificando un único parámetro cada vez. En este trabajo se corrigieron parámetros relevantes como \texttt{bypass_guidance_embedding}, \texttt{quantize_te} o \texttt{timestep_type}, pero futuras iteraciones podrían estudiar con más detalle el efecto del número de pasos, el \textit{learning rate}, el rango LoRA, el tamaño del dataset, el tipo de captions o el LoRA previo utilizado. Esto permitiría obtener conclusiones más precisas sobre qué configuración funciona mejor para pixel-art.

\medskip

Otra posible mejora sería conservar de forma estructurada todos los metadatos de generación. Actualmente el sistema organiza las imágenes generadas por modelo y conserva información como la ruta y la semilla, pero sería recomendable almacenar también el \textit{prompt} completo, los parámetros de inferencia, el modelo seleccionado, la fecha de generación y la versión exacta del LoRA. Esta información facilitaría repetir comparaciones, generar informes automáticos y reconstruir con precisión el origen de cada imagen.

\medskip

La evaluación automática también podría ampliarse. Las métricas actuales permiten analizar construcción pixel-art, lectura visual y paleta compacta, pero todavía no valoran aspectos como la coherencia semántica con el \textit{prompt}, la composición, la presencia de errores anatómicos o la similitud con ejemplos del dataset. En futuras versiones se podrían combinar métricas técnicas con modelos de visión que comparen la imagen generada con el texto introducido, siempre manteniendo la revisión humana como criterio final.

\medskip

En relación con la interfaz web, una línea de evolución sería mejorar la experiencia de comparación entre modelos. El sistema ya permite generar imágenes con distintos LoRA y revisar sus métricas, pero podría añadirse una vista específica de comparación donde el usuario seleccione varios modelos, un mismo \textit{prompt} y una misma \textit{seed}, generando automáticamente una cuadrícula con resultados y puntuaciones. Esto facilitaría la evaluación de nuevas iteraciones sin tener que organizar manualmente las imágenes.

\medskip

También sería posible incorporar funcionalidades de postprocesado orientadas a pixel-art. Por ejemplo, reducción controlada de paleta, ajuste de resolución, eliminación de suavizado, cuantización de colores o conversión a una cuadrícula de píxel más regular. Estas técnicas no sustituirían al entrenamiento, pero podrían ayudar a mejorar la presentación final de algunas imágenes generadas y permitir comparar el resultado antes y después del postprocesado.

\medskip

Otra línea de trabajo sería probar el pipeline con otros modelos base o con otras técnicas de adaptación. Aunque este proyecto se ha centrado en FLUX.1-dev y LoRA, la estructura desarrollada podría adaptarse a otros modelos texto-imagen o a variantes de entrenamiento más específicas. Esto permitiría comprobar si el comportamiento observado depende del modelo base utilizado o del propio enfoque de preparación de datos y captions.

\medskip

Finalmente, el pipeline podría generalizarse a otros estilos visuales. Aunque el trabajo se ha orientado al pixel-art, muchas de las fases desarrolladas son reutilizables: recopilación de imágenes, revisión, captioning, entrenamiento LoRA, generación y evaluación. Adaptando los criterios de filtrado, los captions y las métricas de evaluación, el sistema podría aplicarse a otros estilos gráficos, como ilustración vectorial, arte isométrico, sprites de videojuegos, iconografía o estilos artísticos concretos.

\medskip

En conjunto, estas líneas futuras permitirían convertir el sistema desarrollado en una herramienta más completa para la experimentación con modelos generativos especializados. El trabajo actual deja una base funcional y trazable sobre la que seguir ampliando datasets, mejorando entrenamientos y evaluando de forma más precisa la fidelidad a estilos visuales específicos.



%%\nocite{*} %incluye TODOS los documentos de la base de datos bibliográfica sean o no citados en el texto
\bibliography{bibliografia/bibliografia}
\addcontentsline{toc}{chapter}{Bibliografía} %sustituir bibliografía con el nombre del fichero bibtex con la bibliografía
\bibliographystyle{apalike}
%


\appendix
\chapter{Estructura del repositorio}
\label{anexo:estructura_repositorio}

El repositorio del proyecto se organiza sobre una base existente de \texttt{ai-toolkit}, ampliada con una interfaz web, scripts propios para la preparación del dataset, generación automática de captions, ejecución de entrenamientos y generación de imágenes. El Listado~\ref{lst:estructura_repositorio} muestra una vista resumida de las carpetas y archivos más relevantes para comprender la estructura del sistema.

\begin{lstlisting}[caption={Estructura principal del repositorio del proyecto. Fuente: elaboración propia.}, label={lst:estructura_repositorio}]
ai-toolkit/
|-- config/
|   |-- iter_01_train.yaml
|   |-- iter_02_train.yaml
|   |-- iter_03_Arreglo_train.yaml
|   `-- train_lora_pixelart_tfg.yaml
|
|-- datasets/
|   |-- iter_01/
|   |   `-- images/
|   |       |-- *.png
|   |       `-- *.txt
|   |
|   |-- iter_02/
|   |   |-- raw/
|   |   |-- images/
|   |   |-- captions.csv
|   |   `-- notes.md
|   |
|   |-- iter_03/
|   |   `-- images/
|   |       |-- *.png
|   |       |-- *.txt
|   |       `-- captions.csv
|   |
|   |-- iter_03_captions_medios/
|   |   `-- images/
|   |       |-- *.png
|   |       |-- *.txt
|   |       `-- captions.csv
|   |
|   `-- pixilart/
|       |-- raw/
|       |-- filtered/
|       |   |-- accepted/
|       |   |-- rejected/
|       |   |-- manual_accepted/
|       |   `-- metadata/
|       |-- reports/
|       `-- scrapes/
|
|-- caption_groups/
|   `-- <grupo>/
|       |-- images/
|       |-- captions.csv
|       |-- captions.auto.csv
|       |-- manifest.json
|       `-- progress.json
|
|-- scripts/
|   |-- tfg/
|   |   |-- generate_captions.py
|   |   |-- generate_api.py
|   |   |-- evaluate_style.py
|   |   |-- scale_pixelart.py
|   |   |-- sync_captions.py
|   |   
|   |
|   `-- scrapers/
|       `-- pixilart_scraper/
|           |-- run_spider.py
|           |-- pixilart_scraper/
|           |   |-- spiders/
|           |   |   `-- pixilart_spider.py
|           |   |-- pipelines.py
|           |   |-- items.py
|           |   |-- settings.py
|           |   `-- utils.py
|           `-- README.md
|
|-- ui/
|   |-- package.json
|   |-- prisma/
|   |   `-- schema.prisma
|   |
|   `-- src/
|       |-- app/
|       |   |-- scraper/
|       |   |-- captions/
|       |   |-- datasets/
|       |   |-- jobs/
|       |   |-- trainings/
|       |   |-- generate/
|       |   `-- api/
|       |
|       |-- server/
|       |   |-- scraper.ts
|       |   |-- captions.ts
|       |   `-- imageEvaluations.ts
|       |
|       |-- components/
|       `-- utils/
|
|-- outputs/
|   |-- flux_base/
|   |   `-- generated/
|   |
|   |-- iter_01/
|   |   |-- generated/
|   |   `-- pixel_art_lora_v1/
|   |       |-- samples/
|   |       |-- checkpoints/
|   |       |-- config.yaml
|   |       `-- pixel_art_lora_v1.safetensors
|   |
|   |-- iter_02/
|   |   |-- generated/
|   |   `-- pixel_art_lora_v2/
|   |       |-- samples/
|   |       |-- checkpoints/
|   |       |-- config.yaml
|   |       `-- pixel_art_lora_v2.safetensors
|   |
|   |-- iter_03/
|   |   |-- generated/
|   |   `-- pixel_art_lora_v3/
|   |
|   |-- iter_03_Arreglo/
|   |   |-- generated/
|   |   `-- pixel_art_lora_v3_Arreglo/
|   |
|   |-- iter_04/
|   |   |-- generated/
|   |   `-- pixel_art_lora_v4/
|   |
|   `-- .image-evaluations.json
|
|-- output/
|   |-- pixel_art_lora_v3_Arreglo/
|   `-- pixel_art_lora_v4/
|
|-- toolkit/
|-- extensions/
|-- extensions_built_in/
|-- jobs/
|-- assets/
|-- venv/
|
|-- aitk_db.db
|-- run.py
|-- run_modal.py
|-- requirements.txt
|-- README_TFG_PIPELINE.md
`-- README_ORIGINAL_AI_TOOLKIT.md
\end{lstlisting}

En el esquema se han omitido carpetas internas de control o caché, como \texttt{.git/}, \texttt{\_\_pycache\_\_/}, \texttt{tmp/} y otros archivos auxiliares, ya que no aportan información relevante sobre el diseño funcional del sistema. La estructura muestra únicamente los elementos necesarios para entender cómo se relacionan la interfaz, los scripts, los datasets, las configuraciones de entrenamiento y los artefactos generados.


\chapter*{Glosario de términos y abreviaturas}
\addcontentsline{toc}{chapter}{Glosario}

\begin{longtable}{|L{0.25\textwidth}|L{0.65\textwidth}|}
\hline
\textbf{Término} & \textbf{Definición} \\
\hline

IA generativa & Rama de la inteligencia artificial orientada a crear contenido nuevo, como texto, imágenes, audio o código, a partir de patrones aprendidos durante el entrenamiento. \\
\hline

Modelo de difusión & Modelo generativo que produce imágenes mediante un proceso progresivo de eliminación de ruido hasta obtener una imagen coherente. \\
\hline

FLUX.1-dev & Modelo generativo texto-imagen utilizado como base en el proyecto para entrenar adaptadores LoRA orientados al estilo pixel-art. \\
\hline

LoRA & Técnica de adaptación eficiente que permite ajustar un modelo preentrenado entrenando un conjunto reducido de parámetros adicionales, sin modificar por completo el modelo base. \\
\hline

Fine-tuning & Proceso de continuar el entrenamiento de un modelo preentrenado con un nuevo conjunto de datos para adaptarlo a una tarea, dominio o estilo concreto. \\
\hline

Pixel-art & Estilo de arte digital basado en el uso visible del píxel, formas simplificadas, bordes definidos y paletas de color limitadas. \\
\hline

Dataset & Conjunto de imágenes y descripciones utilizadas para entrenar o adaptar un modelo generativo. \\
\hline

Caption & Descripción textual asociada a una imagen del dataset. En este proyecto se utiliza para indicar al modelo qué representa cada imagen y qué rasgos visuales debe aprender. \\
\hline

Trigger word & Palabra clave utilizada de forma constante en captions y prompts para asociar un término concreto al estilo aprendido. En este trabajo se utiliza \texttt{pixelart}. \\
\hline

Prompt & Entrada textual introducida por el usuario para indicar al modelo qué imagen debe generar. \\
\hline

Inferencia & Proceso de generación de una imagen a partir de un modelo ya entrenado y un prompt de entrada. \\
\hline

Seed & Valor numérico utilizado para controlar la aleatoriedad de una generación. Permite repetir o comparar resultados bajo condiciones similares. \\
\hline

Guidance scale & Parámetro de inferencia que controla cuánto debe ajustarse la generación al prompt introducido. \\
\hline

Checkpoint & Copia intermedia de un modelo o adaptador guardada durante el entrenamiento, útil para recuperar una ejecución o analizar su evolución. \\
\hline

Sample & Imagen generada durante el entrenamiento para observar visualmente la evolución del modelo en distintos pasos. \\
\hline

Pipeline & Secuencia organizada de fases que transforma una entrada inicial en un resultado final. En este proyecto incluye scraping, revisión, captioning, entrenamiento, generación y evaluación. \\
\hline

Scraping & Proceso automatizado de recopilación de imágenes desde una fuente web, utilizado en este trabajo para ampliar el dataset a partir de Pixilart. \\
\hline

Qwen2.5-VL & Modelo visión-lenguaje utilizado para generar captions automáticos a partir de imágenes. \\
\hline

Modal & Plataforma cloud utilizada para ejecutar entrenamientos en GPU remota cuando los recursos locales no son suficientes. \\
\hline

GPU & Unidad de procesamiento gráfico utilizada para acelerar operaciones intensivas de entrenamiento e inferencia en modelos de inteligencia artificial. \\
\hline

VRAM & Memoria de vídeo disponible en la GPU. Su capacidad condiciona el tamaño de los modelos, datasets y configuraciones que pueden entrenarse localmente. \\
\hline

CUDA & Plataforma de computación paralela de NVIDIA que permite utilizar la GPU para acelerar procesos de entrenamiento e inferencia. \\
\hline

safetensors & Formato de archivo utilizado para almacenar pesos de modelos o adaptadores de forma eficiente y segura. \\
\hline

Antialiasing & Técnica de suavizado de bordes que reduce el efecto escalonado de una imagen. En pixel-art suele evitarse porque puede difuminar la estructura del píxel. \\
\hline

Nearest-neighbor & Método de escalado de imágenes que conserva bloques de píxel definidos, utilizado para redimensionar imágenes pixel-art sin introducir suavizado. \\
\hline

Paleta limitada & Conjunto reducido de colores utilizado en una imagen. Es una característica habitual del pixel-art y ayuda a mantener una estética limpia y coherente. \\
\hline

\end{longtable}
\end{document}
 
